% =====================================================================
%  Initial research proposal (working draft, internal)
%
%  Files:  main.tex                    this document
%          preamble_maps.tex           colours / macros / TikZ styles
%          background_boundaries.tex   Appendix A.1 figure
%          background_maps.tex         Appendix A.2 equation maps
%          references.bib
%  Build:  latexmk -pdf main.tex   (or pdflatex, bibtex, pdflatex x2)
%
%  \todo{...} and \todoitems{...} print in red. Set \todofalse to hide them.
% =====================================================================
\documentclass[11pt]{article}
\usepackage[a4paper,margin=2.2cm]{geometry}
\usepackage{amsmath,amssymb,array,booktabs,tabularx,enumitem}
\usepackage[numbers,sort&compress]{natbib}

% equation-map part headings sit under Appendix A.2
\newcommand{\sect}[1]{\subsubsection*{#1}}
% =====================================================================
%  preamble_maps.tex -- packages, colours, macros and TikZ styles for the
%  equation maps. \input from the preamble of any document that
%  % =====================================================================
%  background_maps.tex -- equation maps (body only).
%  Online: Liniger MPCC (C++ master). Offline: TUM global_racetrajectory_optimization.
%  \input inside a document whose preamble % =====================================================================
%  preamble_maps.tex -- packages, colours, macros and TikZ styles for the
%  equation maps. \input from the preamble of any document that
%  \input{background_maps}. Requires amsmath, amssymb, array.
% =====================================================================
\usepackage{xcolor}
\usepackage{tikz}
\usetikzlibrary{arrows.meta,calc,decorations.pathmorphing,decorations.pathreplacing}

% ---------------- role colours ----------------------------------------
\definecolor{cS}{HTML}{0072B2}  % decision state
\definecolor{cU}{HTML}{C76E00}  % control input
\definecolor{cR}{HTML}{00876C}  % reference / track data
\definecolor{cP}{HTML}{7A7A7A}  % fixed parameter
\definecolor{cD}{HTML}{8E44AD}  % derived quantity
\definecolor{cL}{HTML}{C0392B}  % limit / bound
\newcommand{\xS}[1]{{\color{cS}#1}}
\newcommand{\xU}[1]{{\color{cU}#1}}
\newcommand{\xR}[1]{{\color{cR}#1}}
\newcommand{\xP}[1]{{\color{cP}#1}}
\newcommand{\xD}[1]{{\color{cD}#1}}
\newcommand{\xL}[1]{{\color{cL}#1}}

% ---------------- layout ----------------------------------------------
\newenvironment{panel}[2]{%
  \par\addvspace{5pt}\noindent\begin{minipage}{\linewidth}%
  {\sffamily\small\bfseries\color{black!70}#1}\enspace{\sffamily\small #2}\par\vspace{4pt}}%
  {\end{minipage}\par\vspace{6pt}\noindent{\color{black!12}\rule{\linewidth}{0.5pt}}\par}
\newenvironment{figside}[1][0.53]{\begin{minipage}[c]{#1\linewidth}\centering}{\end{minipage}\hfill}
\newenvironment{eqside}[1][0.45]{\begin{minipage}[c]{#1\linewidth}\small
  \setlength{\abovedisplayskip}{3pt}\setlength{\belowdisplayskip}{3pt}%
  \setlength{\abovedisplayshortskip}{2pt}\setlength{\belowdisplayshortskip}{2pt}}{\end{minipage}}
\newcommand{\deps}[1]{\par\smallskip{\sffamily\scriptsize\bfseries\color{black!45}DEPENDS}\enspace{\small#1}}
\newcommand{\restr}[1]{\par\smallskip{\sffamily\scriptsize\bfseries\color{cL}RESTRAINS}\enspace{\small#1}}
\newcommand{\chip}[2]{\tikz[baseline=-0.6ex]\fill[#1] (0,0) circle (3pt);\,{\small\sffamily\color{#1}#2}\quad}
\providecommand{\sect}[1]{\par\vspace{8pt}{\large\sffamily\bfseries #1}\par\vspace{2pt}\noindent{\color{black!40}\rule{\linewidth}{0.8pt}}\par}

% ---------------- tikz styles -----------------------------------------
\tikzset{
  every picture/.style={line cap=round,line join=round,>={Stealth[length=4.5pt,width=3.2pt]}},
  ax/.style={->,black!45,line width=0.5pt},
  body/.style={draw=black!65,fill=black!5,line width=0.6pt,rounded corners=1.2pt},
  whl/.style={fill=black!75,draw=none,rounded corners=0.6pt,minimum width=5.5mm,minimum height=2mm,inner sep=0pt},
  vS/.style={->,cS,line width=0.9pt},
  vD/.style={->,cD,line width=0.9pt},
  vec/.style={->,line width=0.8pt},
  refl/.style={cR,line width=1pt},
  bnd/.style={black!55,line width=0.7pt},
  liml/.style={cL,line width=0.8pt},
  guide/.style={black!35,dash pattern=on 2pt off 1.5pt,line width=0.4pt},
  dim/.style={<->,black!50,line width=0.4pt,>={Stealth[length=3pt,width=2.4pt]}},
  lab/.style={font=\footnotesize,inner sep=1pt},
  sm/.style={font=\scriptsize,inner sep=1pt},
  dot/.style={circle,fill,inner sep=0pt,minimum size=3.4pt},
  odot/.style={circle,draw,fill=white,inner sep=0pt,minimum size=3.8pt,line width=0.6pt},
}
% small car icon: \caricon{(x,y)}{heading deg}{scale}
\newcommand{\caricon}[3]{%
  \begin{scope}[shift={#1},rotate=#2,scale=#3]
    \draw[body] (-0.45,-0.2) rectangle (0.45,0.2);
    \fill[black!55] (0.45,0) -- (0.28,0.11) -- (0.28,-0.11) -- cycle;
  \end{scope}}
% dependency-matrix row: {row}{label}{direct cols}{via cols}  (0 = none)
\newcommand{\mrow}[4]{%
  \draw[black!7,line width=3.6pt] (0.6,-#1) -- (\ncols+0.4,-#1);
  \node[anchor=east,font=\footnotesize,inner sep=1pt] at (0.45,-#1) {$#2$};
  \foreach \c in {#3}{\ifnum\c>0 \fill[black!85] (\c,-#1) circle (2.5pt);\fi}
  \foreach \c in {#4}{\ifnum\c>0 \draw[black!60,line width=0.6pt,fill=white] (\c,-#1) circle (2.2pt);\fi}}
\newcommand{\mcol}[3]{\node[font=\footnotesize,text=#3] at (#1,0.85) {$#2$};}

% =====================================================================
.
%  \sect{} renders the two part headings; redefine it before % =====================================================================
%  preamble_maps.tex -- packages, colours, macros and TikZ styles for the
%  equation maps. \input from the preamble of any document that
%  \input{background_maps}. Requires amsmath, amssymb, array.
% =====================================================================
\usepackage{xcolor}
\usepackage{tikz}
\usetikzlibrary{arrows.meta,calc,decorations.pathmorphing,decorations.pathreplacing}

% ---------------- role colours ----------------------------------------
\definecolor{cS}{HTML}{0072B2}  % decision state
\definecolor{cU}{HTML}{C76E00}  % control input
\definecolor{cR}{HTML}{00876C}  % reference / track data
\definecolor{cP}{HTML}{7A7A7A}  % fixed parameter
\definecolor{cD}{HTML}{8E44AD}  % derived quantity
\definecolor{cL}{HTML}{C0392B}  % limit / bound
\newcommand{\xS}[1]{{\color{cS}#1}}
\newcommand{\xU}[1]{{\color{cU}#1}}
\newcommand{\xR}[1]{{\color{cR}#1}}
\newcommand{\xP}[1]{{\color{cP}#1}}
\newcommand{\xD}[1]{{\color{cD}#1}}
\newcommand{\xL}[1]{{\color{cL}#1}}

% ---------------- layout ----------------------------------------------
\newenvironment{panel}[2]{%
  \par\addvspace{5pt}\noindent\begin{minipage}{\linewidth}%
  {\sffamily\small\bfseries\color{black!70}#1}\enspace{\sffamily\small #2}\par\vspace{4pt}}%
  {\end{minipage}\par\vspace{6pt}\noindent{\color{black!12}\rule{\linewidth}{0.5pt}}\par}
\newenvironment{figside}[1][0.53]{\begin{minipage}[c]{#1\linewidth}\centering}{\end{minipage}\hfill}
\newenvironment{eqside}[1][0.45]{\begin{minipage}[c]{#1\linewidth}\small
  \setlength{\abovedisplayskip}{3pt}\setlength{\belowdisplayskip}{3pt}%
  \setlength{\abovedisplayshortskip}{2pt}\setlength{\belowdisplayshortskip}{2pt}}{\end{minipage}}
\newcommand{\deps}[1]{\par\smallskip{\sffamily\scriptsize\bfseries\color{black!45}DEPENDS}\enspace{\small#1}}
\newcommand{\restr}[1]{\par\smallskip{\sffamily\scriptsize\bfseries\color{cL}RESTRAINS}\enspace{\small#1}}
\newcommand{\chip}[2]{\tikz[baseline=-0.6ex]\fill[#1] (0,0) circle (3pt);\,{\small\sffamily\color{#1}#2}\quad}
\providecommand{\sect}[1]{\par\vspace{8pt}{\large\sffamily\bfseries #1}\par\vspace{2pt}\noindent{\color{black!40}\rule{\linewidth}{0.8pt}}\par}

% ---------------- tikz styles -----------------------------------------
\tikzset{
  every picture/.style={line cap=round,line join=round,>={Stealth[length=4.5pt,width=3.2pt]}},
  ax/.style={->,black!45,line width=0.5pt},
  body/.style={draw=black!65,fill=black!5,line width=0.6pt,rounded corners=1.2pt},
  whl/.style={fill=black!75,draw=none,rounded corners=0.6pt,minimum width=5.5mm,minimum height=2mm,inner sep=0pt},
  vS/.style={->,cS,line width=0.9pt},
  vD/.style={->,cD,line width=0.9pt},
  vec/.style={->,line width=0.8pt},
  refl/.style={cR,line width=1pt},
  bnd/.style={black!55,line width=0.7pt},
  liml/.style={cL,line width=0.8pt},
  guide/.style={black!35,dash pattern=on 2pt off 1.5pt,line width=0.4pt},
  dim/.style={<->,black!50,line width=0.4pt,>={Stealth[length=3pt,width=2.4pt]}},
  lab/.style={font=\footnotesize,inner sep=1pt},
  sm/.style={font=\scriptsize,inner sep=1pt},
  dot/.style={circle,fill,inner sep=0pt,minimum size=3.4pt},
  odot/.style={circle,draw,fill=white,inner sep=0pt,minimum size=3.8pt,line width=0.6pt},
}
% small car icon: \caricon{(x,y)}{heading deg}{scale}
\newcommand{\caricon}[3]{%
  \begin{scope}[shift={#1},rotate=#2,scale=#3]
    \draw[body] (-0.45,-0.2) rectangle (0.45,0.2);
    \fill[black!55] (0.45,0) -- (0.28,0.11) -- (0.28,-0.11) -- cycle;
  \end{scope}}
% dependency-matrix row: {row}{label}{direct cols}{via cols}  (0 = none)
\newcommand{\mrow}[4]{%
  \draw[black!7,line width=3.6pt] (0.6,-#1) -- (\ncols+0.4,-#1);
  \node[anchor=east,font=\footnotesize,inner sep=1pt] at (0.45,-#1) {$#2$};
  \foreach \c in {#3}{\ifnum\c>0 \fill[black!85] (\c,-#1) circle (2.5pt);\fi}
  \foreach \c in {#4}{\ifnum\c>0 \draw[black!60,line width=0.6pt,fill=white] (\c,-#1) circle (2.2pt);\fi}}
\newcommand{\mcol}[3]{\node[font=\footnotesize,text=#3] at (#1,0.85) {$#2$};}

% =====================================================================

%  to change how they look (e.g. as \subsection).
% =====================================================================
{\setlength{\parindent}{0pt}%
\chip{cS}{decision state}\chip{cU}{control input}\chip{cR}{reference / track data}%
\chip{cP}{fixed parameter}\chip{cD}{derived}\chip{cL}{limit / bound}

% =====================================================================
\sect{Online \quad MPCC \quad (receding horizon, 60 steps $\times$ 20 ms)}

% ---------------------------------------------------------------- A0
\begin{panel}{A0}{Variables and who depends on whom}
\begin{figside}[0.60]
\begin{tikzpicture}[x=0.5cm,y=0.37cm]
  \def\ncols{14}
  \mcol{1}{X}{cS}\mcol{2}{Y}{cS}\mcol{3}{\varphi}{cS}\mcol{4}{v_x}{cS}\mcol{5}{v_y}{cS}
  \mcol{6}{r}{cS}\mcol{7}{s}{cS}\mcol{8}{D}{cS}\mcol{9}{\delta}{cS}\mcol{10}{v_s}{cS}
  \mcol{11}{\dot D}{cU}\mcol{12}{\dot\delta}{cU}\mcol{13}{\dot v_s}{cU}\mcol{14}{\mathbf p_{\!\mathrm{ref}}}{cR}
  \mrow{1}{\xD{\alpha_f}}{4,5,6,9}{0}
  \mrow{2}{\xD{\alpha_r}}{4,5,6}{0}
  \mrow{3}{\xD{F_{f,y}}}{0}{4,5,6,9}
  \mrow{4}{\xD{F_{r,y}}}{0}{4,5,6}
  \mrow{5}{\xD{F_{r,x}}}{4,8}{0}
  \mrow{6}{\xD{F_{\mathrm{fric}}}}{4}{0}
  \mrow{7.4}{\xS{\dot X,\dot Y}}{3,4,5}{0}
  \mrow{8.4}{\xS{\dot\varphi}}{6}{0}
  \mrow{9.4}{\xS{\dot v_x}}{5,6,9}{4,8}
  \mrow{10.4}{\xS{\dot v_y}}{4,6,9}{5}
  \mrow{11.4}{\xS{\dot r}}{9}{4,5,6}
  \mrow{12.4}{\xS{\dot s}}{10}{0}
  \mrow{13.4}{\xS{\dot D,\dot\delta,\dot v_s}}{11,12,13}{0}
  \mrow{14.8}{\xD{e_c,\,e_l}}{1,2,7,14}{0}
  \mrow{15.8}{\xL{\text{track}}}{1,2,14}{0}
  \mrow{16.8}{\xL{\text{ellipse}}}{0}{4,5,6,8}
  \mrow{17.8}{\xL{|\alpha_f|}}{0}{4,5,6,9}
  \mrow{18.8}{\beta\ \text{cost}}{4,5}{0}
  \mrow{19.8}{\varphi\ \text{cost}}{3,7,14}{0}
  \mrow{20.8}{\text{progress}}{10}{0}
  \draw[black!30] (0.5,-6.7) -- (14.5,-6.7);
  \draw[black!30] (0.5,-14.1) -- (14.5,-14.1);
  \draw[black!30] (10.5,1.4) -- (10.5,-21.3);
\end{tikzpicture}
\end{figside}
\begin{eqside}[0.38]
\[\xS{\mathbf x} = [\,\xS{X\ Y\ \varphi\ v_x\ v_y\ r\ s\ D\ \delta\ v_s}\,]^{\top}\]
\[\xU{\mathbf u} = [\,\xU{\dot D\ \ \dot\delta\ \ \dot v_s}\,]^{\top}\]
\[\xS{\dot s}=\xS{v_s},\quad \xS{\dot D}=\xU{u_1},\quad \xS{\dot\delta}=\xU{u_2},\quad \xS{\dot v_s}=\xU{u_3}\]
\medskip
{\footnotesize
\tikz\fill[black!85] (0,0) circle (2.5pt);\ appears in the formula\par
\tikz\draw[black!60,line width=0.6pt] (0,0) circle (2.2pt);\ enters through $\alpha$ or $F$\par}
\restr{inputs are \emph{rates}, so rate costs and slew bounds act directly on $\xU{\mathbf u}$.}
\end{eqside}

\par\medskip
{\footnotesize\renewcommand{\arraystretch}{1.12}%
\begin{minipage}[t]{0.335\linewidth}
{\sffamily\scriptsize\bfseries\color{cS}STATE $\mathbf x$}\par\smallskip
\begin{tabular}{@{}l@{\ \ }>{\raggedright\arraybackslash}p{4.35cm}@{}}
$\xS{X,Y}$       & car position, world frame [m]\\
$\xS{\varphi}$   & heading (yaw angle), world frame [rad]\\
$\xS{v_x}$       & forward velocity, body frame [m/s]\\
$\xS{v_y}$       & sideways velocity, body frame [m/s]\\
$\xS{r}$         & yaw rate [rad/s]\\
$\xS{s}$         & progress: virtual arc length along the reference [m]\\
$\xS{D}$         & motor duty cycle, $[-0.1,1]$; $<0$ brakes\\
$\xS{\delta}$    & front steering angle, $|\delta|\le 0.35$ rad\\
$\xS{v_s}$       & progress speed $\dot s\ge 0$ [m/s]
\end{tabular}
\end{minipage}\hfill
\begin{minipage}[t]{0.32\linewidth}
{\sffamily\scriptsize\bfseries\color{cU}INPUT $\mathbf u$}\par\smallskip
\begin{tabular}{@{}l@{\ \ }>{\raggedright\arraybackslash}p{4.05cm}@{}}
$\xU{\dot D}$         & duty-cycle rate\\
$\xU{\dot\delta}$     & steering rate\\
$\xU{\dot v_s}$       & progress acceleration
\end{tabular}\par\medskip
{\sffamily\scriptsize\bfseries\color{cR}REFERENCE}\par\smallskip
\begin{tabular}{@{}l@{\ \ }>{\raggedright\arraybackslash}p{4.05cm}@{}}
$\xR{\mathbf p_{\mathrm{ref}}}$ & centreline spline $(X_{\mathrm{ref}}(s),Y_{\mathrm{ref}}(s))$ with heading $\Phi(s)$; fixed data
\end{tabular}\par\medskip
{\sffamily\scriptsize\bfseries\color{cD}DERIVED}\par\smallskip
\begin{tabular}{@{}l@{\ \ }>{\raggedright\arraybackslash}p{4.05cm}@{}}
$\xD{\alpha_f,\alpha_r}$       & front / rear tire slip angle\\
$\xD{F_{f,y},F_{r,y}}$         & front / rear lateral tire force\\
$\xD{F_{r,x}}$                 & rear drive force from the motor\\
$\xD{F_{\mathrm{fric}}}$       & rolling resistance + air drag
\end{tabular}
\end{minipage}\hfill
\begin{minipage}[t]{0.31\linewidth}
{\sffamily\scriptsize\bfseries\color{black!60}COST / CONSTRAINT}\par\smallskip
\begin{tabular}{@{}l@{\ \ }>{\raggedright\arraybackslash}p{3.55cm}@{}}
$\dot{(\,\cdot\,)}$        & time derivative: the state equations\\
$\xD{e_c}$                 & contouring error: sideways offset from $\mathbf p_{\mathrm{ref}}(s)$\\
$\xD{e_l}$                 & lag error: along-track gap between $s$ and the car\\
\xL{track}                 & car inside the slab around the reference\\
\xL{ellipse}               & rear friction ellipse\\
$\xL{|\alpha_f|}$          & front slip-angle limit\\
$\beta$ cost               & sideslip $\arctan(v_y/v_x)$ penalty\\
$\varphi$ cost             & heading vs.\ track heading $\Phi(s)$\\
progress                   & reward $-q_{v_s}v_s$
\end{tabular}
\end{minipage}}
\end{panel}

% ---------------------------------------------------------------- A1
\begin{panel}{A1}{Body velocity $\to$ world pose}
\begin{figside}
\begin{tikzpicture}
  \draw[ax] (0,0) -- (1.0,0) node[sm,right]{$X$};
  \draw[ax] (0,0) -- (0,1.0) node[sm,above]{$Y$};
  \coordinate (C) at (2.9,0.45);
  % world components of v (dotted)
  \draw[black!55,dotted,line width=0.9pt] (C) -- ++(1.282,0) node[sm,below,midway,yshift=-2pt]{$\dot X$}
        -- ++(0,1.779) node[sm,right,pos=0.8]{$\dot Y$};
  % heading angle, measured behind the car
  \draw[guide] (C) -- ++(180:1.5);
  \draw[guide] (C) -- ++(210:1.5);
  \draw[cS,line width=0.6pt] ($(C)+(180:1.15)$) arc (180:210:1.15);
  \node[lab,text=cS] at ($(C)+(195:1.42)$) {$\varphi$};
  \begin{scope}[shift={(C)},rotate=30]
    \draw[body] (-0.55,-0.22) rectangle (0.55,0.22);
    \fill[black!55] (0.55,0) -- (0.4,0.1) -- (0.4,-0.1) -- cycle;
    \draw[vS] (0,0) -- (2.0,0) node[lab,below right=0pt]{$v_x$};
    \draw[vS] (0,0) -- (0,0.9) node[lab,above left=0pt]{$v_y$};
    \draw[->,black!75,line width=0.6pt] (0,0) -- (2.0,0.9) node[lab,above=1pt]{$\mathbf v$};
  \end{scope}
  \draw[cS,->,line width=0.7pt] ($(C)+(245:0.7)$) arc (245:305:0.7);
  \node[lab,text=cS] at ($(C)+(275:0.92)$) {$r$};
\end{tikzpicture}
\end{figside}
\begin{eqside}
\[\begin{aligned}
\xS{\dot X} &= \xS{v_x}\cos\xS{\varphi} - \xS{v_y}\sin\xS{\varphi}\\
\xS{\dot Y} &= \xS{v_x}\sin\xS{\varphi} + \xS{v_y}\cos\xS{\varphi}\\
\xS{\dot\varphi} &= \xS{r}
\end{aligned}\]
\deps{$(\dot X,\dot Y)\leftarrow \xS{v_x,v_y,\varphi}$;\quad $\dot\varphi\leftarrow\xS{r}$}
\restr{pose = integral of the body velocity rotated by the heading.}
\end{eqside}
\end{panel}

% ---------------------------------------------------------------- A2
\begin{panel}{A2}{Slip angles}
\begin{figside}
\begin{tikzpicture}
  \coordinate (R) at (0,0); \coordinate (G) at (1.5,0); \coordinate (F) at (3.2,0);
  \draw[black!40,line width=1.4pt] (R) -- (F);
  \draw[guide] (F) -- ++(0:1.7);
  \draw[guide] ($(F)+(208:0.9)$) -- ++(28:2.7);
  \node[whl] at (R) {};
  \node[whl,rotate=28] at (F) {};
  \fill (G) circle (1.6pt);
  % rear
  \draw[vS] (R) -- ++(-16:1.55) node[sm,below=2pt]{$(v_x,\ v_y{-}l_r r)$};
  \draw[cD,line width=0.8pt] ($(R)+(-16:1.1)$) arc (-16:0:1.1);
  \node[lab,text=cD] at ($(R)+(-6:1.35)$) {$\alpha_r$};
  % front
  \draw[vS] (F) -- ++(10:1.6) node[sm,right]{$(v_x,\ v_y{+}l_f r)$};
  \draw[cD,line width=0.8pt] ($(F)+(10:1.25)$) arc (10:28:1.25);
  \node[lab,text=cD] at ($(F)+(19:1.5)$) {$\alpha_f$};
  \draw[cS,line width=0.6pt] ($(F)+(180:0.75)$) arc (180:208:0.75);
  \node[lab,text=cS] at ($(F)+(194:1.0)$) {$\delta$};
  % dims
  \draw[dim] (0,-1.05) -- (1.5,-1.05) node[sm,midway,below]{$l_r$};
  \draw[dim] (1.5,-1.05) -- (3.2,-1.05) node[sm,midway,below]{$l_f$};
\end{tikzpicture}
\end{figside}
\begin{eqside}
\[\begin{aligned}
\xD{\alpha_f} &= \xS{\delta} - \arctan\frac{\xS{v_y}+\xP{l_f}\,\xS{r}}{\xS{v_x}}\\[2pt]
\xD{\alpha_r} &= -\arctan\frac{\xS{v_y}-\xP{l_r}\,\xS{r}}{\xS{v_x}}
\end{aligned}\]
\deps{$\alpha_f\leftarrow\xS{v_x,v_y,r,\delta}$;\quad $\alpha_r\leftarrow\xS{v_x,v_y,r}$}
\restr{nothing alone: angle between where a tire points and where it moves (bounded in A8).}
\end{eqside}
\end{panel}

% ---------------------------------------------------------------- A3
\begin{panel}{A3}{Tire and drivetrain forces}
\begin{figside}[0.56]
\begin{tikzpicture}
  \begin{scope}[xscale=1.35,yscale=7]
    \fill[cL!10] (0.6,-0.215) rectangle (1.65,0.215);
    \fill[cL!10] (-1.65,-0.215) rectangle (-0.6,0.215);
    \draw[ax] (-1.7,0) -- (1.8,0) node[sm,right]{$\alpha$};
    \draw[ax] (0,-0.22) -- (0,0.25) node[sm,above]{$F_y$};
    \draw[guide] (-1.65,0.192) -- (1.65,0.192);
    \node[sm,text=cP,anchor=south east] at (-0.05,0.193) {$D_f$};
    \draw[black!50,line width=0.5pt] (-0.15,-0.0891) -- (0.26,0.1545);
    \node[sm,text=cP,anchor=east] at (-0.05,0.12) {$B_fC_fD_f$};
    \draw[cD,line width=1pt,domain=-1.6:1.6,samples=100,smooth] plot (\x,{0.192*sin(1.2*atan(2.579*\x))});
    \draw[liml] (0.6,-0.215) -- (0.6,0.215);
    \draw[liml] (-0.6,-0.215) -- (-0.6,0.215);
    \node[sm,text=cL,above] at (0.6,0.215) {$\alpha_{\max}$};
  \end{scope}
  \begin{scope}[shift={(3.35,-0.65)},xscale=0.6,yscale=7]
    \draw[ax] (0,0) -- (3.9,0) node[sm,right]{$v_x$};
    \draw[ax] (0,-0.09) -- (0,0.34) node[sm,above]{$F_x$};
    \draw[cD,line width=1pt,domain=0:3.5] plot (\x,{0.287-0.0545*\x}) node[sm,right]{$D{=}1$};
    \draw[cD!50,line width=1pt,domain=0:3.5] plot (\x,{0.5*(0.287-0.0545*\x)}) node[sm,right]{$D{=}\tfrac12$};
    \draw[cD,dashed,line width=0.8pt,domain=0:3.5] plot (\x,{-0.0518-0.00035*\x*\x}) node[sm,right]{$F_{\mathrm{fric}}$};
  \end{scope}
\end{tikzpicture}
\end{figside}
\begin{eqside}[0.42]
\[\begin{aligned}
\xD{F_{f,y}} &= \xP{D_f}\sin\!\big(\xP{C_f}\arctan(\xP{B_f}\,\xD{\alpha_f})\big)\\
\xD{F_{r,y}} &= \xP{D_r}\sin\!\big(\xP{C_r}\arctan(\xP{B_r}\,\xD{\alpha_r})\big)\\
\xD{F_{r,x}} &= (\xP{C_{m1}}-\xP{C_{m2}}\,\xS{v_x})\,\xS{D}\\
\xD{F_{\mathrm{fric}}} &= -\xP{C_{r0}}-\xP{C_{r2}}\,\xS{v_x}^2
\end{aligned}\]
{\footnotesize $\xP{D_f},\xP{D_r}$: peak force $\approx\mu F_z$ \quad $\xS D$: duty cycle $\in[-0.1,1]$}
\deps{$F_y\leftarrow\xD{\alpha}$;\quad $F_{r,x}\leftarrow\xS{D,v_x}$;\quad $F_{\mathrm{fric}}\leftarrow\xS{v_x}$}
\restr{lateral force saturates at $D$; drive force fades with speed; front tire has no $F_x$.}
\end{eqside}
\end{panel}

% ---------------------------------------------------------------- A4
\begin{panel}{A4}{Newton--Euler in the body frame}
\begin{figside}
\begin{tikzpicture}
  \coordinate (R) at (0,0); \coordinate (G) at (1.6,0); \coordinate (F) at (3.4,0);
  \draw[black!40,line width=1.4pt] (R) -- (F);
  \node[whl] at (R) {};
  \node[whl,rotate=22] at (F) {};
  \fill (G) circle (1.6pt);
  \draw[vD] (R) -- ++(0,1.0) node[lab,above]{$F_{r,y}$};
  \draw[vD] (-1.1,0) -- (-0.33,0);
  \node[lab,text=cD,below] at (-0.75,-0.05) {$F_{r,x}$};
  \draw[vD] (F) -- ++(112:1.05) node[lab,above]{$F_{f,y}$};
  \draw[vD,dashed] (G) -- ++(-0.75,0);
  \node[lab,text=cD,above] at (1.15,0.05) {$F_{\mathrm{fric}}$};
  \draw[cS,->,line width=0.7pt] ($(G)+(-150:0.5)$) arc (-150:-30:0.5);
  \node[lab,text=cS] at ($(G)+(-90:0.72)$) {$r$};
  \draw[dim] (0,-1.15) -- (1.6,-1.15) node[sm,midway,below]{$l_r$};
  \draw[dim] (1.6,-1.15) -- (3.4,-1.15) node[sm,midway,below]{$l_f$};
\end{tikzpicture}
\end{figside}
\begin{eqside}
\[\begin{aligned}
\xP m\,\xS{\dot v_x} &= \xD{F_{r,x}}+\xD{F_{\mathrm{fric}}}-\xD{F_{f,y}}\sin\xS\delta+\xP m\,\xS{v_y}\,\xS r\\
\xP m\,\xS{\dot v_y} &= \xD{F_{r,y}}+\xD{F_{f,y}}\cos\xS\delta-\xP m\,\xS{v_x}\,\xS r\\
\xP{I_z}\,\xS{\dot r} &= \xD{F_{f,y}}\,\xP{l_f}\cos\xS\delta-\xD{F_{r,y}}\,\xP{l_r}
\end{aligned}\]
{\footnotesize $\pm m\,v\,r$ terms: the body frame itself rotates.}
\deps{$(\dot v_x,\dot v_y,\dot r)\leftarrow\xD{F}$'s, $\xS{\delta}$, $\xS{v_x,v_y,r}$}
\restr{accelerations = forces / inertia. Linearized + ZOH-discretized each QP (A9).}
\end{eqside}
\end{panel}

% ---------------------------------------------------------------- A5
\begin{panel}{A5}{Progress $s$, contouring error $e_c$, lag error $e_l$}
\begin{figside}
\begin{tikzpicture}
  \begin{scope}[shift={(0,-4.4)}]
    \draw[refl] (127:5) arc (127:58:5) node[sm,right]{$\mathbf p_{\mathrm{ref}}(\cdot)$};
    \coordinate (P) at (102:5);
    \coordinate (S) at (88:5);
    \coordinate (p) at (88:4.15);
    \coordinate (N) at ($(P)+(-78:1)$);
    \coordinate (Q) at ($(P)!(p)!(N)$);
    \draw[guide] ($(P)+(192:0.6)$) -- ($(P)+(12:2.9)$);
    \draw[vec,black!70] (P) -- ++(12:0.65) node[sm,above]{$\mathbf t$};
    \draw[black!45,dotted,line width=0.8pt] (p) -- (S);
    \draw[cD,line width=1.2pt] (P) -- (Q) node[lab,midway,left=1pt]{$e_c$};
    \draw[cD,line width=1.2pt] (Q) -- (p) node[lab,midway,below=1pt]{$e_l$};
    \caricon{(p)}{-2}{0.65}
    \node[dot,cR] at (P) {};
    \node[sm,text=cR,above left] at (P) {$s$};
    \node[odot,draw=cR] at (S) {};
    \node[sm,text=cR,above] at (S) {$s^{\ast}$};
    \node[sm,below right=4pt] at (p) {$(X,Y)$};
  \end{scope}
\end{tikzpicture}
\end{figside}
\begin{eqside}
\[\begin{aligned}
\xD{e_c} &= \phantom{-}\xR{\mathbf n(\xS s)}^{\!\top}\big(\xS{\mathbf p}-\xR{\mathbf p_{\mathrm{ref}}(\xS s)}\big)\\
\xD{e_l} &= -\xR{\mathbf t(\xS s)}^{\!\top}\big(\xS{\mathbf p}-\xR{\mathbf p_{\mathrm{ref}}(\xS s)}\big)
\end{aligned}\]
{\footnotesize $\xR{\mathbf t}=(\cos\Phi,\sin\Phi)$,\ \ $\xR{\mathbf n}=(\sin\Phi,-\cos\Phi)$ (right),\ \ $\xS{\mathbf p}=(\xS X,\xS Y)$}
\deps{$e_c,e_l\leftarrow\xS{X,Y,s}$ and the spline $\xR{\mathbf p_{\mathrm{ref}}}$}
\restr{nothing alone; $e_l\!\approx\!0 \Leftrightarrow s\!\approx\!s^{\ast}$ (projection), with no inner $\arg\min$.}
\end{eqside}
\end{panel}

% ---------------------------------------------------------------- A6
\begin{panel}{A6}{Stage cost as springs}
\begin{figside}
\begin{tikzpicture}
  \draw[refl] (-0.2,0) -- (5.6,0);
  \coordinate (P) at (3.5,0); \coordinate (Fp) at (1.9,0); \coordinate (p) at (1.9,1.15);
  \draw[cD,line width=1.1pt,decorate,decoration={zigzag,segment length=3.2pt,amplitude=2.6pt}] (Fp) -- (P);
  \draw[cD!75,line width=0.5pt,decorate,decoration={coil,aspect=0.45,segment length=6pt,amplitude=3pt}] (Fp) -- (p);
  \caricon{(p)}{0}{0.8}
  \node[dot,cR] at (P) {};
  \node[sm,text=cR,below=2pt] at (P) {$s$};
  \draw[vS,line width=1.2pt] (P) -- ++(1.15,0);
  \node[lab,text=cS,above] at ($(P)+(0.62,0.05)$) {$q_{v_s}v_s$};
  \node[lab,text=cD,below] at ($(Fp)!0.5!(P)+(0,-0.12)$) {$q_l\,e_l^2$\ \ \textcolor{cP}{\scriptsize$q_l{=}500$}};
  \node[lab,text=cD,left=5pt] at ($(Fp)!0.5!(p)+(0,0.12)$) {$q_c\,e_c^2$};
  \node[sm,text=cP,left=5pt] at ($(Fp)!0.5!(p)+(0,-0.22)$) {$q_c{=}0.1$};
\end{tikzpicture}
\end{figside}
\begin{eqside}
\[J=\sum_{k=0}^{N}\Big(\xP{q_c}\xD{e_c}^2+\xP{q_l}\xD{e_l}^2-\xP{q_{v_s}}\xS{v_s}\Big)+\text{reg.}\]
{\footnotesize reg.\ $= \xP{q_\beta}\beta^2+\xP{q_r}\xS r^2+\xP{q_\mu}(\xS\varphi-\xR\Phi)^2+\|(\xS{D,\delta,v_s})\|^2_{\xP{R}}+\|\xU{\mathbf u}\|^2_{\xP{R_\Delta}}$\\
\phantom{reg.\ }$\beta=\arctan(\xS{v_y}/\xS{v_x})$;\ \ $q_c,q_r\times10$ at $k=N$;\ \ $q_{v_s}{=}0.02$}
\deps{$e_c,e_l$ (A5), $\xS{v_s}$, $\xS{v_x,v_y,r}$, $\xU{\mathbf u}$}
\restr{soft only: $s$ pinned to the car (stiff), lateral position nearly free (weak), progress rewarded.}
\end{eqside}
\end{panel}

% ---------------------------------------------------------------- A7
\begin{panel}{A7}{Track constraint (linearized slab)}
\begin{figside}
\begin{tikzpicture}
  \begin{scope}[shift={(0,-4.3)}]
    \coordinate (P) at (95:4.5);
    \coordinate (Ao) at (95:5.25);
    \coordinate (Ai) at (95:3.75);
    \path ($(Ao)+(185:2.6)$) coordinate (o1) ($(Ao)+(5:2.6)$) coordinate (o2)
          ($(Ai)+(185:2.6)$) coordinate (i1) ($(Ai)+(5:2.6)$) coordinate (i2);
    \fill[cS!7] (o1) -- (o2) -- (i2) -- (i1) -- cycle;
    \begin{scope}
      \clip (o1) -- (o2) -- (i2) -- (i1) -- cycle;
      \fill[cL!20,even odd rule] (0,0) circle (7) (0,0) circle (5.25);
    \end{scope}
    \draw[bnd] (128:5.25) arc (128:60:5.25);
    \draw[bnd] (128:3.75) arc (128:60:3.75);
    \draw[cR,dash pattern=on 4pt off 2pt,line width=0.8pt] (128:4.5) arc (128:60:4.5);
    \draw[liml] (o1) -- (o2);
    \draw[liml] (i1) -- (i2);
    \draw[vec,black!70] (P) -- ++(95:0.5) node[sm,right]{$\bar{\mathbf n}$};
    \node[dot,cR] at (P) {};
    \node[sm,text=cR,below right] at (P) {$\bar s_k$};
    \draw[dim] ($(P)+(185:1.7)$) -- ($(Ao)+(185:1.7)$) node[sm,midway,left]{$R_{\mathrm{out}}$};
    \draw[dim] ($(P)+(185:1.7)$) -- ($(Ai)+(185:1.7)$) node[sm,midway,left]{$R_{\mathrm{in}}$};
    \node[dot,cS] at ($(P)+(5:1.3)+(95:-0.42)$) {};
    \node[sm,text=cS,below] at ($(P)+(5:1.3)+(95:-0.42)$) {$\mathbf p_k$};
    \node[sm,text=cL] at ($(o2)+(-0.15,-0.35)$) {\tiny lin.\ error};
  \end{scope}
\end{tikzpicture}
\end{figside}
\begin{eqside}
\[-\xL{R_{\mathrm{in}}}\ \le\ \xR{\bar{\mathbf n}}^{\!\top}\big(\xS{\mathbf p_k}-\xR{\mathbf p_{\mathrm{ref}}(\bar s_k)}\big)\ \le\ \xL{R_{\mathrm{out}}}\quad(+\,\xi_k\ge 0)\]
{\footnotesize $\bar s_k$: progress of the \emph{previous guess} (fixed in this QP); $\xR{\bar{\mathbf n}}$: left normal there;\\
$R_{\mathrm{in}}{=}R_{\mathrm{out}}{=}0.14$\,m constant; car treated as a point.}
\deps{$\xS{X_k,Y_k}$ only}
\restr{car point inside a straight slab; soft (slack $\xi$, penalized); can overshoot the real edge far from $\bar s_k$.}
\end{eqside}
\end{panel}

% ---------------------------------------------------------------- A8
\begin{panel}{A8}{Tire constraints}
\begin{figside}
\begin{tikzpicture}[scale=0.72]
  \fill[cS!7] (0,0) ellipse (1.83 and 1.65);
  \draw[ax] (-2.3,0) -- (2.9,0) node[sm,right]{$F_{r,x}$};
  \draw[ax] (0,-2.0) -- (0,2.3) node[sm,above]{$F_{r,y}$};
  \draw[guide] (0,0) circle (1.737);
  \draw[liml] (0,0) ellipse (1.83 and 1.65);
  \draw[vD] (0,0) -- (0.9,1.2);
  \node[dot,cD] at (0.9,1.2) {};
  \node[sm,text=cD,right=2pt] at (0.9,1.2) {$(F_{r,x},F_{r,y})$};
  \node[sm,text=cL,anchor=north west] at (1.86,-0.3) {$e_\varepsilon D_r/e_{\mathrm{long}}$};
  \node[sm,text=cL,anchor=south east] at (-0.05,1.68) {$e_\varepsilon D_r$};
  \node[sm,text=black!55,anchor=north east] at (-1.25,-1.25) {$D_r$};
\end{tikzpicture}
\end{figside}
\begin{eqside}
\[\Big(\tfrac{\xP{e_{\mathrm{long}}}\,\xD{F_{r,x}}}{\xP{F_{N,r}}}\Big)^2+\Big(\tfrac{\xD{F_{r,y}}}{\xP{F_{N,r}}}\Big)^2\ \le\ \Big(\tfrac{\xL{e_\varepsilon}\,\xP{D_r}}{\xP{F_{N,r}}}\Big)^2\]
\[|\xD{\alpha_f}|\ \le\ \xL{\alpha_{\max}}\]
{\footnotesize $e_{\mathrm{long}}{=}0.9$,\ $e_\varepsilon{=}0.95$,\ $\alpha_{\max}{=}0.6$; $F_{N,r}$ static load. Both first-order in $\xS{\mathbf x}$ around the guess.}
\deps{ellipse $\leftarrow\xS{v_x,v_y,r,D}$;\quad $|\alpha_f|\leftarrow\xS{v_x,v_y,r,\delta}$}
\restr{rear: drive + lateral force share one ellipse at 95\,\% of peak. Front: slip below $\alpha_{\max}$ (A3, red band). Soft.}
\end{eqside}
\end{panel}

% ---------------------------------------------------------------- A9
\begin{panel}{A9}{One QP per SQP step, trust region on $s$}
\begin{figside}
\begin{tikzpicture}[x=0.52cm,y=1cm]
  \fill[cL!12] (0,0.0) -- (11,2.18) -- (11,2.68) -- (0,0.5) -- cycle;
  \draw[ax] (-0.3,0) -- (12.4,0) node[sm,right]{$k$};
  \draw[ax] (0,-0.1) -- (0,2.8) node[sm,left]{$s_k$};
  \draw[black!40,line width=0.5pt] (0,0.25) -- (11,2.43);
  \foreach \k in {1,...,7,11} {
    \pgfmathsetmacro\sb{0.25+0.198*\k}
    \pgfmathsetmacro\so{\sb+0.17*sin(40*\k)}
    \node[odot,draw=black!45] at (\k,\sb) {};
    \node[dot,cS] at (\k,\so) {};
  }
  \node[dot,cS] at (0,0.25) {};
  \node[sm,left] at (0,0.25) {$\hat{\mathbf x}$};
  \node[sm] at (9,1.4) {$\cdots$};
  \foreach \k/\l in {0/0,1/1,2/2,11/N} \node[sm,below] at (\k,-0.05) {$\l$};
  \draw[decorate,decoration={brace,amplitude=4pt},black!50] (0.2,3.05) -- (11,3.05) node[sm,midway,above=4pt]{$N\,T_s=60\times20\,\text{ms}$};
  \node[sm,text=cL] at (6.3,2.55) {$|s_k-\bar s_k|\le 0.2$\,m};
  \node[sm,text=black!50] at (8.6,0.95) {$\bar s_k$ (guess)};
\end{tikzpicture}
\end{figside}
\begin{eqside}
\[\begin{aligned}
\min\ &\textstyle\sum_k \tfrac12\xS{\mathbf x_k}^{\!\top}\!Q_k\xS{\mathbf x_k}+\mathbf q_k^{\!\top}\xS{\mathbf x_k}+\tfrac12\xU{\mathbf u_k}^{\!\top}\!R\,\xU{\mathbf u_k}+\text{slack}\\
\text{s.t. }&\xS{\mathbf x_0}=\hat{\mathbf x},\quad \xS{\mathbf x_{k+1}}=A_k\xS{\mathbf x_k}+B_k\xU{\mathbf u_k}+\mathbf g_k\\
&\text{A7, A8 (soft)},\ \ \text{box bounds on }\xS{\mathbf x},\xU{\mathbf u}\\
&|\xS{s_k}-\bar s_k|\le 0.2\,\text{m}
\end{aligned}\]
{\footnotesize $(A_k,B_k,\mathbf g_k)$: exact ZOH of the Jacobian at $(\bar{\mathbf x}_k,\bar{\mathbf u}_k)$.\\
$\bar{\mathbf x}\leftarrow 0.8\,\mathbf x^{\ast}+0.2\,\bar{\mathbf x}$, \ 2 iterations, apply $\xU{\mathbf u_0}$, shift.}
\restr{each QP trusts its linearization only near the previous guess.}
\end{eqside}
\end{panel}

% =====================================================================
\clearpage
\sect{Offline \quad TUM \texttt{global\_racetrajectory\_optimization} \quad (whole lap)}

% ---------------------------------------------------------------- B0
\begin{panel}{B0}{Two pipelines}
\centering
\begin{tikzpicture}[
  bx/.style={draw=black!35,rounded corners=2pt,inner sep=3.5pt,font=\footnotesize,align=center,fill=#1!9},
  fl/.style={->,black!55,line width=0.6pt}]
  \node[font=\sffamily\small\bfseries,anchor=west] at (-1.8,0.95) {mincurv};
  \node[font=\footnotesize\itshape,text=black!50,anchor=west] at (-0.2,0.95) {path first, speed second};
  \node[bx=cR] (t1) at (0,0) {$\xR{\mathbf p_i,\ \mathbf n_i}$\\ $\xR{w_{\ell,i},\ w_{r,i}}$};
  \node[bx=cS] (q1) at (3.4,0) {QP over $\xS{\boldsymbol\alpha}$\\ $\min\sum\kappa_i^2$};
  \node[bx=cD] (r1) at (6.6,0) {$\xD{\mathbf r(s)},\ \xD{\kappa(s)}$};
  \node[bx=cS] (v1) at (10.2,0) {fwd/bwd pass\\ $\xS{v(s)},\ \xD{a_x(s)}$};
  \node[bx=cP] (g1) at (10.2,-1.3) {$\text{ggv}(v),\ a_{x,\mathrm{mach}}(v),\ c_d,\ m,\ \mu$};
  \draw[fl] (t1) -- (q1); \draw[fl] (q1) -- (r1); \draw[fl] (r1) -- (v1); \draw[fl] (g1) -- (v1);
  \draw[fl,dashed] (r1.north) to[bend right=35] node[sm,above]{IQP: re-linearize} (q1.north);
  \node[font=\sffamily\small\bfseries,anchor=west] at (-1.8,-2.35) {mintime};
  \node[font=\footnotesize\itshape,text=black!50,anchor=west] at (-0.2,-2.35) {path and speed in one problem};
  \node[bx=cR] (t2) at (0,-3.4) {$\xR{\kappa_{\mathrm{ref}}(s)}$\\ $\xR{w_\ell(s),\ w_r(s)}$};
  \node[bx=cS] (n2) at (5.0,-3.4) {NLP (collocation + IPOPT)\\ $\min\int_0^L \tfrac{dt}{ds}\,ds$\\ over $\xS{\mathbf x(s)},\ \xU{\mathbf u(s)}$};
  \node[bx=cD] (o2) at (10.2,-3.4) {$\xD{\mathbf r(s)},\ \xS{v(s)}$};
  \node[bx=cP] (p2) at (5.0,-4.95) {double track, Pacejka, load transfer, aero, $\mu$};
  \draw[fl] (t2) -- (n2); \draw[fl] (n2) -- (o2); \draw[fl] (p2) -- (n2);
\end{tikzpicture}
\restr{\texttt{mincurv}: speed never feeds back into geometry. \texttt{mintime}: geometry and speed traded off jointly.}
\end{panel}

% ---------------------------------------------------------------- B1
\begin{panel}{B1}{mincurv: the only variable is a lateral shift $\alpha_i$}
\begin{figside}
\begin{tikzpicture}
  \begin{scope}[shift={(0,-5.65)}]
    \draw[bnd] (121:7) arc (121:59:7);
    \draw[bnd] (121:5) arc (121:59:5);
    \draw[refl,dash pattern=on 4pt off 2pt] (121:6) arc (121:59:6);
    \draw[cS,line width=1pt] plot[smooth] coordinates {(121:6.62) (115:6.6) (105:6.15) (95:5.45) (85:5.45) (75:6.15) (65:6.6) (59:6.62)};
    \foreach \t/\a in {115/-0.6,105/-0.15,95/0.55,85/0.55,75/-0.15,65/-0.6} {
      \draw[->,black!40,line width=0.5pt] (\t:6) -- (\t:5.5);
      \node[dot,cR] at (\t:6) {};
      \pgfmathsetmacro\rr{6-\a}
      \node[dot,cS] at (\t:\rr) {};
    }
    \draw[cL!60,line width=3pt,opacity=0.5] (77.5:6.9) -- (77.5:5.1);
    \draw[cL,line width=0.6pt] ($(77.5:6.9)+(167.5:0.12)$) -- ($(77.5:6.9)+(-12.5:0.12)$);
    \draw[cL,line width=0.6pt] ($(77.5:5.1)+(167.5:0.12)$) -- ($(77.5:5.1)+(-12.5:0.12)$);
    \node[sm,text=cL,right=3pt] at (77.5:5.25) {$[-\bar w_\ell,\ \bar w_r]$};
    \node[sm,text=cR,above] at (95:6.02) {$\mathbf p_i$};
    \node[sm,text=cS,below] at (95:5.42) {$\mathbf r_i$};
    \node[sm,text=black!55,left] at (115:5.6) {$\mathbf n_i$};
    \node[lab,text=cS,right] at (91.5:5.75) {$\alpha_i$};
    \draw[dim] (119:6) -- (119:7) node[sm,midway,left]{$w_\ell$};
    \draw[dim] (119:6) -- (119:5) node[sm,midway,left]{$w_r$};
    \draw[->,cR,line width=0.8pt] (63:6) arc (63:58:6);
  \end{scope}
\end{tikzpicture}
\end{figside}
\begin{eqside}
\[\xD{\mathbf r_i}=\xR{\mathbf p_i}+\xS{\alpha_i}\,\xR{\mathbf n_i}\]
\[-\big(\xR{w_{\ell,i}}-\tfrac{\xP{w_{\mathrm{veh}}}}{2}\big)\ \le\ \xS{\alpha_i}\ \le\ \xR{w_{r,i}}-\tfrac{\xP{w_{\mathrm{veh}}}}{2}\]
{\footnotesize $\xR{\mathbf n_i}$ points \emph{right};\quad $\bar w=w-w_{\mathrm{veh}}/2$;\quad decision vector $\xS{\boldsymbol\alpha}\in\mathbb R^{N}$.}
\deps{$\mathbf r_i\leftarrow\xS{\alpha_i}$ and fixed track data}
\restr{each point slides only along its own normal, inside the track minus half the car width.}
\end{eqside}
\end{panel}

% ---------------------------------------------------------------- B2
\begin{panel}{B2}{mincurv: curvature, linearized}
\begin{figside}
\begin{tikzpicture}[scale=1.05]
  \draw[cR,line width=0.8pt,dash pattern=on 4pt off 2pt,domain=-2.2:2.5,samples=50,smooth] plot (\x,{0.22*\x*\x-0.75});
  \draw[cS,line width=1pt,domain=-2.2:2.5,samples=50,smooth] plot (\x,{0.22*\x*\x});
  \coordinate (Ar) at (1.2,-0.4332);
  \coordinate (A) at (1.2,0.3168);
  \draw[vec,black!55] (Ar) -- ++(0.95,0.5016) node[sm,right]{$(\xP{x'},\xP{y'})$ frozen};
  \node[dot,cR] at (Ar) {};
  \draw[vec,cS] (A) -- ++(0,0.85) node[sm,above]{$(\xS{x''},\xS{y''})$};
  \draw[guide] (A) -- ++(-0.7004,1.3265) node[sm,left]{$1/\kappa_i$};
  \draw[cD,dashed,line width=0.6pt] (-0.335,3.224) ++(-98:3.287) arc (-98:-26:3.287);
  \node[dot,cS] at (A) {};
  \node[sm,text=cS,right=2pt] at (A) {$\mathbf r_i$};
  \node[sm,text=cR] at (-1.6,-0.05) {ref.};
  \node[sm,text=cS] at (-1.75,1.25) {raceline};
\end{tikzpicture}
\end{figside}
\begin{eqside}
\[\xD{\kappa_i}=\frac{\xP{x_i'}\,\xS{y_i''}-\xP{y_i'}\,\xS{x_i''}}{\big(\xP{x_i'}^2+\xP{y_i'}^2\big)^{3/2}}
\ \Rightarrow\ \xD{\boldsymbol\kappa}\approx\xR{\boldsymbol\kappa_{\mathrm{ref}}}+\xP{E_\kappa}\,\xS{\boldsymbol\alpha}\]
\[\min_{\xS{\boldsymbol\alpha}}\ \sum_i\xD{\kappa_i}^2=\tfrac12\xS{\boldsymbol\alpha}^{\!\top}\!H\xS{\boldsymbol\alpha}+\mathbf f^{\top}\xS{\boldsymbol\alpha}
\quad\text{s.t.}\ \ |\xD{\kappa_i}|\le\xL{\kappa_{\max}}\ \ \&\ \text{B1}\]
{\footnotesize $x',y'$ taken from the reference spline (frozen); $x'',y''$ are linear in $\boldsymbol\alpha$ through the spline system. IQP: refreeze at the new line, re-solve.}
\deps{$\kappa_i\leftarrow\xS{\boldsymbol\alpha}$ (neighbours too, via splines)}
\restr{a QP: smoothest line through the corridor; no speed, mass or tire in it.}
\end{eqside}
\end{panel}

% ---------------------------------------------------------------- B3
\begin{panel}{B3}{Velocity profile: acceleration left after cornering}
\begin{figside}
\begin{tikzpicture}
  \begin{scope}
    \clip (-2.0,-1.6) rectangle (2.0,0.8);
    \fill[cS!7] (0,0) ellipse (1.6 and 1.3);
  \end{scope}
  \draw[ax] (-2.1,0) -- (2.4,0) node[sm,right]{$a_y$};
  \draw[ax] (0,-1.7) -- (0,1.85) node[sm,above]{$a_x$};
  \draw[cL,line width=0.9pt] (0,0) ellipse (1.6 and 1.3);
  \draw[cL!70,dashed,line width=0.7pt] (1.6,0) -- (0,1.3) -- (-1.6,0) -- (0,-1.3) -- cycle;
  \draw[liml] (-1.9,0.8) -- (1.9,0.8) node[sm,right]{$a_{x,\mathrm{mach}}(v)$};
  \draw[cD,line width=0.8pt] (0.75,-1.6) -- (0.75,1.6) node[sm,above]{$a_y{=}v^2\kappa$};
  \node[odot,draw=cD] at (0.75,1.148) {};
  \node[dot,cD] at (0.75,-1.148) {};
  \node[dot,cD] at (0.75,0.8) {};
  \node[sm,text=cL,anchor=north west] at (1.62,-0.18) {$\mu a_{y,\max}$};
  \node[sm,text=cL,above left] at (0,1.3) {$\mu a_{x,\max}$};
  \node[sm,text=cL] at (-1.45,1.05) {$p{=}2$};
  \node[sm,text=cL!70] at (-0.55,-0.6) {$p{=}1$};
\end{tikzpicture}
\end{figside}
\begin{eqside}
\[\xD{a_y}=\xS v^2\,\xD{\kappa}\]
\[\xD{a_x^{\mathrm{avail}}}=\xL{\mu\,a_{x,\max}(v)}\Big[1-\Big(\frac{\xD{a_y}}{\xL{\mu\,a_{y,\max}(v)}}\Big)^{\xP p}\Big]^{1/\xP p}\]
\[\begin{aligned}
\text{accel: }&\min\!\big(\xD{a_x^{\mathrm{avail}}},\,\xL{a_{x,\mathrm{mach}}(v)}\big)-\tfrac{\xP{c_d}}{\xP m}\xS v^2\\
\text{brake: }&\xD{a_x^{\mathrm{avail}}}+\tfrac{\xP{c_d}}{\xP m}\xS v^2\quad(\text{magnitude})
\end{aligned}\]
{\footnotesize ggv table: $(v,a_{x,\max},a_{y,\max})$; $p$ = \texttt{dyn\_model\_exp} $\in[1,2]$ (default 1).}
\deps{$a_x\leftarrow\xS v,\ \xD\kappa$ (from B2)}
\restr{longitudinal accel.\ allowed by grip after cornering, and by the motor.}
\end{eqside}
\end{panel}

% ---------------------------------------------------------------- B4
\begin{panel}{B4}{Velocity profile: forward--backward pass}
\begin{figside}
\begin{tikzpicture}[x=0.72cm,y=0.36cm]
  \draw[ax] (0,0) -- (10.7,0) node[sm,right]{$s$};
  \draw[ax] (0,0) -- (0,7.2) node[sm,above]{$v$};
  \draw[cL,dash pattern=on 3pt off 2pt,line width=0.6pt,domain=0:10,samples=150,smooth]
     plot (\x,{6-4*exp(-(((\x-3)/0.6)^2))-3*exp(-(((\x-7.5)/0.6)^2))});
  \draw[cS!45,dashed,line width=0.7pt,domain=3:7.5,samples=40] plot (\x,{sqrt(4+5*(\x-3))});
  \draw[cS!45,dashed,line width=0.7pt,domain=7.5:10,samples=30] plot (\x,{sqrt(9+5*(\x-7.5))});
  \draw[cL!55,dashed,line width=0.7pt,domain=0.3:3,samples=30] plot (\x,{sqrt(4+8*(3-\x))});
  \draw[cL!55,dashed,line width=0.7pt,domain=4.2:7.5,samples=30] plot (\x,{sqrt(9+8*(7.5-\x))});
  \draw[cS,line width=1.3pt,domain=0:10,samples=200]
     plot (\x,{min(6-4*exp(-(((\x-3)/0.6)^2))-3*exp(-(((\x-7.5)/0.6)^2)),
                   (\x<3)*sqrt(4+8*abs(3-\x)) + (\x>=3)*(\x<7.5)*min(sqrt(4+5*abs(\x-3)),sqrt(9+8*abs(7.5-\x))) + (\x>=7.5)*sqrt(9+5*abs(\x-7.5)))});
  \node[sm,text=cL,anchor=south] at (5.2,6.05) {$v_\kappa=\sqrt{\mu a_{y,\max}/|\kappa|}$};
  \node[sm,text=cS!70] at (5.35,3.3) {fwd};
  \node[sm,text=cL!80] at (6.45,4.6) {bwd};
  \foreach \s in {3,7.5} \draw[black!40] (\s,-0.25) -- (\s,0.25);
  \node[sm,below] at (3,-0.2) {apex};
  \node[sm,below] at (7.5,-0.2) {apex};
\end{tikzpicture}
\end{figside}
\begin{eqside}
\[\text{fwd: }\ \xS{v_{i+1}}=\sqrt{\xS{v_i}^2+2\,\xD{a_x^{+}}(\xS{v_i},\xD{\kappa_i})\,\xR{\ell_i}}\]
\[\text{bwd: }\ \xS{v_{i}}=\sqrt{\xS{v_{i+1}}^2+2\,|\xD{a_x^{-}}(\xS{v_{i+1}},\xD{\kappa_{i+1}})|\,\xR{\ell_i}}\]
\[\xS{v}=\min\big(v_\kappa,\ \text{fwd},\ \text{bwd}\big)\]
{\footnotesize $\ell_i$: segment length; $a_x^{\pm}$ from B3. No optimizer: a sweep.}
\deps{$v\leftarrow\xD\kappa(s)$, ggv, motor, drag}
\restr{speed $\le$ cornering limit, and reachable by accelerating out of / braking into each apex.}
\end{eqside}
\end{panel}

% ---------------------------------------------------------------- B5
\begin{panel}{B5}{mintime: who depends on whom}
\begin{figside}[0.56]
\begin{tikzpicture}[x=0.55cm,y=0.37cm]
  \def\ncols{11}
  \mcol{1}{v}{cS}\mcol{2}{\beta}{cS}\mcol{3}{\omega_z}{cS}\mcol{4}{n}{cS}\mcol{5}{\xi}{cS}
  \mcol{6}{\delta}{cU}\mcol{7}{F_d}{cU}\mcol{8}{F_b}{cU}\mcol{9}{\gamma_y}{cU}
  \mcol{10}{\kappa}{cR}\mcol{11}{w}{cR}
  \mrow{1}{\xD{S_f}}{1,2,4,5,10}{0}
  \mrow{2}{\xD{\alpha_{ij}}}{1,2,3,6}{0}
  \mrow{3}{\xD{F_{z,ij}}}{1,7,8,9}{0}
  \mrow{4}{\xD{F_{x,ij}}}{7,8}{0}
  \mrow{5}{\xD{F_{y,ij}}}{0}{1,2,3,6,7,8,9}
  \mrow{6.4}{\xS{v'}}{2,6}{1,3,4,5,7,8,9,10}
  \mrow{7.4}{\xS{\beta'}}{1,2,3,6}{4,5,7,8,9,10}
  \mrow{8.4}{\xS{\omega_z'}}{6}{1,2,3,4,5,7,8,9,10}
  \mrow{9.4}{\xS{n'}}{1,2,5}{4,10}
  \mrow{10.4}{\xS{\xi'}}{3,10}{1,2,4,5}
  \mrow{11.8}{T_{\mathrm{lap}}}{0}{1,2,4,5,10}
  \mrow{12.8}{\xL{\text{track}}}{4,11}{0}
  \mrow{13.8}{\xL{\text{circles}}}{0}{1,2,3,6,7,8,9}
  \mrow{14.8}{\xL{\gamma_y\ \text{eq.}}}{9}{1,2,3,6,7,8}
  \mrow{15.8}{\xL{\text{power}}}{1,7}{0}
  \mrow{16.8}{\xL{F_dF_b}}{7,8}{0}
  \mrow{17.8}{\xL{\text{rates}}}{6,7,8}{1,4,10}
  \mrow{18.8}{\xL{\text{periodic}}}{1,2,3,4,5}{0}
  \draw[black!30] (0.5,-5.7) -- (11.5,-5.7);
  \draw[black!30] (0.5,-11.1) -- (11.5,-11.1);
  \draw[black!30] (5.5,1.4) -- (5.5,-19.3);
  \draw[black!30] (9.5,1.4) -- (9.5,-19.3);
\end{tikzpicture}
\end{figside}
\begin{eqside}[0.40]
{\footnotesize
$(\cdot)'=d(\cdot)/ds$\par\medskip
\tikz\fill[black!85] (0,0) circle (2.5pt);\ appears in the formula\par
\tikz\draw[black!60,line width=0.6pt] (0,0) circle (2.2pt);\ enters through $S_f$, $\alpha$ or $F$\par\medskip
Every $d/ds$ row inherits $n,\xi,\kappa$ through $S_f=dt/ds$.}
\restr{compared with A0: no $X,Y$ states and no progress variable; position is $(s,n)$ on the fixed reference.}
\end{eqside}
\end{panel}

% ---------------------------------------------------------------- B6
\begin{panel}{B6}{mintime: curvilinear coordinates, time as a function of $s$}
\begin{figside}
\begin{tikzpicture}
  \coordinate (O) at (0,0);
  \coordinate (P) at (0:4);
  \coordinate (C) at (0:2.85);
  \draw[refl] (-32:4) arc (-32:42:4) node[sm,above]{ref.};
  \node[dot,black!55] at (O) {};
  \draw[guide] (O) -- (-22:4) node[sm,pos=0.33,below,text=cR]{$1/\kappa(s)$};
  \node[dot,cR] at (P) {};
  \node[sm,text=cR,right] at (P) {$s$};
  \draw[guide] (2.85,-1.0) -- (2.85,2.4);
  \draw[guide] (4,-0.25) -- (4,-1.15);
  \draw[guide] (C) -- ++(108:2.3);
  \draw[dim] (2.85,-1.1) -- (4,-1.1) node[sm,midway,below,text=cS]{$n$};
  \begin{scope}[shift={(C)},rotate=108]
    \draw[body] (-0.72,-0.28) rectangle (0.72,0.28);
    \fill[black!55] (0.72,0) -- (0.52,0.12) -- (0.52,-0.12) -- cycle;
  \end{scope}
  \draw[vS] (C) -- ++(122:2.15) node[lab,above]{$v$};
  \draw[cS,line width=0.6pt] ($(C)+(90:1.3)$) arc (90:108:1.3);
  \node[lab,text=cS] at ($(C)+(99:1.08)$) {$\xi$};
  \draw[cS,line width=0.6pt] ($(C)+(108:1.75)$) arc (108:122:1.75);
  \node[lab,text=cS] at ($(C)+(115:2.0)$) {$\beta$};
  \draw[cS,->,line width=0.7pt] ($(C)+(195:0.95)$) arc (195:255:0.95);
  \node[lab,text=cS] at ($(C)+(222:1.22)$) {$\omega_z$};
\end{tikzpicture}
\end{figside}
\begin{eqside}
\[\xD{S_f}\equiv\frac{dt}{ds}=\frac{1-\xS n\,\xR\kappa}{\xS v\cos(\xS\xi+\xS\beta)}\]
\[\frac{d\xS n}{ds}=\xD{S_f}\,\xS v\sin(\xS\xi+\xS\beta),\qquad \frac{d\xS\xi}{ds}=\xD{S_f}\,\xS{\omega_z}-\xR\kappa\]
\[T_{\mathrm{lap}}=\int_0^{L}\xD{S_f}\,ds\]
{\footnotesize $\xS{\mathbf x}=(v,\beta,\omega_z,n,\xi)$,\ $\xU{\mathbf u}=(\delta,F_d,F_b,\gamma_y)$;\ \ $n>0$ = \emph{left} (opposite sign to $\alpha$ in B1).}
\deps{$S_f\leftarrow\xS{v,\beta,n,\xi},\ \xR\kappa$}
\restr{time is not a variable: lap time integrates $dt/ds$ over the fixed reference length $L$.}
\end{eqside}
\end{panel}

% ---------------------------------------------------------------- B7
\begin{panel}{B7}{mintime: double-track loads}
\begin{figside}
\begin{tikzpicture}
  \draw[body] (-0.55,-0.55) rectangle (3.85,0.55);
  \draw[black!35] (0,-0.8) -- (0,0.8);
  \draw[black!35] (3.3,-0.8) -- (3.3,0.8);
  \draw[cL!80,line width=0.6pt,fill=cL!6] (3.3,-0.8) circle (0.55);
  \draw[cL!80,line width=0.6pt,fill=cL!6] (3.3,0.8) circle (0.38);
  \draw[cL!80,line width=0.6pt,fill=cL!6] (0,-0.8) circle (0.40);
  \draw[cL!80,line width=0.6pt,fill=cL!6] (0,0.8) circle (0.24);
  \node[whl] at (0,0.8) {}; \node[whl] at (0,-0.8) {};
  \node[whl,rotate=15] at (3.3,0.8) {}; \node[whl,rotate=15] at (3.3,-0.8) {};
  \fill (1.54,0) circle (1.6pt);
  \draw[vD] (1.54,0) -- ++(-0.8,0) node[sm,left]{$a_x$};
  \draw[vD] (1.54,0) -- ++(0,0.5) node[sm,above]{$a_y$};
  \node[sm,text=cL] at (4.45,-0.8) {$\propto F_{z,fr}$};
  \draw[dim] (0,-1.55) -- (1.54,-1.55) node[sm,midway,below]{$l_r$};
  \draw[dim] (1.54,-1.55) -- (3.3,-1.55) node[sm,midway,below]{$l_f$};
  \draw[dim] (-1.0,-0.8) -- (-1.0,0.8) node[sm,midway,left]{$b$};
  \node[sm,text=black!50] at (1.65,1.15) {braking + left turn};
\end{tikzpicture}
\end{figside}
\begin{eqside}[0.46]
\[\xD{F_{z,ij}}=\xP{F^{\mathrm{st}}_{z,ij}}+\tfrac12\xP{c_{l,i}}\xS v^2
\mp\tfrac{\xP h}{2\xP l}\big(\xU{F_d}{+}\xU{F_b}{-}\xP{c_d}\xS v^2{-}\xP{F_{\mathrm{roll}}}\big)
\mp\xP{k_i}\,\xU{\gamma_y}\]
{\footnotesize long.: $-$front/$+$rear;\ lat.: $-$left/$+$right;\ $k_f{=}k_{\mathrm{roll}}$, $k_r{=}1{-}k_{\mathrm{roll}}$}
\[\xU{\gamma_y}=\frac{\xP h}{\xP{\bar b}}\,\xP m\,\xD{a_y}\quad(\text{equality constraint})\]
{\footnotesize $\xD{F_{x,ij}}$: drive/brake split $\xP{k_{\mathrm{drive}}},\xP{k_{\mathrm{brake}}}$ minus rolling resistance.}
\deps{$F_z\leftarrow\xS v,\ \xU{F_d,F_b,\gamma_y}$}
\restr{each wheel's load (hence grip) follows the accelerations; $\gamma_y$ is a decision variable closed by an equality.}
\end{eqside}
\end{panel}

% ---------------------------------------------------------------- B8
\begin{panel}{B8}{mintime: tire model and friction circles}
\begin{figside}
\begin{tikzpicture}
  \begin{scope}[xscale=7,yscale=1.6]
    \draw[ax] (0,0) -- (0.38,0) node[sm,right]{$\alpha$};
    \draw[ax] (0,0) -- (0,1.3) node[sm,above]{$F_y/D$};
    \draw[guide] (0,1) -- (0.36,1);
    \draw[cD,line width=1pt,domain=0:0.35,samples=150,smooth] plot (\x,{sin(2.5*atan(10*\x-1.0*(10*\x-rad(atan(10*\x)))))});
    \node[sm,left] at (0,1) {$1$};
  \end{scope}
  \begin{scope}[shift={(3.55,0)},xscale=0.25,yscale=0.25]
    \draw[ax] (0,0) -- (10,0) node[sm,right]{$F_z$};
    \draw[ax] (0,0) -- (0,9.6) node[sm,above]{$D$};
    \draw[guide] (0,0) -- (9,9) node[sm,above]{$\mu F_z$};
    \draw[cD,line width=1pt,domain=0:9,samples=30] plot (\x,{\x*(1-0.1*\x/3)});
    \node[sm,text=cD] at (8.2,4.9) {$\varepsilon<0$};
  \end{scope}
\end{tikzpicture}
\end{figside}
\begin{eqside}
\[\xD{D_{ij}}=\xP\mu\,\xD{F_{z,ij}}\Big(1+\xP\varepsilon\,\frac{\xD{F_{z,ij}}}{\xP{F_{z0}}}\Big)\]
\[\xD{F_{y,ij}}=\xD{D_{ij}}\sin\!\Big(\xP C\arctan\big(\xP B\xD\alpha-\xP E(\xP B\xD\alpha-\arctan\xP B\xD\alpha)\big)\Big)\]
\[\Big(\frac{\xD{F_{x,ij}}}{\xP\mu\,\xD{F_{z,ij}}}\Big)^2+\Big(\frac{\xD{F_{y,ij}}}{\xP\mu\,\xD{F_{z,ij}}}\Big)^2\le\xL{1}\qquad\forall\,ij\]
\deps{$\alpha_{ij}\leftarrow\xS{v,\beta,\omega_z},\xU\delta$;\ \ $F_{z,ij}$ from B7}
\restr{each tire inside its own circle sized by its load (B7 circles); load transfer costs total grip ($\varepsilon<0$).}
\end{eqside}
\end{panel}

% ---------------------------------------------------------------- B9
\begin{panel}{B9}{mintime: the NLP}
\begin{figside}[0.40]
\begin{tikzpicture}[x=1cm,y=1cm]
  \draw[ax] (-0.2,0) -- (5.3,0) node[sm,right]{$s$};
  \foreach \t in {0.1127,0.5,0.8873} \draw[black!25] ({0.4+4.2*\t},0) -- ({0.4+4.2*\t},2.2);
  \draw[cS,line width=1pt,domain=0:1,samples=40] plot ({0.4+4.2*\x},{1+2.4*\x-1.4*\x*\x});
  \foreach \t/\m in {0.1127/0.4963,0.5/0.2381,0.8873/-0.0201} {
    \pgfmathsetmacro\xx{0.4+4.2*\t}
    \pgfmathsetmacro\yy{1+2.4*\t-1.4*\t*\t}
    \draw[cD,line width=1.1pt] ({\xx-0.6},{\yy-0.6*\m+0.12}) -- ({\xx+0.6},{\yy+0.6*\m+0.12});
    \node[dot,cS] at (\xx,\yy) {};
  }
  \node[dot,cS] at (0.4,1) {};
  \node[odot,draw=cS] at (4.6,2) {};
  \node[sm,text=cS,left] at (0.4,1) {$X_k$};
  \node[sm,text=cS,right] at (4.6,2) {$X_{k+1}$};
  \node[sm,text=cD,above] at (2.5,2.3) {slope $=h_k\mathbf f(X_{k,j},U_k,\kappa)$};
  \draw[cU,line width=1pt] (-0.1,-0.75) -- (0.4,-0.75) (0.4,-0.5) -- (4.6,-0.5) (4.6,-0.85) -- (5.1,-0.85);
  \node[sm,text=cU,above] at (2.5,-0.5) {$U_k$};
  \draw[dim] (0.4,-1.15) -- (4.6,-1.15) node[sm,midway,below]{$h_k$};
\end{tikzpicture}
\end{figside}
\begin{eqside}[0.58]
\[\min_{\xS{\mathbf x},\xU{\mathbf u}}\ \int_0^{L}\xD{S_f}\,ds+\xP{\lambda_F}\|\Delta\xU F\|^2+\xP{\lambda_\delta}\|\Delta\xU\delta\|^2\]
\footnotesize
\[\begin{aligned}
&\text{dynamics:}  && h_k\,\mathbf f(X_{k,j},\xU{U_k},\xR{\kappa})=\textstyle\sum_r C_{rj}X_{k,r},\ \ 3\text{ Legendre pts}\\
&\text{track:}     && -\xR{w_r}+\tfrac{\xP{w_{\mathrm{opt}}}}{2}\le\xS n\le\xR{w_\ell}-\tfrac{\xP{w_{\mathrm{opt}}}}{2}\\
&\text{grip, loads:}&& \text{B8 circles},\ \ \gamma_y\ \text{equality (B7)}\\
&\text{power:}     && \xS v\,\xU{F_d}\le\xL{P_{\max}},\qquad -\epsilon\le\xU{F_d}\,\xU{F_b}\le 0\\
&\text{actuators:} && \tfrac{\Delta\xU\delta}{\Delta t}\in\pm\tfrac{\xL{\delta_{\max}}}{\xP{t_\delta}},\ \
                      \tfrac{\Delta\xU{F_d}}{\Delta t}\le\tfrac{\xL{F_{d,\max}}}{\xP{t_d}},\ \
                      \tfrac{\Delta\xU{F_b}}{\Delta t}\ge-\tfrac{\xL{F_{b,\max}}}{\xP{t_b}}\\
&\text{closed lap:}&& \xS{\mathbf x}(0)=\xS{\mathbf x}(L)\\
&\text{bounds:}    && \xS v\ge 1,\ \ |\xS\beta|,|\xS{\omega_z}|,|\xS\xi|\le\tfrac{\pi}{2},\ \ \xU{F_d}\ge0\ge\xU{F_b}
\end{aligned}\]
\small
\restr{dynamics enforced only at collocation points; the lap must close; everything else is a path constraint.}
\end{eqside}
\end{panel}
}
. Requires amsmath, amssymb, array.
% =====================================================================
\usepackage{xcolor}
\usepackage{tikz}
\usetikzlibrary{arrows.meta,calc,decorations.pathmorphing,decorations.pathreplacing}

% ---------------- role colours ----------------------------------------
\definecolor{cS}{HTML}{0072B2}  % decision state
\definecolor{cU}{HTML}{C76E00}  % control input
\definecolor{cR}{HTML}{00876C}  % reference / track data
\definecolor{cP}{HTML}{7A7A7A}  % fixed parameter
\definecolor{cD}{HTML}{8E44AD}  % derived quantity
\definecolor{cL}{HTML}{C0392B}  % limit / bound
\newcommand{\xS}[1]{{\color{cS}#1}}
\newcommand{\xU}[1]{{\color{cU}#1}}
\newcommand{\xR}[1]{{\color{cR}#1}}
\newcommand{\xP}[1]{{\color{cP}#1}}
\newcommand{\xD}[1]{{\color{cD}#1}}
\newcommand{\xL}[1]{{\color{cL}#1}}

% ---------------- layout ----------------------------------------------
\newenvironment{panel}[2]{%
  \par\addvspace{5pt}\noindent\begin{minipage}{\linewidth}%
  {\sffamily\small\bfseries\color{black!70}#1}\enspace{\sffamily\small #2}\par\vspace{4pt}}%
  {\end{minipage}\par\vspace{6pt}\noindent{\color{black!12}\rule{\linewidth}{0.5pt}}\par}
\newenvironment{figside}[1][0.53]{\begin{minipage}[c]{#1\linewidth}\centering}{\end{minipage}\hfill}
\newenvironment{eqside}[1][0.45]{\begin{minipage}[c]{#1\linewidth}\small
  \setlength{\abovedisplayskip}{3pt}\setlength{\belowdisplayskip}{3pt}%
  \setlength{\abovedisplayshortskip}{2pt}\setlength{\belowdisplayshortskip}{2pt}}{\end{minipage}}
\newcommand{\deps}[1]{\par\smallskip{\sffamily\scriptsize\bfseries\color{black!45}DEPENDS}\enspace{\small#1}}
\newcommand{\restr}[1]{\par\smallskip{\sffamily\scriptsize\bfseries\color{cL}RESTRAINS}\enspace{\small#1}}
\newcommand{\chip}[2]{\tikz[baseline=-0.6ex]\fill[#1] (0,0) circle (3pt);\,{\small\sffamily\color{#1}#2}\quad}
\providecommand{\sect}[1]{\par\vspace{8pt}{\large\sffamily\bfseries #1}\par\vspace{2pt}\noindent{\color{black!40}\rule{\linewidth}{0.8pt}}\par}

% ---------------- tikz styles -----------------------------------------
\tikzset{
  every picture/.style={line cap=round,line join=round,>={Stealth[length=4.5pt,width=3.2pt]}},
  ax/.style={->,black!45,line width=0.5pt},
  body/.style={draw=black!65,fill=black!5,line width=0.6pt,rounded corners=1.2pt},
  whl/.style={fill=black!75,draw=none,rounded corners=0.6pt,minimum width=5.5mm,minimum height=2mm,inner sep=0pt},
  vS/.style={->,cS,line width=0.9pt},
  vD/.style={->,cD,line width=0.9pt},
  vec/.style={->,line width=0.8pt},
  refl/.style={cR,line width=1pt},
  bnd/.style={black!55,line width=0.7pt},
  liml/.style={cL,line width=0.8pt},
  guide/.style={black!35,dash pattern=on 2pt off 1.5pt,line width=0.4pt},
  dim/.style={<->,black!50,line width=0.4pt,>={Stealth[length=3pt,width=2.4pt]}},
  lab/.style={font=\footnotesize,inner sep=1pt},
  sm/.style={font=\scriptsize,inner sep=1pt},
  dot/.style={circle,fill,inner sep=0pt,minimum size=3.4pt},
  odot/.style={circle,draw,fill=white,inner sep=0pt,minimum size=3.8pt,line width=0.6pt},
}
% small car icon: \caricon{(x,y)}{heading deg}{scale}
\newcommand{\caricon}[3]{%
  \begin{scope}[shift={#1},rotate=#2,scale=#3]
    \draw[body] (-0.45,-0.2) rectangle (0.45,0.2);
    \fill[black!55] (0.45,0) -- (0.28,0.11) -- (0.28,-0.11) -- cycle;
  \end{scope}}
% dependency-matrix row: {row}{label}{direct cols}{via cols}  (0 = none)
\newcommand{\mrow}[4]{%
  \draw[black!7,line width=3.6pt] (0.6,-#1) -- (\ncols+0.4,-#1);
  \node[anchor=east,font=\footnotesize,inner sep=1pt] at (0.45,-#1) {$#2$};
  \foreach \c in {#3}{\ifnum\c>0 \fill[black!85] (\c,-#1) circle (2.5pt);\fi}
  \foreach \c in {#4}{\ifnum\c>0 \draw[black!60,line width=0.6pt,fill=white] (\c,-#1) circle (2.2pt);\fi}}
\newcommand{\mcol}[3]{\node[font=\footnotesize,text=#3] at (#1,0.85) {$#2$};}

% =====================================================================

\usetikzlibrary{fit,positioning}

\usepackage{url}
\usepackage[hidelinks]{hyperref}

% ---------------- visible TODO notes -----------------------------------
\newif\iftodo \todotrue
\DeclareRobustCommand{\todo}[1]{\iftodo\ifmmode\text{\sffamily\small\color{cL}[TODO: #1]}%
  \else\leavevmode{\sffamily\small\color{cL}[TODO: #1]}\fi\fi}
\newcommand{\todoitems}[1]{\iftodo\par{\small\sffamily\color{cL}%
  \begin{itemize}[leftmargin=1.5em,itemsep=0pt,topsep=2pt,label=--]#1\end{itemize}}\fi}

% ---------------- notation ---------------------------------------------
\newcommand{\bx}{\mathbf x}
\newcommand{\bp}{\mathbf p}
\newcommand{\bv}{\mathbf v}
\newcommand{\bu}{\mathbf u}
\newcommand{\tauc}{\boldsymbol\tau}
\newcommand{\ok}{\checkmark}

\title{LOORacer: Learning-driven, Opponent-aware,\\ Optimal autonomous racing agent}
\author{Yiran Liu}
\date{Initial proposal \textendash\ draft on \today}

\begin{document}
\maketitle

% =====================================================================
\begin{abstract}
Trajectory optimization has been the dominant method in autonomous racing, but it lacks the ability to react in real time in head-to-head competition.
Recent learning-based approaches seek to address the gap, but they fall short in racing-time optimization.
We propose a learning-driven framework for autonomous racing policies that are both optimal and opponent-aware. 
We first learn a trajectory latent capturing trajectories that are generally feasible for a vehicle's physical constraints. 
Generative policies are pretrained using real-world professional driver demonstrations. 
Additionally, learned feasibility critics are parameterized from actual vehicle motion measurements and judge the policy on execution feasibility. They may also provide learning signals for subsequent reinforcement learning.

The goal is to develop policies that match top drivers in strategic thinking, produce time-optimal trajectories, and engage responsively in execution.
\end{abstract}

% =====================================================================
\section{Introduction}

\paragraph{Foundations.}
The kinematics of a racing vehicle are well-modelled by vehicle dynamics theories; exploiting optimal programming on such model allows us to produce best-time solutions to the racing problem. In [cite], an offline planner constructs the global trajectory, while an online controller replans over a short horizon to correct tracking error. In robot learning, a broad spectrum of generative policies [cite] predicts actions given observations, capturing the strategic decision and multimodel behavior of human experts. 

\paragraph{Gap.}
\todo{what current work does not cover, with Table~\ref{tab:gap};
RL racing agents \citep{wurman2022outracing}}

\begin{table}[h!]
\centering\small
\begin{tabular}{@{}lcccc@{}}
\toprule
 & Optimal & Learned prior & Opponent-aware & Real-time\\
\midrule
Offline trajectory optimization & \ok & -- & -- & --\\
Online contouring MPC           & horizon only & -- & partial & \ok\\
Imitation-learned driving policies & -- & \ok & \ok & \ok\\
Reinforcement-learning racing agents & \ok & -- & \ok & \ok\\
\midrule
This proposal (goal)            & \ok & \ok & \ok & \ok\\
\bottomrule
\end{tabular}
\caption{Properties of existing approaches (defined in Section~\ref{sec:props}).
\todo{verify each cell}}
\label{tab:gap}
\end{table}

% =====================================================================

\pagebreak
\section{Problem and scope}\label{sec:problem}

\subsection{Problem definition}

A vehicle races on a known race track, either against opponents or by itself. Let $\mathcal V$ denote the vehicle kinematics space, and $\mathcal P$ denote pose space, they combine to be state space $\mathcal X$. 
\[
\bp_t\in\mathcal P,\qquad \bv_t\in\mathcal V, \qquad \bx_t=[\,\bp_t\ \ \bv_t\,]\in\mathcal X
\]

In [cite], the vehicle functions in 2D track, and has kinematics state as: body velocity $[ v_x\ \ v_y ]$, yaw rate $r$, duty cycle $D$  
and steering $\delta_t$. See Appendix~\ref{app:maps} for details.
\[
\bp_t \in \mathrm{SE}(2)=[\,p_x\ \ p_y\ \ \varphi\,],\qquad
\bv_t=[\,v_x\ \ v_y \ \ r\ \ D\ \ \delta\,],
\]

The race track geometry is defined as a subset $\mathcal T$ of the pose space
$\mathcal P$:
\[
\mathcal T\subseteq\mathcal P,\qquad \bp_t\in\mathcal T\iff\text{"vehicle in track"}
\]

The vehicle can be modeled with fair accuracy by a known dynamics model $F_\theta$ with
vehicle parameters $\theta$ and control input $\bu_t=[\,\dot D_t\ \ \dot\delta_t\,]\in\mathcal U$:
\[
\bx_{t+\Delta t}=F_\theta(\bx_t,\bu_t,\Delta t).
\]

The physical constraints of the vehicle makes a subset of $\mathcal V$, denoting possible vehicle state:
\[
\mathcal V_{\mathrm{feas}}\subseteq\mathcal V,\qquad
\bp_t\in\mathcal T\ \wedge\ \bv_t\in\mathcal V_{\mathrm{feas}}\iff
\text{"vehicle can be in this state"}
\]

A trajectory is a sequence of $H$ future states sampled every $\Delta t$:
\[
\tauc_t=(\bx_{t+\Delta t},\,\bx_{t+2\Delta t},\,\dots,\,\bx_{t+H\Delta t}).
\]

Each opponent $j=1,\dots,M$ is observed by sensors to have estimation
$\boldsymbol\mu^j_t$ and uncertainty $\Sigma^j_t$. The policy generates a trajectory from both past the vehicle states, the track and the opponent observations:
\[
\tauc_t\sim\pi\big(\,\cdot\mid\bx_t,\ \mathcal T,\ \{(\boldsymbol\mu^j_t,\Sigma^j_t)\}_{j=1}^{M}\big),
\]

An online controller rolls out the trajectory by choosing $\bu_t$ to minimze tracking error.

\subsection{Defining properties}\label{sec:props}
\begin{description}[leftmargin=1.2em,itemsep=2pt]
  \item[Optimal.]        \todo{maximizes a racing objective, not merely feasible}
  \item[Learned prior.]  \todo{operates in a latent space learned from data}
  \item[Opponent-aware.] \todo{conditions on opponent observations with uncertainty; no collisions}
\end{description}
\todo{real-time requirement}

\subsection{Goals}\label{sec:goals}
\todo{goals, kept blurry: one policy that is optimal, operates in a learned prior, is
opponent-aware and runs in real time; transfer across vehicle types as a further goal}

\subsection{Assumptions}
\todoitems{
  \item opponent states observed as (position, uncertainty) $\Rightarrow$ perception not in scope
  \item vehicle capability = practical constraint model regressed from observed behaviour
  \item vehicle class: 2026 Formula~1 car (matches the data)
}

% =====================================================================
\section{Method}\label{sec:method}

\begin{figure}[h!]
\centering
\begin{tikzpicture}[
  box/.style={draw=black!45,rounded corners=2pt,fill=#1!8,align=center,
              font=\footnotesize,inner sep=4pt,minimum height=0.9cm,minimum width=2.0cm},
  box/.default=black,
  fl/.style={->,black!60,line width=0.6pt},
  lab/.style={font=\scriptsize,text=black!60,inner sep=1pt}]
  % data branch
  \node[box=cR] (tel)  at (0,0)      {F1 telemetry};
  \node[box=cR] (cur)  at (2.6,0)    {curation\\pipeline};
  \node[box=cR] (demo) at (5.3,0)    {curated\\demonstrations};
  \node[box=cP] (cm)   at (5.3,-2.35) {practical\\constraint model};
  % policy
  \node[box=cS,minimum width=2.4cm] (pol) at (9.0,0) {policy $\pi$};
  \node[box=cU] (tau)  at (12.6,0)   {trajectory $\tauc$};
  \node[box,font=\scriptsize,minimum height=0.6cm] (obs) at (9.0,1.5)
       {ego state, track, opponents $(\boldsymbol\mu,\Sigma)$};
  % critics
  \node[box=cL] (fc) at (12.6,-1.75) {feasibility critic};
  \node[box=cL] (ec) at (12.6,-2.95) {execution critic\\[-1pt]\scriptsize online MPC + simulation};
  \node[draw=black!35,dashed,rounded corners=3pt,inner sep=5pt,fit=(fc)(ec)] (crit) {};
  % arrows
  \draw[fl] (tel) -- (cur);
  \draw[fl] (cur) -- (demo);
  \draw[fl] (demo) -- node[lab,right]{regression} (cm);
  \draw[fl] (demo) -- node[lab,above]{imitation} (pol);
  \draw[fl] (obs) -- (pol);
  \draw[fl] (pol) -- (tau);
  \draw[fl] (tau) -- (crit.north -| tau);
  \draw[fl] (cm) -- node[lab,above]{parameterizes} (crit.west |- cm);
  \draw[fl,dashed] (crit.west |- fc) -| node[lab,pos=0.75,left,align=right]{evaluation /\\learning signal} (pol.south);
\end{tikzpicture}
\caption{Pipeline overview (placeholder). \todo{revise}}
\label{fig:pipeline}
\end{figure}

\subsection{Initialization from driver telemetry}
\todo{imitation learning from professional-driver telemetry of the 2026 season
\citep{tracinginsights2026,fastf1}}

\paragraph{Data accuracy.}
The telemetry is not first-hand vehicle measurement: (a) it does not contain the full state
$\bx_t$, so some components must be inferred; (b) it has a low sampling rate and moderate
noise. Data curation is therefore likely required and is part of what we develop.

\subsection{Two critics}\label{sec:critics}
\todo{both parameterized by the practical constraint model regressed from observed vehicle
behaviour; used for result measurement and future IL; classical constraint forms in
Appendix~\ref{app:classical}}
\begin{enumerate}[leftmargin=1.5em,itemsep=2pt]
  \item \textbf{Feasibility critic.} \todo{how feasible a trajectory is to execute}
  \item \textbf{Execution critic.}   \todo{simulation with online MPC; works well when
        (a) inference time on practical hardware fits, (b) continuity--reactiveness
        balance holds}
\end{enumerate}

\subsection{Learning methods}
\todo{set of methods adapted from SOTA robot policy learning, basic first; examples only,
no verdict: behaviour cloning, action chunking \citep{zhao2023act}, diffusion
\citep{chi2023diffusion}, flow matching \citep{lipman2022flow}, later A2A \citep{jia2026a2a}}

% =====================================================================
\section{Research questions}\label{sec:rq}
\todo{list the research questions; initial suggestions:}
\todoitems{
  \item Data: can telemetry with missing state components, low sampling rate and noise be
        curated into demonstrations usable for imitation learning?
  \item Constraint model: how well does a constraint model regressed from observed behaviour
        separate trajectories the vehicle executed from ones it did not (e.g.\ on held-out laps)?
  \item Methods: which learning methods adapted from robot learning produce trajectories
        that pass the feasibility critic?
  \item Execution: which of them can be executed by an online MPC in real time on practical
        hardware, and how do they trade continuity against reactiveness?
  \item Learned prior: does operating in a learned latent space change feasibility compared
        with generating trajectories directly?
  \item Opponents: how does conditioning on opponent observations and their uncertainty affect
        collisions and race outcome?
  \item Optimality: how close to optimal does imitation alone get, and what is needed beyond it?
  \item Transfer (further): what carries over to another vehicle type?
}

% =====================================================================
% TODO: remove this \nocite once every entry is cited in the text.
\nocite{heilmeier2019minimum,christ2019timeoptimal,liniger2015optimization,zhao2023act,
chi2023diffusion,lipman2022flow,black2025rtc,jia2026a2a,wurman2022outracing,
tracinginsights2026,fastf1}
\bibliographystyle{plainnat}
\bibliography{references}

% =====================================================================
\clearpage
\appendix
\newgeometry{margin=1.5cm}
\section{Classical racing optimizers}\label{app:classical}

\subsection{Constraint boundaries}\label{app:boundaries}
% =====================================================================
%  background_boundaries.tex -- one condensed figure: the constraint
%  boundaries of the three classical optimizers. Needs preamble_maps.tex,
%  tabularx, booktabs, array. Full equations: background_maps.tex.
% =====================================================================
\begin{figure}[h!]
\centering
\tikzset{
  mt/.style={black,line width=0.9pt},                          % TUM mintime
  mp/.style={black,line width=0.9pt,dash pattern=on 4pt off 2pt}, % MPCC
  mc/.style={black,line width=1.0pt,dotted},                   % TUM mincurv + speed profile
}
% legend
\begin{tikzpicture}
  \draw[mp] (0,0) -- (0.7,0); \node[sm,anchor=west] at (0.75,0) {MPCC (online)};
  \draw[mc] (3.3,0) -- (4.0,0); \node[sm,anchor=west] at (4.05,0) {\texttt{mincurv} + speed profile};
  \draw[mt] (8.0,0) -- (8.7,0); \node[sm,anchor=west] at (8.75,0) {\texttt{mintime}};
  \fill[cS!12] (11.0,-0.12) rectangle (11.6,0.12); \node[sm,anchor=west] at (11.65,0) {feasible};
\end{tikzpicture}\par\medskip
% (a) track
\begin{tikzpicture}
  \begin{scope}[shift={(0,-3.6)}]
    \fill[cS!12] (119:4.5) arc (119:61:4.5) -- (61:3.5) arc (61:119:3.5) -- cycle;
    \draw[bnd] (119:4.75) arc (119:61:4.75);
    \draw[bnd] (119:3.25) arc (119:61:3.25);
    \draw[cR,dash pattern=on 3pt off 2pt,line width=0.6pt] (119:4) arc (119:61:4);
    \draw[mp,cL] (-1.35,4.75) -- (1.35,4.75);
    \draw[mp,cL] (-1.35,3.25) -- (1.35,3.25);
    \node[dot,cR] at (90:4) {};
    \node[sm,text=cR,below=1pt] at (90:4) {$\bar s_k$};
    \draw[dim] (110:4.5) -- (110:4.75);
    \node[sm,text=black!60,anchor=east] at (110:4.85) {$\tfrac{w_{\mathrm{veh}}}{2}$};
  \end{scope}
  \node[sm] at (0,-0.95) {(a) track};
\end{tikzpicture}\hfill
% (b) grip
\begin{tikzpicture}[scale=0.95]
  \fill[cS!12] (0,0) circle (1);
  \draw[ax] (-1.45,0) -- (1.5,0) node[sm,right]{$a_y$};
  \draw[ax] (0,-1.3) -- (0,1.35) node[sm,above]{$a_x$};
  \draw[mt,cL] (0,0) circle (1);
  \draw[mp,cL] (0,0) ellipse (1.056 and 0.95);
  \draw[mc,cL] (1,0) -- (0,1) -- (-1,0) -- (0,-1) -- cycle;
  \node[sm] at (0,-1.65) {(b) grip};
\end{tikzpicture}\hfill
% (c) drive / brake
\begin{tikzpicture}[x=1.15cm,y=1.0cm]
  \fill[cS!12] plot[domain=0:3,samples=40] (\x,{min(1.0,0.8/max(\x,0.01))}) -- (3,-1) -- (0,-1) -- cycle;
  \draw[ax] (0,0) -- (3.3,0) node[sm,right]{$v$};
  \draw[ax] (0,-1.3) -- (0,1.3) node[sm,above]{$F_x$};
  \draw[mt,cL,domain=0:3,samples=60] plot (\x,{min(1.0,0.8/max(\x,0.01))});
  \draw[mp,cL,domain=0:3] plot (\x,{0.9-0.2*\x});
  \draw[mc,cL,domain=0:3,samples=60] plot (\x,{min(0.95,0.9/max(\x,0.01))-0.03*\x*\x});
  \draw[mt,cL] (0,-1.0) -- (3,-1.0);
  \draw[mp,cL] (0,-0.12) -- (3,-0.12);
  \draw[mc,cL,domain=0:3] plot (\x,{-0.9-0.03*\x*\x});
  \node[sm,text=cL,anchor=south west] at (1.45,0.72) {$P_{\max}/v$};
  \node[sm,text=black!60] at (2.6,-0.35) {brake};
  \node[sm] at (1.5,-1.65) {(c) drive / brake};
\end{tikzpicture}

\medskip
{\footnotesize\renewcommand{\arraystretch}{1.35}
\begin{tabularx}{\linewidth}{@{}>{\sffamily}l
  >{\raggedright\arraybackslash}X >{\raggedright\arraybackslash}X >{\raggedright\arraybackslash}X@{}}
\toprule
 & MPCC (online, 1.2\,s) & \texttt{mincurv} + speed profile & \texttt{mintime} (whole lap)\\
\midrule
track &
  $-\xL{R_{\mathrm{in}}}\le\bar{\mathbf n}^{\top}(\mathbf p_k-\mathbf p_{\mathrm{ref}}(\bar s_k))\le\xL{R_{\mathrm{out}}}$; slab at the guess, soft, car = point &
  $-\bar w_{\ell,i}\le\alpha_i\le\bar w_{r,i}$, \ $\bar w=w-\tfrac{w_{\mathrm{veh}}}{2}$ &
  $-w_r+\tfrac{w_{\mathrm{opt}}}{2}\le n\le w_\ell-\tfrac{w_{\mathrm{opt}}}{2}$\\
grip &
  rear: $(e_{\mathrm{long}}F_{r,x})^2+F_{r,y}^2\le(\xL{e_\varepsilon}D_r)^2$; front: $|\alpha_f|\le\xL{\alpha_{\max}}$; soft, linearized &
  $|\kappa_i|\le\xL{\kappa_{\max}}$; then $\big(\tfrac{|a_x|}{\mu a_{x,\max}(v)}\big)^{p}+\big(\tfrac{|a_y|}{\mu a_{y,\max}(v)}\big)^{p}\le1$, default $p{=}1$ &
  per tire: $F_{x,ij}^2+F_{y,ij}^2\le(\mu F_{z,ij})^2$, $F_z$ with aero and load transfer\\
drive / brake &
  $-0.1\le D\le1$, \ $F_{r,x}=(C_{m1}-C_{m2}v_x)D$; no front brake &
  $a_x\le\xL{a_{x,\mathrm{mach}}(v)}-\tfrac{c_d}{m}v^2$ &
  $vF_d\le\xL{P_{\max}}$, \ $F_d\ge0\ge F_b\ge-\xL{F_{b,\max}}$, \ $F_dF_b\approx0$\\
rates &
  box bounds on $\dot D,\dot\delta,\dot v_s$ &
  -- &
  $\tfrac{\Delta\delta}{\Delta t}\in\pm\tfrac{\delta_{\max}}{t_\delta}$, similar for $F_d,F_b$\\
structure &
  $|s_k-\bar s_k|\le0.2$\,m; dynamics linearized per QP &
  no dynamics: geometry, then a speed sweep &
  $\mathbf x(0)=\mathbf x(L)$; dynamics at collocation points\\
\bottomrule
\end{tabularx}}
\caption{Constraint boundaries of the classical optimizers, condensed. Shaded: feasible
region. Shapes are schematic; sizes differ per optimizer. Full equations, variables and
dependencies: Appendix~\ref{app:maps}.}
\label{fig:boundaries}
\end{figure}


\clearpage
\subsection{Equation maps}\label{app:maps}
% =====================================================================
%  background_maps.tex -- equation maps (body only).
%  Online: Liniger MPCC (C++ master). Offline: TUM global_racetrajectory_optimization.
%  \input inside a document whose preamble % =====================================================================
%  preamble_maps.tex -- packages, colours, macros and TikZ styles for the
%  equation maps. \input from the preamble of any document that
%  % =====================================================================
%  background_maps.tex -- equation maps (body only).
%  Online: Liniger MPCC (C++ master). Offline: TUM global_racetrajectory_optimization.
%  \input inside a document whose preamble \input{preamble_maps}.
%  \sect{} renders the two part headings; redefine it before \input{preamble_maps}
%  to change how they look (e.g. as \subsection).
% =====================================================================
{\setlength{\parindent}{0pt}%
\chip{cS}{decision state}\chip{cU}{control input}\chip{cR}{reference / track data}%
\chip{cP}{fixed parameter}\chip{cD}{derived}\chip{cL}{limit / bound}

% =====================================================================
\sect{Online \quad MPCC \quad (receding horizon, 60 steps $\times$ 20 ms)}

% ---------------------------------------------------------------- A0
\begin{panel}{A0}{Variables and who depends on whom}
\begin{figside}[0.60]
\begin{tikzpicture}[x=0.5cm,y=0.37cm]
  \def\ncols{14}
  \mcol{1}{X}{cS}\mcol{2}{Y}{cS}\mcol{3}{\varphi}{cS}\mcol{4}{v_x}{cS}\mcol{5}{v_y}{cS}
  \mcol{6}{r}{cS}\mcol{7}{s}{cS}\mcol{8}{D}{cS}\mcol{9}{\delta}{cS}\mcol{10}{v_s}{cS}
  \mcol{11}{\dot D}{cU}\mcol{12}{\dot\delta}{cU}\mcol{13}{\dot v_s}{cU}\mcol{14}{\mathbf p_{\!\mathrm{ref}}}{cR}
  \mrow{1}{\xD{\alpha_f}}{4,5,6,9}{0}
  \mrow{2}{\xD{\alpha_r}}{4,5,6}{0}
  \mrow{3}{\xD{F_{f,y}}}{0}{4,5,6,9}
  \mrow{4}{\xD{F_{r,y}}}{0}{4,5,6}
  \mrow{5}{\xD{F_{r,x}}}{4,8}{0}
  \mrow{6}{\xD{F_{\mathrm{fric}}}}{4}{0}
  \mrow{7.4}{\xS{\dot X,\dot Y}}{3,4,5}{0}
  \mrow{8.4}{\xS{\dot\varphi}}{6}{0}
  \mrow{9.4}{\xS{\dot v_x}}{5,6,9}{4,8}
  \mrow{10.4}{\xS{\dot v_y}}{4,6,9}{5}
  \mrow{11.4}{\xS{\dot r}}{9}{4,5,6}
  \mrow{12.4}{\xS{\dot s}}{10}{0}
  \mrow{13.4}{\xS{\dot D,\dot\delta,\dot v_s}}{11,12,13}{0}
  \mrow{14.8}{\xD{e_c,\,e_l}}{1,2,7,14}{0}
  \mrow{15.8}{\xL{\text{track}}}{1,2,14}{0}
  \mrow{16.8}{\xL{\text{ellipse}}}{0}{4,5,6,8}
  \mrow{17.8}{\xL{|\alpha_f|}}{0}{4,5,6,9}
  \mrow{18.8}{\beta\ \text{cost}}{4,5}{0}
  \mrow{19.8}{\varphi\ \text{cost}}{3,7,14}{0}
  \mrow{20.8}{\text{progress}}{10}{0}
  \draw[black!30] (0.5,-6.7) -- (14.5,-6.7);
  \draw[black!30] (0.5,-14.1) -- (14.5,-14.1);
  \draw[black!30] (10.5,1.4) -- (10.5,-21.3);
\end{tikzpicture}
\end{figside}
\begin{eqside}[0.38]
\[\xS{\mathbf x} = [\,\xS{X\ Y\ \varphi\ v_x\ v_y\ r\ s\ D\ \delta\ v_s}\,]^{\top}\]
\[\xU{\mathbf u} = [\,\xU{\dot D\ \ \dot\delta\ \ \dot v_s}\,]^{\top}\]
\[\xS{\dot s}=\xS{v_s},\quad \xS{\dot D}=\xU{u_1},\quad \xS{\dot\delta}=\xU{u_2},\quad \xS{\dot v_s}=\xU{u_3}\]
\medskip
{\footnotesize
\tikz\fill[black!85] (0,0) circle (2.5pt);\ appears in the formula\par
\tikz\draw[black!60,line width=0.6pt] (0,0) circle (2.2pt);\ enters through $\alpha$ or $F$\par}
\restr{inputs are \emph{rates}, so rate costs and slew bounds act directly on $\xU{\mathbf u}$.}
\end{eqside}

\par\medskip
{\footnotesize\renewcommand{\arraystretch}{1.12}%
\begin{minipage}[t]{0.335\linewidth}
{\sffamily\scriptsize\bfseries\color{cS}STATE $\mathbf x$}\par\smallskip
\begin{tabular}{@{}l@{\ \ }>{\raggedright\arraybackslash}p{4.35cm}@{}}
$\xS{X,Y}$       & car position, world frame [m]\\
$\xS{\varphi}$   & heading (yaw angle), world frame [rad]\\
$\xS{v_x}$       & forward velocity, body frame [m/s]\\
$\xS{v_y}$       & sideways velocity, body frame [m/s]\\
$\xS{r}$         & yaw rate [rad/s]\\
$\xS{s}$         & progress: virtual arc length along the reference [m]\\
$\xS{D}$         & motor duty cycle, $[-0.1,1]$; $<0$ brakes\\
$\xS{\delta}$    & front steering angle, $|\delta|\le 0.35$ rad\\
$\xS{v_s}$       & progress speed $\dot s\ge 0$ [m/s]
\end{tabular}
\end{minipage}\hfill
\begin{minipage}[t]{0.32\linewidth}
{\sffamily\scriptsize\bfseries\color{cU}INPUT $\mathbf u$}\par\smallskip
\begin{tabular}{@{}l@{\ \ }>{\raggedright\arraybackslash}p{4.05cm}@{}}
$\xU{\dot D}$         & duty-cycle rate\\
$\xU{\dot\delta}$     & steering rate\\
$\xU{\dot v_s}$       & progress acceleration
\end{tabular}\par\medskip
{\sffamily\scriptsize\bfseries\color{cR}REFERENCE}\par\smallskip
\begin{tabular}{@{}l@{\ \ }>{\raggedright\arraybackslash}p{4.05cm}@{}}
$\xR{\mathbf p_{\mathrm{ref}}}$ & centreline spline $(X_{\mathrm{ref}}(s),Y_{\mathrm{ref}}(s))$ with heading $\Phi(s)$; fixed data
\end{tabular}\par\medskip
{\sffamily\scriptsize\bfseries\color{cD}DERIVED}\par\smallskip
\begin{tabular}{@{}l@{\ \ }>{\raggedright\arraybackslash}p{4.05cm}@{}}
$\xD{\alpha_f,\alpha_r}$       & front / rear tire slip angle\\
$\xD{F_{f,y},F_{r,y}}$         & front / rear lateral tire force\\
$\xD{F_{r,x}}$                 & rear drive force from the motor\\
$\xD{F_{\mathrm{fric}}}$       & rolling resistance + air drag
\end{tabular}
\end{minipage}\hfill
\begin{minipage}[t]{0.31\linewidth}
{\sffamily\scriptsize\bfseries\color{black!60}COST / CONSTRAINT}\par\smallskip
\begin{tabular}{@{}l@{\ \ }>{\raggedright\arraybackslash}p{3.55cm}@{}}
$\dot{(\,\cdot\,)}$        & time derivative: the state equations\\
$\xD{e_c}$                 & contouring error: sideways offset from $\mathbf p_{\mathrm{ref}}(s)$\\
$\xD{e_l}$                 & lag error: along-track gap between $s$ and the car\\
\xL{track}                 & car inside the slab around the reference\\
\xL{ellipse}               & rear friction ellipse\\
$\xL{|\alpha_f|}$          & front slip-angle limit\\
$\beta$ cost               & sideslip $\arctan(v_y/v_x)$ penalty\\
$\varphi$ cost             & heading vs.\ track heading $\Phi(s)$\\
progress                   & reward $-q_{v_s}v_s$
\end{tabular}
\end{minipage}}
\end{panel}

% ---------------------------------------------------------------- A1
\begin{panel}{A1}{Body velocity $\to$ world pose}
\begin{figside}
\begin{tikzpicture}
  \draw[ax] (0,0) -- (1.0,0) node[sm,right]{$X$};
  \draw[ax] (0,0) -- (0,1.0) node[sm,above]{$Y$};
  \coordinate (C) at (2.9,0.45);
  % world components of v (dotted)
  \draw[black!55,dotted,line width=0.9pt] (C) -- ++(1.282,0) node[sm,below,midway,yshift=-2pt]{$\dot X$}
        -- ++(0,1.779) node[sm,right,pos=0.8]{$\dot Y$};
  % heading angle, measured behind the car
  \draw[guide] (C) -- ++(180:1.5);
  \draw[guide] (C) -- ++(210:1.5);
  \draw[cS,line width=0.6pt] ($(C)+(180:1.15)$) arc (180:210:1.15);
  \node[lab,text=cS] at ($(C)+(195:1.42)$) {$\varphi$};
  \begin{scope}[shift={(C)},rotate=30]
    \draw[body] (-0.55,-0.22) rectangle (0.55,0.22);
    \fill[black!55] (0.55,0) -- (0.4,0.1) -- (0.4,-0.1) -- cycle;
    \draw[vS] (0,0) -- (2.0,0) node[lab,below right=0pt]{$v_x$};
    \draw[vS] (0,0) -- (0,0.9) node[lab,above left=0pt]{$v_y$};
    \draw[->,black!75,line width=0.6pt] (0,0) -- (2.0,0.9) node[lab,above=1pt]{$\mathbf v$};
  \end{scope}
  \draw[cS,->,line width=0.7pt] ($(C)+(245:0.7)$) arc (245:305:0.7);
  \node[lab,text=cS] at ($(C)+(275:0.92)$) {$r$};
\end{tikzpicture}
\end{figside}
\begin{eqside}
\[\begin{aligned}
\xS{\dot X} &= \xS{v_x}\cos\xS{\varphi} - \xS{v_y}\sin\xS{\varphi}\\
\xS{\dot Y} &= \xS{v_x}\sin\xS{\varphi} + \xS{v_y}\cos\xS{\varphi}\\
\xS{\dot\varphi} &= \xS{r}
\end{aligned}\]
\deps{$(\dot X,\dot Y)\leftarrow \xS{v_x,v_y,\varphi}$;\quad $\dot\varphi\leftarrow\xS{r}$}
\restr{pose = integral of the body velocity rotated by the heading.}
\end{eqside}
\end{panel}

% ---------------------------------------------------------------- A2
\begin{panel}{A2}{Slip angles}
\begin{figside}
\begin{tikzpicture}
  \coordinate (R) at (0,0); \coordinate (G) at (1.5,0); \coordinate (F) at (3.2,0);
  \draw[black!40,line width=1.4pt] (R) -- (F);
  \draw[guide] (F) -- ++(0:1.7);
  \draw[guide] ($(F)+(208:0.9)$) -- ++(28:2.7);
  \node[whl] at (R) {};
  \node[whl,rotate=28] at (F) {};
  \fill (G) circle (1.6pt);
  % rear
  \draw[vS] (R) -- ++(-16:1.55) node[sm,below=2pt]{$(v_x,\ v_y{-}l_r r)$};
  \draw[cD,line width=0.8pt] ($(R)+(-16:1.1)$) arc (-16:0:1.1);
  \node[lab,text=cD] at ($(R)+(-6:1.35)$) {$\alpha_r$};
  % front
  \draw[vS] (F) -- ++(10:1.6) node[sm,right]{$(v_x,\ v_y{+}l_f r)$};
  \draw[cD,line width=0.8pt] ($(F)+(10:1.25)$) arc (10:28:1.25);
  \node[lab,text=cD] at ($(F)+(19:1.5)$) {$\alpha_f$};
  \draw[cS,line width=0.6pt] ($(F)+(180:0.75)$) arc (180:208:0.75);
  \node[lab,text=cS] at ($(F)+(194:1.0)$) {$\delta$};
  % dims
  \draw[dim] (0,-1.05) -- (1.5,-1.05) node[sm,midway,below]{$l_r$};
  \draw[dim] (1.5,-1.05) -- (3.2,-1.05) node[sm,midway,below]{$l_f$};
\end{tikzpicture}
\end{figside}
\begin{eqside}
\[\begin{aligned}
\xD{\alpha_f} &= \xS{\delta} - \arctan\frac{\xS{v_y}+\xP{l_f}\,\xS{r}}{\xS{v_x}}\\[2pt]
\xD{\alpha_r} &= -\arctan\frac{\xS{v_y}-\xP{l_r}\,\xS{r}}{\xS{v_x}}
\end{aligned}\]
\deps{$\alpha_f\leftarrow\xS{v_x,v_y,r,\delta}$;\quad $\alpha_r\leftarrow\xS{v_x,v_y,r}$}
\restr{nothing alone: angle between where a tire points and where it moves (bounded in A8).}
\end{eqside}
\end{panel}

% ---------------------------------------------------------------- A3
\begin{panel}{A3}{Tire and drivetrain forces}
\begin{figside}[0.56]
\begin{tikzpicture}
  \begin{scope}[xscale=1.35,yscale=7]
    \fill[cL!10] (0.6,-0.215) rectangle (1.65,0.215);
    \fill[cL!10] (-1.65,-0.215) rectangle (-0.6,0.215);
    \draw[ax] (-1.7,0) -- (1.8,0) node[sm,right]{$\alpha$};
    \draw[ax] (0,-0.22) -- (0,0.25) node[sm,above]{$F_y$};
    \draw[guide] (-1.65,0.192) -- (1.65,0.192);
    \node[sm,text=cP,anchor=south east] at (-0.05,0.193) {$D_f$};
    \draw[black!50,line width=0.5pt] (-0.15,-0.0891) -- (0.26,0.1545);
    \node[sm,text=cP,anchor=east] at (-0.05,0.12) {$B_fC_fD_f$};
    \draw[cD,line width=1pt,domain=-1.6:1.6,samples=100,smooth] plot (\x,{0.192*sin(1.2*atan(2.579*\x))});
    \draw[liml] (0.6,-0.215) -- (0.6,0.215);
    \draw[liml] (-0.6,-0.215) -- (-0.6,0.215);
    \node[sm,text=cL,above] at (0.6,0.215) {$\alpha_{\max}$};
  \end{scope}
  \begin{scope}[shift={(3.35,-0.65)},xscale=0.6,yscale=7]
    \draw[ax] (0,0) -- (3.9,0) node[sm,right]{$v_x$};
    \draw[ax] (0,-0.09) -- (0,0.34) node[sm,above]{$F_x$};
    \draw[cD,line width=1pt,domain=0:3.5] plot (\x,{0.287-0.0545*\x}) node[sm,right]{$D{=}1$};
    \draw[cD!50,line width=1pt,domain=0:3.5] plot (\x,{0.5*(0.287-0.0545*\x)}) node[sm,right]{$D{=}\tfrac12$};
    \draw[cD,dashed,line width=0.8pt,domain=0:3.5] plot (\x,{-0.0518-0.00035*\x*\x}) node[sm,right]{$F_{\mathrm{fric}}$};
  \end{scope}
\end{tikzpicture}
\end{figside}
\begin{eqside}[0.42]
\[\begin{aligned}
\xD{F_{f,y}} &= \xP{D_f}\sin\!\big(\xP{C_f}\arctan(\xP{B_f}\,\xD{\alpha_f})\big)\\
\xD{F_{r,y}} &= \xP{D_r}\sin\!\big(\xP{C_r}\arctan(\xP{B_r}\,\xD{\alpha_r})\big)\\
\xD{F_{r,x}} &= (\xP{C_{m1}}-\xP{C_{m2}}\,\xS{v_x})\,\xS{D}\\
\xD{F_{\mathrm{fric}}} &= -\xP{C_{r0}}-\xP{C_{r2}}\,\xS{v_x}^2
\end{aligned}\]
{\footnotesize $\xP{D_f},\xP{D_r}$: peak force $\approx\mu F_z$ \quad $\xS D$: duty cycle $\in[-0.1,1]$}
\deps{$F_y\leftarrow\xD{\alpha}$;\quad $F_{r,x}\leftarrow\xS{D,v_x}$;\quad $F_{\mathrm{fric}}\leftarrow\xS{v_x}$}
\restr{lateral force saturates at $D$; drive force fades with speed; front tire has no $F_x$.}
\end{eqside}
\end{panel}

% ---------------------------------------------------------------- A4
\begin{panel}{A4}{Newton--Euler in the body frame}
\begin{figside}
\begin{tikzpicture}
  \coordinate (R) at (0,0); \coordinate (G) at (1.6,0); \coordinate (F) at (3.4,0);
  \draw[black!40,line width=1.4pt] (R) -- (F);
  \node[whl] at (R) {};
  \node[whl,rotate=22] at (F) {};
  \fill (G) circle (1.6pt);
  \draw[vD] (R) -- ++(0,1.0) node[lab,above]{$F_{r,y}$};
  \draw[vD] (-1.1,0) -- (-0.33,0);
  \node[lab,text=cD,below] at (-0.75,-0.05) {$F_{r,x}$};
  \draw[vD] (F) -- ++(112:1.05) node[lab,above]{$F_{f,y}$};
  \draw[vD,dashed] (G) -- ++(-0.75,0);
  \node[lab,text=cD,above] at (1.15,0.05) {$F_{\mathrm{fric}}$};
  \draw[cS,->,line width=0.7pt] ($(G)+(-150:0.5)$) arc (-150:-30:0.5);
  \node[lab,text=cS] at ($(G)+(-90:0.72)$) {$r$};
  \draw[dim] (0,-1.15) -- (1.6,-1.15) node[sm,midway,below]{$l_r$};
  \draw[dim] (1.6,-1.15) -- (3.4,-1.15) node[sm,midway,below]{$l_f$};
\end{tikzpicture}
\end{figside}
\begin{eqside}
\[\begin{aligned}
\xP m\,\xS{\dot v_x} &= \xD{F_{r,x}}+\xD{F_{\mathrm{fric}}}-\xD{F_{f,y}}\sin\xS\delta+\xP m\,\xS{v_y}\,\xS r\\
\xP m\,\xS{\dot v_y} &= \xD{F_{r,y}}+\xD{F_{f,y}}\cos\xS\delta-\xP m\,\xS{v_x}\,\xS r\\
\xP{I_z}\,\xS{\dot r} &= \xD{F_{f,y}}\,\xP{l_f}\cos\xS\delta-\xD{F_{r,y}}\,\xP{l_r}
\end{aligned}\]
{\footnotesize $\pm m\,v\,r$ terms: the body frame itself rotates.}
\deps{$(\dot v_x,\dot v_y,\dot r)\leftarrow\xD{F}$'s, $\xS{\delta}$, $\xS{v_x,v_y,r}$}
\restr{accelerations = forces / inertia. Linearized + ZOH-discretized each QP (A9).}
\end{eqside}
\end{panel}

% ---------------------------------------------------------------- A5
\begin{panel}{A5}{Progress $s$, contouring error $e_c$, lag error $e_l$}
\begin{figside}
\begin{tikzpicture}
  \begin{scope}[shift={(0,-4.4)}]
    \draw[refl] (127:5) arc (127:58:5) node[sm,right]{$\mathbf p_{\mathrm{ref}}(\cdot)$};
    \coordinate (P) at (102:5);
    \coordinate (S) at (88:5);
    \coordinate (p) at (88:4.15);
    \coordinate (N) at ($(P)+(-78:1)$);
    \coordinate (Q) at ($(P)!(p)!(N)$);
    \draw[guide] ($(P)+(192:0.6)$) -- ($(P)+(12:2.9)$);
    \draw[vec,black!70] (P) -- ++(12:0.65) node[sm,above]{$\mathbf t$};
    \draw[black!45,dotted,line width=0.8pt] (p) -- (S);
    \draw[cD,line width=1.2pt] (P) -- (Q) node[lab,midway,left=1pt]{$e_c$};
    \draw[cD,line width=1.2pt] (Q) -- (p) node[lab,midway,below=1pt]{$e_l$};
    \caricon{(p)}{-2}{0.65}
    \node[dot,cR] at (P) {};
    \node[sm,text=cR,above left] at (P) {$s$};
    \node[odot,draw=cR] at (S) {};
    \node[sm,text=cR,above] at (S) {$s^{\ast}$};
    \node[sm,below right=4pt] at (p) {$(X,Y)$};
  \end{scope}
\end{tikzpicture}
\end{figside}
\begin{eqside}
\[\begin{aligned}
\xD{e_c} &= \phantom{-}\xR{\mathbf n(\xS s)}^{\!\top}\big(\xS{\mathbf p}-\xR{\mathbf p_{\mathrm{ref}}(\xS s)}\big)\\
\xD{e_l} &= -\xR{\mathbf t(\xS s)}^{\!\top}\big(\xS{\mathbf p}-\xR{\mathbf p_{\mathrm{ref}}(\xS s)}\big)
\end{aligned}\]
{\footnotesize $\xR{\mathbf t}=(\cos\Phi,\sin\Phi)$,\ \ $\xR{\mathbf n}=(\sin\Phi,-\cos\Phi)$ (right),\ \ $\xS{\mathbf p}=(\xS X,\xS Y)$}
\deps{$e_c,e_l\leftarrow\xS{X,Y,s}$ and the spline $\xR{\mathbf p_{\mathrm{ref}}}$}
\restr{nothing alone; $e_l\!\approx\!0 \Leftrightarrow s\!\approx\!s^{\ast}$ (projection), with no inner $\arg\min$.}
\end{eqside}
\end{panel}

% ---------------------------------------------------------------- A6
\begin{panel}{A6}{Stage cost as springs}
\begin{figside}
\begin{tikzpicture}
  \draw[refl] (-0.2,0) -- (5.6,0);
  \coordinate (P) at (3.5,0); \coordinate (Fp) at (1.9,0); \coordinate (p) at (1.9,1.15);
  \draw[cD,line width=1.1pt,decorate,decoration={zigzag,segment length=3.2pt,amplitude=2.6pt}] (Fp) -- (P);
  \draw[cD!75,line width=0.5pt,decorate,decoration={coil,aspect=0.45,segment length=6pt,amplitude=3pt}] (Fp) -- (p);
  \caricon{(p)}{0}{0.8}
  \node[dot,cR] at (P) {};
  \node[sm,text=cR,below=2pt] at (P) {$s$};
  \draw[vS,line width=1.2pt] (P) -- ++(1.15,0);
  \node[lab,text=cS,above] at ($(P)+(0.62,0.05)$) {$q_{v_s}v_s$};
  \node[lab,text=cD,below] at ($(Fp)!0.5!(P)+(0,-0.12)$) {$q_l\,e_l^2$\ \ \textcolor{cP}{\scriptsize$q_l{=}500$}};
  \node[lab,text=cD,left=5pt] at ($(Fp)!0.5!(p)+(0,0.12)$) {$q_c\,e_c^2$};
  \node[sm,text=cP,left=5pt] at ($(Fp)!0.5!(p)+(0,-0.22)$) {$q_c{=}0.1$};
\end{tikzpicture}
\end{figside}
\begin{eqside}
\[J=\sum_{k=0}^{N}\Big(\xP{q_c}\xD{e_c}^2+\xP{q_l}\xD{e_l}^2-\xP{q_{v_s}}\xS{v_s}\Big)+\text{reg.}\]
{\footnotesize reg.\ $= \xP{q_\beta}\beta^2+\xP{q_r}\xS r^2+\xP{q_\mu}(\xS\varphi-\xR\Phi)^2+\|(\xS{D,\delta,v_s})\|^2_{\xP{R}}+\|\xU{\mathbf u}\|^2_{\xP{R_\Delta}}$\\
\phantom{reg.\ }$\beta=\arctan(\xS{v_y}/\xS{v_x})$;\ \ $q_c,q_r\times10$ at $k=N$;\ \ $q_{v_s}{=}0.02$}
\deps{$e_c,e_l$ (A5), $\xS{v_s}$, $\xS{v_x,v_y,r}$, $\xU{\mathbf u}$}
\restr{soft only: $s$ pinned to the car (stiff), lateral position nearly free (weak), progress rewarded.}
\end{eqside}
\end{panel}

% ---------------------------------------------------------------- A7
\begin{panel}{A7}{Track constraint (linearized slab)}
\begin{figside}
\begin{tikzpicture}
  \begin{scope}[shift={(0,-4.3)}]
    \coordinate (P) at (95:4.5);
    \coordinate (Ao) at (95:5.25);
    \coordinate (Ai) at (95:3.75);
    \path ($(Ao)+(185:2.6)$) coordinate (o1) ($(Ao)+(5:2.6)$) coordinate (o2)
          ($(Ai)+(185:2.6)$) coordinate (i1) ($(Ai)+(5:2.6)$) coordinate (i2);
    \fill[cS!7] (o1) -- (o2) -- (i2) -- (i1) -- cycle;
    \begin{scope}
      \clip (o1) -- (o2) -- (i2) -- (i1) -- cycle;
      \fill[cL!20,even odd rule] (0,0) circle (7) (0,0) circle (5.25);
    \end{scope}
    \draw[bnd] (128:5.25) arc (128:60:5.25);
    \draw[bnd] (128:3.75) arc (128:60:3.75);
    \draw[cR,dash pattern=on 4pt off 2pt,line width=0.8pt] (128:4.5) arc (128:60:4.5);
    \draw[liml] (o1) -- (o2);
    \draw[liml] (i1) -- (i2);
    \draw[vec,black!70] (P) -- ++(95:0.5) node[sm,right]{$\bar{\mathbf n}$};
    \node[dot,cR] at (P) {};
    \node[sm,text=cR,below right] at (P) {$\bar s_k$};
    \draw[dim] ($(P)+(185:1.7)$) -- ($(Ao)+(185:1.7)$) node[sm,midway,left]{$R_{\mathrm{out}}$};
    \draw[dim] ($(P)+(185:1.7)$) -- ($(Ai)+(185:1.7)$) node[sm,midway,left]{$R_{\mathrm{in}}$};
    \node[dot,cS] at ($(P)+(5:1.3)+(95:-0.42)$) {};
    \node[sm,text=cS,below] at ($(P)+(5:1.3)+(95:-0.42)$) {$\mathbf p_k$};
    \node[sm,text=cL] at ($(o2)+(-0.15,-0.35)$) {\tiny lin.\ error};
  \end{scope}
\end{tikzpicture}
\end{figside}
\begin{eqside}
\[-\xL{R_{\mathrm{in}}}\ \le\ \xR{\bar{\mathbf n}}^{\!\top}\big(\xS{\mathbf p_k}-\xR{\mathbf p_{\mathrm{ref}}(\bar s_k)}\big)\ \le\ \xL{R_{\mathrm{out}}}\quad(+\,\xi_k\ge 0)\]
{\footnotesize $\bar s_k$: progress of the \emph{previous guess} (fixed in this QP); $\xR{\bar{\mathbf n}}$: left normal there;\\
$R_{\mathrm{in}}{=}R_{\mathrm{out}}{=}0.14$\,m constant; car treated as a point.}
\deps{$\xS{X_k,Y_k}$ only}
\restr{car point inside a straight slab; soft (slack $\xi$, penalized); can overshoot the real edge far from $\bar s_k$.}
\end{eqside}
\end{panel}

% ---------------------------------------------------------------- A8
\begin{panel}{A8}{Tire constraints}
\begin{figside}
\begin{tikzpicture}[scale=0.72]
  \fill[cS!7] (0,0) ellipse (1.83 and 1.65);
  \draw[ax] (-2.3,0) -- (2.9,0) node[sm,right]{$F_{r,x}$};
  \draw[ax] (0,-2.0) -- (0,2.3) node[sm,above]{$F_{r,y}$};
  \draw[guide] (0,0) circle (1.737);
  \draw[liml] (0,0) ellipse (1.83 and 1.65);
  \draw[vD] (0,0) -- (0.9,1.2);
  \node[dot,cD] at (0.9,1.2) {};
  \node[sm,text=cD,right=2pt] at (0.9,1.2) {$(F_{r,x},F_{r,y})$};
  \node[sm,text=cL,anchor=north west] at (1.86,-0.3) {$e_\varepsilon D_r/e_{\mathrm{long}}$};
  \node[sm,text=cL,anchor=south east] at (-0.05,1.68) {$e_\varepsilon D_r$};
  \node[sm,text=black!55,anchor=north east] at (-1.25,-1.25) {$D_r$};
\end{tikzpicture}
\end{figside}
\begin{eqside}
\[\Big(\tfrac{\xP{e_{\mathrm{long}}}\,\xD{F_{r,x}}}{\xP{F_{N,r}}}\Big)^2+\Big(\tfrac{\xD{F_{r,y}}}{\xP{F_{N,r}}}\Big)^2\ \le\ \Big(\tfrac{\xL{e_\varepsilon}\,\xP{D_r}}{\xP{F_{N,r}}}\Big)^2\]
\[|\xD{\alpha_f}|\ \le\ \xL{\alpha_{\max}}\]
{\footnotesize $e_{\mathrm{long}}{=}0.9$,\ $e_\varepsilon{=}0.95$,\ $\alpha_{\max}{=}0.6$; $F_{N,r}$ static load. Both first-order in $\xS{\mathbf x}$ around the guess.}
\deps{ellipse $\leftarrow\xS{v_x,v_y,r,D}$;\quad $|\alpha_f|\leftarrow\xS{v_x,v_y,r,\delta}$}
\restr{rear: drive + lateral force share one ellipse at 95\,\% of peak. Front: slip below $\alpha_{\max}$ (A3, red band). Soft.}
\end{eqside}
\end{panel}

% ---------------------------------------------------------------- A9
\begin{panel}{A9}{One QP per SQP step, trust region on $s$}
\begin{figside}
\begin{tikzpicture}[x=0.52cm,y=1cm]
  \fill[cL!12] (0,0.0) -- (11,2.18) -- (11,2.68) -- (0,0.5) -- cycle;
  \draw[ax] (-0.3,0) -- (12.4,0) node[sm,right]{$k$};
  \draw[ax] (0,-0.1) -- (0,2.8) node[sm,left]{$s_k$};
  \draw[black!40,line width=0.5pt] (0,0.25) -- (11,2.43);
  \foreach \k in {1,...,7,11} {
    \pgfmathsetmacro\sb{0.25+0.198*\k}
    \pgfmathsetmacro\so{\sb+0.17*sin(40*\k)}
    \node[odot,draw=black!45] at (\k,\sb) {};
    \node[dot,cS] at (\k,\so) {};
  }
  \node[dot,cS] at (0,0.25) {};
  \node[sm,left] at (0,0.25) {$\hat{\mathbf x}$};
  \node[sm] at (9,1.4) {$\cdots$};
  \foreach \k/\l in {0/0,1/1,2/2,11/N} \node[sm,below] at (\k,-0.05) {$\l$};
  \draw[decorate,decoration={brace,amplitude=4pt},black!50] (0.2,3.05) -- (11,3.05) node[sm,midway,above=4pt]{$N\,T_s=60\times20\,\text{ms}$};
  \node[sm,text=cL] at (6.3,2.55) {$|s_k-\bar s_k|\le 0.2$\,m};
  \node[sm,text=black!50] at (8.6,0.95) {$\bar s_k$ (guess)};
\end{tikzpicture}
\end{figside}
\begin{eqside}
\[\begin{aligned}
\min\ &\textstyle\sum_k \tfrac12\xS{\mathbf x_k}^{\!\top}\!Q_k\xS{\mathbf x_k}+\mathbf q_k^{\!\top}\xS{\mathbf x_k}+\tfrac12\xU{\mathbf u_k}^{\!\top}\!R\,\xU{\mathbf u_k}+\text{slack}\\
\text{s.t. }&\xS{\mathbf x_0}=\hat{\mathbf x},\quad \xS{\mathbf x_{k+1}}=A_k\xS{\mathbf x_k}+B_k\xU{\mathbf u_k}+\mathbf g_k\\
&\text{A7, A8 (soft)},\ \ \text{box bounds on }\xS{\mathbf x},\xU{\mathbf u}\\
&|\xS{s_k}-\bar s_k|\le 0.2\,\text{m}
\end{aligned}\]
{\footnotesize $(A_k,B_k,\mathbf g_k)$: exact ZOH of the Jacobian at $(\bar{\mathbf x}_k,\bar{\mathbf u}_k)$.\\
$\bar{\mathbf x}\leftarrow 0.8\,\mathbf x^{\ast}+0.2\,\bar{\mathbf x}$, \ 2 iterations, apply $\xU{\mathbf u_0}$, shift.}
\restr{each QP trusts its linearization only near the previous guess.}
\end{eqside}
\end{panel}

% =====================================================================
\clearpage
\sect{Offline \quad TUM \texttt{global\_racetrajectory\_optimization} \quad (whole lap)}

% ---------------------------------------------------------------- B0
\begin{panel}{B0}{Two pipelines}
\centering
\begin{tikzpicture}[
  bx/.style={draw=black!35,rounded corners=2pt,inner sep=3.5pt,font=\footnotesize,align=center,fill=#1!9},
  fl/.style={->,black!55,line width=0.6pt}]
  \node[font=\sffamily\small\bfseries,anchor=west] at (-1.8,0.95) {mincurv};
  \node[font=\footnotesize\itshape,text=black!50,anchor=west] at (-0.2,0.95) {path first, speed second};
  \node[bx=cR] (t1) at (0,0) {$\xR{\mathbf p_i,\ \mathbf n_i}$\\ $\xR{w_{\ell,i},\ w_{r,i}}$};
  \node[bx=cS] (q1) at (3.4,0) {QP over $\xS{\boldsymbol\alpha}$\\ $\min\sum\kappa_i^2$};
  \node[bx=cD] (r1) at (6.6,0) {$\xD{\mathbf r(s)},\ \xD{\kappa(s)}$};
  \node[bx=cS] (v1) at (10.2,0) {fwd/bwd pass\\ $\xS{v(s)},\ \xD{a_x(s)}$};
  \node[bx=cP] (g1) at (10.2,-1.3) {$\text{ggv}(v),\ a_{x,\mathrm{mach}}(v),\ c_d,\ m,\ \mu$};
  \draw[fl] (t1) -- (q1); \draw[fl] (q1) -- (r1); \draw[fl] (r1) -- (v1); \draw[fl] (g1) -- (v1);
  \draw[fl,dashed] (r1.north) to[bend right=35] node[sm,above]{IQP: re-linearize} (q1.north);
  \node[font=\sffamily\small\bfseries,anchor=west] at (-1.8,-2.35) {mintime};
  \node[font=\footnotesize\itshape,text=black!50,anchor=west] at (-0.2,-2.35) {path and speed in one problem};
  \node[bx=cR] (t2) at (0,-3.4) {$\xR{\kappa_{\mathrm{ref}}(s)}$\\ $\xR{w_\ell(s),\ w_r(s)}$};
  \node[bx=cS] (n2) at (5.0,-3.4) {NLP (collocation + IPOPT)\\ $\min\int_0^L \tfrac{dt}{ds}\,ds$\\ over $\xS{\mathbf x(s)},\ \xU{\mathbf u(s)}$};
  \node[bx=cD] (o2) at (10.2,-3.4) {$\xD{\mathbf r(s)},\ \xS{v(s)}$};
  \node[bx=cP] (p2) at (5.0,-4.95) {double track, Pacejka, load transfer, aero, $\mu$};
  \draw[fl] (t2) -- (n2); \draw[fl] (n2) -- (o2); \draw[fl] (p2) -- (n2);
\end{tikzpicture}
\restr{\texttt{mincurv}: speed never feeds back into geometry. \texttt{mintime}: geometry and speed traded off jointly.}
\end{panel}

% ---------------------------------------------------------------- B1
\begin{panel}{B1}{mincurv: the only variable is a lateral shift $\alpha_i$}
\begin{figside}
\begin{tikzpicture}
  \begin{scope}[shift={(0,-5.65)}]
    \draw[bnd] (121:7) arc (121:59:7);
    \draw[bnd] (121:5) arc (121:59:5);
    \draw[refl,dash pattern=on 4pt off 2pt] (121:6) arc (121:59:6);
    \draw[cS,line width=1pt] plot[smooth] coordinates {(121:6.62) (115:6.6) (105:6.15) (95:5.45) (85:5.45) (75:6.15) (65:6.6) (59:6.62)};
    \foreach \t/\a in {115/-0.6,105/-0.15,95/0.55,85/0.55,75/-0.15,65/-0.6} {
      \draw[->,black!40,line width=0.5pt] (\t:6) -- (\t:5.5);
      \node[dot,cR] at (\t:6) {};
      \pgfmathsetmacro\rr{6-\a}
      \node[dot,cS] at (\t:\rr) {};
    }
    \draw[cL!60,line width=3pt,opacity=0.5] (77.5:6.9) -- (77.5:5.1);
    \draw[cL,line width=0.6pt] ($(77.5:6.9)+(167.5:0.12)$) -- ($(77.5:6.9)+(-12.5:0.12)$);
    \draw[cL,line width=0.6pt] ($(77.5:5.1)+(167.5:0.12)$) -- ($(77.5:5.1)+(-12.5:0.12)$);
    \node[sm,text=cL,right=3pt] at (77.5:5.25) {$[-\bar w_\ell,\ \bar w_r]$};
    \node[sm,text=cR,above] at (95:6.02) {$\mathbf p_i$};
    \node[sm,text=cS,below] at (95:5.42) {$\mathbf r_i$};
    \node[sm,text=black!55,left] at (115:5.6) {$\mathbf n_i$};
    \node[lab,text=cS,right] at (91.5:5.75) {$\alpha_i$};
    \draw[dim] (119:6) -- (119:7) node[sm,midway,left]{$w_\ell$};
    \draw[dim] (119:6) -- (119:5) node[sm,midway,left]{$w_r$};
    \draw[->,cR,line width=0.8pt] (63:6) arc (63:58:6);
  \end{scope}
\end{tikzpicture}
\end{figside}
\begin{eqside}
\[\xD{\mathbf r_i}=\xR{\mathbf p_i}+\xS{\alpha_i}\,\xR{\mathbf n_i}\]
\[-\big(\xR{w_{\ell,i}}-\tfrac{\xP{w_{\mathrm{veh}}}}{2}\big)\ \le\ \xS{\alpha_i}\ \le\ \xR{w_{r,i}}-\tfrac{\xP{w_{\mathrm{veh}}}}{2}\]
{\footnotesize $\xR{\mathbf n_i}$ points \emph{right};\quad $\bar w=w-w_{\mathrm{veh}}/2$;\quad decision vector $\xS{\boldsymbol\alpha}\in\mathbb R^{N}$.}
\deps{$\mathbf r_i\leftarrow\xS{\alpha_i}$ and fixed track data}
\restr{each point slides only along its own normal, inside the track minus half the car width.}
\end{eqside}
\end{panel}

% ---------------------------------------------------------------- B2
\begin{panel}{B2}{mincurv: curvature, linearized}
\begin{figside}
\begin{tikzpicture}[scale=1.05]
  \draw[cR,line width=0.8pt,dash pattern=on 4pt off 2pt,domain=-2.2:2.5,samples=50,smooth] plot (\x,{0.22*\x*\x-0.75});
  \draw[cS,line width=1pt,domain=-2.2:2.5,samples=50,smooth] plot (\x,{0.22*\x*\x});
  \coordinate (Ar) at (1.2,-0.4332);
  \coordinate (A) at (1.2,0.3168);
  \draw[vec,black!55] (Ar) -- ++(0.95,0.5016) node[sm,right]{$(\xP{x'},\xP{y'})$ frozen};
  \node[dot,cR] at (Ar) {};
  \draw[vec,cS] (A) -- ++(0,0.85) node[sm,above]{$(\xS{x''},\xS{y''})$};
  \draw[guide] (A) -- ++(-0.7004,1.3265) node[sm,left]{$1/\kappa_i$};
  \draw[cD,dashed,line width=0.6pt] (-0.335,3.224) ++(-98:3.287) arc (-98:-26:3.287);
  \node[dot,cS] at (A) {};
  \node[sm,text=cS,right=2pt] at (A) {$\mathbf r_i$};
  \node[sm,text=cR] at (-1.6,-0.05) {ref.};
  \node[sm,text=cS] at (-1.75,1.25) {raceline};
\end{tikzpicture}
\end{figside}
\begin{eqside}
\[\xD{\kappa_i}=\frac{\xP{x_i'}\,\xS{y_i''}-\xP{y_i'}\,\xS{x_i''}}{\big(\xP{x_i'}^2+\xP{y_i'}^2\big)^{3/2}}
\ \Rightarrow\ \xD{\boldsymbol\kappa}\approx\xR{\boldsymbol\kappa_{\mathrm{ref}}}+\xP{E_\kappa}\,\xS{\boldsymbol\alpha}\]
\[\min_{\xS{\boldsymbol\alpha}}\ \sum_i\xD{\kappa_i}^2=\tfrac12\xS{\boldsymbol\alpha}^{\!\top}\!H\xS{\boldsymbol\alpha}+\mathbf f^{\top}\xS{\boldsymbol\alpha}
\quad\text{s.t.}\ \ |\xD{\kappa_i}|\le\xL{\kappa_{\max}}\ \ \&\ \text{B1}\]
{\footnotesize $x',y'$ taken from the reference spline (frozen); $x'',y''$ are linear in $\boldsymbol\alpha$ through the spline system. IQP: refreeze at the new line, re-solve.}
\deps{$\kappa_i\leftarrow\xS{\boldsymbol\alpha}$ (neighbours too, via splines)}
\restr{a QP: smoothest line through the corridor; no speed, mass or tire in it.}
\end{eqside}
\end{panel}

% ---------------------------------------------------------------- B3
\begin{panel}{B3}{Velocity profile: acceleration left after cornering}
\begin{figside}
\begin{tikzpicture}
  \begin{scope}
    \clip (-2.0,-1.6) rectangle (2.0,0.8);
    \fill[cS!7] (0,0) ellipse (1.6 and 1.3);
  \end{scope}
  \draw[ax] (-2.1,0) -- (2.4,0) node[sm,right]{$a_y$};
  \draw[ax] (0,-1.7) -- (0,1.85) node[sm,above]{$a_x$};
  \draw[cL,line width=0.9pt] (0,0) ellipse (1.6 and 1.3);
  \draw[cL!70,dashed,line width=0.7pt] (1.6,0) -- (0,1.3) -- (-1.6,0) -- (0,-1.3) -- cycle;
  \draw[liml] (-1.9,0.8) -- (1.9,0.8) node[sm,right]{$a_{x,\mathrm{mach}}(v)$};
  \draw[cD,line width=0.8pt] (0.75,-1.6) -- (0.75,1.6) node[sm,above]{$a_y{=}v^2\kappa$};
  \node[odot,draw=cD] at (0.75,1.148) {};
  \node[dot,cD] at (0.75,-1.148) {};
  \node[dot,cD] at (0.75,0.8) {};
  \node[sm,text=cL,anchor=north west] at (1.62,-0.18) {$\mu a_{y,\max}$};
  \node[sm,text=cL,above left] at (0,1.3) {$\mu a_{x,\max}$};
  \node[sm,text=cL] at (-1.45,1.05) {$p{=}2$};
  \node[sm,text=cL!70] at (-0.55,-0.6) {$p{=}1$};
\end{tikzpicture}
\end{figside}
\begin{eqside}
\[\xD{a_y}=\xS v^2\,\xD{\kappa}\]
\[\xD{a_x^{\mathrm{avail}}}=\xL{\mu\,a_{x,\max}(v)}\Big[1-\Big(\frac{\xD{a_y}}{\xL{\mu\,a_{y,\max}(v)}}\Big)^{\xP p}\Big]^{1/\xP p}\]
\[\begin{aligned}
\text{accel: }&\min\!\big(\xD{a_x^{\mathrm{avail}}},\,\xL{a_{x,\mathrm{mach}}(v)}\big)-\tfrac{\xP{c_d}}{\xP m}\xS v^2\\
\text{brake: }&\xD{a_x^{\mathrm{avail}}}+\tfrac{\xP{c_d}}{\xP m}\xS v^2\quad(\text{magnitude})
\end{aligned}\]
{\footnotesize ggv table: $(v,a_{x,\max},a_{y,\max})$; $p$ = \texttt{dyn\_model\_exp} $\in[1,2]$ (default 1).}
\deps{$a_x\leftarrow\xS v,\ \xD\kappa$ (from B2)}
\restr{longitudinal accel.\ allowed by grip after cornering, and by the motor.}
\end{eqside}
\end{panel}

% ---------------------------------------------------------------- B4
\begin{panel}{B4}{Velocity profile: forward--backward pass}
\begin{figside}
\begin{tikzpicture}[x=0.72cm,y=0.36cm]
  \draw[ax] (0,0) -- (10.7,0) node[sm,right]{$s$};
  \draw[ax] (0,0) -- (0,7.2) node[sm,above]{$v$};
  \draw[cL,dash pattern=on 3pt off 2pt,line width=0.6pt,domain=0:10,samples=150,smooth]
     plot (\x,{6-4*exp(-(((\x-3)/0.6)^2))-3*exp(-(((\x-7.5)/0.6)^2))});
  \draw[cS!45,dashed,line width=0.7pt,domain=3:7.5,samples=40] plot (\x,{sqrt(4+5*(\x-3))});
  \draw[cS!45,dashed,line width=0.7pt,domain=7.5:10,samples=30] plot (\x,{sqrt(9+5*(\x-7.5))});
  \draw[cL!55,dashed,line width=0.7pt,domain=0.3:3,samples=30] plot (\x,{sqrt(4+8*(3-\x))});
  \draw[cL!55,dashed,line width=0.7pt,domain=4.2:7.5,samples=30] plot (\x,{sqrt(9+8*(7.5-\x))});
  \draw[cS,line width=1.3pt,domain=0:10,samples=200]
     plot (\x,{min(6-4*exp(-(((\x-3)/0.6)^2))-3*exp(-(((\x-7.5)/0.6)^2)),
                   (\x<3)*sqrt(4+8*abs(3-\x)) + (\x>=3)*(\x<7.5)*min(sqrt(4+5*abs(\x-3)),sqrt(9+8*abs(7.5-\x))) + (\x>=7.5)*sqrt(9+5*abs(\x-7.5)))});
  \node[sm,text=cL,anchor=south] at (5.2,6.05) {$v_\kappa=\sqrt{\mu a_{y,\max}/|\kappa|}$};
  \node[sm,text=cS!70] at (5.35,3.3) {fwd};
  \node[sm,text=cL!80] at (6.45,4.6) {bwd};
  \foreach \s in {3,7.5} \draw[black!40] (\s,-0.25) -- (\s,0.25);
  \node[sm,below] at (3,-0.2) {apex};
  \node[sm,below] at (7.5,-0.2) {apex};
\end{tikzpicture}
\end{figside}
\begin{eqside}
\[\text{fwd: }\ \xS{v_{i+1}}=\sqrt{\xS{v_i}^2+2\,\xD{a_x^{+}}(\xS{v_i},\xD{\kappa_i})\,\xR{\ell_i}}\]
\[\text{bwd: }\ \xS{v_{i}}=\sqrt{\xS{v_{i+1}}^2+2\,|\xD{a_x^{-}}(\xS{v_{i+1}},\xD{\kappa_{i+1}})|\,\xR{\ell_i}}\]
\[\xS{v}=\min\big(v_\kappa,\ \text{fwd},\ \text{bwd}\big)\]
{\footnotesize $\ell_i$: segment length; $a_x^{\pm}$ from B3. No optimizer: a sweep.}
\deps{$v\leftarrow\xD\kappa(s)$, ggv, motor, drag}
\restr{speed $\le$ cornering limit, and reachable by accelerating out of / braking into each apex.}
\end{eqside}
\end{panel}

% ---------------------------------------------------------------- B5
\begin{panel}{B5}{mintime: who depends on whom}
\begin{figside}[0.56]
\begin{tikzpicture}[x=0.55cm,y=0.37cm]
  \def\ncols{11}
  \mcol{1}{v}{cS}\mcol{2}{\beta}{cS}\mcol{3}{\omega_z}{cS}\mcol{4}{n}{cS}\mcol{5}{\xi}{cS}
  \mcol{6}{\delta}{cU}\mcol{7}{F_d}{cU}\mcol{8}{F_b}{cU}\mcol{9}{\gamma_y}{cU}
  \mcol{10}{\kappa}{cR}\mcol{11}{w}{cR}
  \mrow{1}{\xD{S_f}}{1,2,4,5,10}{0}
  \mrow{2}{\xD{\alpha_{ij}}}{1,2,3,6}{0}
  \mrow{3}{\xD{F_{z,ij}}}{1,7,8,9}{0}
  \mrow{4}{\xD{F_{x,ij}}}{7,8}{0}
  \mrow{5}{\xD{F_{y,ij}}}{0}{1,2,3,6,7,8,9}
  \mrow{6.4}{\xS{v'}}{2,6}{1,3,4,5,7,8,9,10}
  \mrow{7.4}{\xS{\beta'}}{1,2,3,6}{4,5,7,8,9,10}
  \mrow{8.4}{\xS{\omega_z'}}{6}{1,2,3,4,5,7,8,9,10}
  \mrow{9.4}{\xS{n'}}{1,2,5}{4,10}
  \mrow{10.4}{\xS{\xi'}}{3,10}{1,2,4,5}
  \mrow{11.8}{T_{\mathrm{lap}}}{0}{1,2,4,5,10}
  \mrow{12.8}{\xL{\text{track}}}{4,11}{0}
  \mrow{13.8}{\xL{\text{circles}}}{0}{1,2,3,6,7,8,9}
  \mrow{14.8}{\xL{\gamma_y\ \text{eq.}}}{9}{1,2,3,6,7,8}
  \mrow{15.8}{\xL{\text{power}}}{1,7}{0}
  \mrow{16.8}{\xL{F_dF_b}}{7,8}{0}
  \mrow{17.8}{\xL{\text{rates}}}{6,7,8}{1,4,10}
  \mrow{18.8}{\xL{\text{periodic}}}{1,2,3,4,5}{0}
  \draw[black!30] (0.5,-5.7) -- (11.5,-5.7);
  \draw[black!30] (0.5,-11.1) -- (11.5,-11.1);
  \draw[black!30] (5.5,1.4) -- (5.5,-19.3);
  \draw[black!30] (9.5,1.4) -- (9.5,-19.3);
\end{tikzpicture}
\end{figside}
\begin{eqside}[0.40]
{\footnotesize
$(\cdot)'=d(\cdot)/ds$\par\medskip
\tikz\fill[black!85] (0,0) circle (2.5pt);\ appears in the formula\par
\tikz\draw[black!60,line width=0.6pt] (0,0) circle (2.2pt);\ enters through $S_f$, $\alpha$ or $F$\par\medskip
Every $d/ds$ row inherits $n,\xi,\kappa$ through $S_f=dt/ds$.}
\restr{compared with A0: no $X,Y$ states and no progress variable; position is $(s,n)$ on the fixed reference.}
\end{eqside}
\end{panel}

% ---------------------------------------------------------------- B6
\begin{panel}{B6}{mintime: curvilinear coordinates, time as a function of $s$}
\begin{figside}
\begin{tikzpicture}
  \coordinate (O) at (0,0);
  \coordinate (P) at (0:4);
  \coordinate (C) at (0:2.85);
  \draw[refl] (-32:4) arc (-32:42:4) node[sm,above]{ref.};
  \node[dot,black!55] at (O) {};
  \draw[guide] (O) -- (-22:4) node[sm,pos=0.33,below,text=cR]{$1/\kappa(s)$};
  \node[dot,cR] at (P) {};
  \node[sm,text=cR,right] at (P) {$s$};
  \draw[guide] (2.85,-1.0) -- (2.85,2.4);
  \draw[guide] (4,-0.25) -- (4,-1.15);
  \draw[guide] (C) -- ++(108:2.3);
  \draw[dim] (2.85,-1.1) -- (4,-1.1) node[sm,midway,below,text=cS]{$n$};
  \begin{scope}[shift={(C)},rotate=108]
    \draw[body] (-0.72,-0.28) rectangle (0.72,0.28);
    \fill[black!55] (0.72,0) -- (0.52,0.12) -- (0.52,-0.12) -- cycle;
  \end{scope}
  \draw[vS] (C) -- ++(122:2.15) node[lab,above]{$v$};
  \draw[cS,line width=0.6pt] ($(C)+(90:1.3)$) arc (90:108:1.3);
  \node[lab,text=cS] at ($(C)+(99:1.08)$) {$\xi$};
  \draw[cS,line width=0.6pt] ($(C)+(108:1.75)$) arc (108:122:1.75);
  \node[lab,text=cS] at ($(C)+(115:2.0)$) {$\beta$};
  \draw[cS,->,line width=0.7pt] ($(C)+(195:0.95)$) arc (195:255:0.95);
  \node[lab,text=cS] at ($(C)+(222:1.22)$) {$\omega_z$};
\end{tikzpicture}
\end{figside}
\begin{eqside}
\[\xD{S_f}\equiv\frac{dt}{ds}=\frac{1-\xS n\,\xR\kappa}{\xS v\cos(\xS\xi+\xS\beta)}\]
\[\frac{d\xS n}{ds}=\xD{S_f}\,\xS v\sin(\xS\xi+\xS\beta),\qquad \frac{d\xS\xi}{ds}=\xD{S_f}\,\xS{\omega_z}-\xR\kappa\]
\[T_{\mathrm{lap}}=\int_0^{L}\xD{S_f}\,ds\]
{\footnotesize $\xS{\mathbf x}=(v,\beta,\omega_z,n,\xi)$,\ $\xU{\mathbf u}=(\delta,F_d,F_b,\gamma_y)$;\ \ $n>0$ = \emph{left} (opposite sign to $\alpha$ in B1).}
\deps{$S_f\leftarrow\xS{v,\beta,n,\xi},\ \xR\kappa$}
\restr{time is not a variable: lap time integrates $dt/ds$ over the fixed reference length $L$.}
\end{eqside}
\end{panel}

% ---------------------------------------------------------------- B7
\begin{panel}{B7}{mintime: double-track loads}
\begin{figside}
\begin{tikzpicture}
  \draw[body] (-0.55,-0.55) rectangle (3.85,0.55);
  \draw[black!35] (0,-0.8) -- (0,0.8);
  \draw[black!35] (3.3,-0.8) -- (3.3,0.8);
  \draw[cL!80,line width=0.6pt,fill=cL!6] (3.3,-0.8) circle (0.55);
  \draw[cL!80,line width=0.6pt,fill=cL!6] (3.3,0.8) circle (0.38);
  \draw[cL!80,line width=0.6pt,fill=cL!6] (0,-0.8) circle (0.40);
  \draw[cL!80,line width=0.6pt,fill=cL!6] (0,0.8) circle (0.24);
  \node[whl] at (0,0.8) {}; \node[whl] at (0,-0.8) {};
  \node[whl,rotate=15] at (3.3,0.8) {}; \node[whl,rotate=15] at (3.3,-0.8) {};
  \fill (1.54,0) circle (1.6pt);
  \draw[vD] (1.54,0) -- ++(-0.8,0) node[sm,left]{$a_x$};
  \draw[vD] (1.54,0) -- ++(0,0.5) node[sm,above]{$a_y$};
  \node[sm,text=cL] at (4.45,-0.8) {$\propto F_{z,fr}$};
  \draw[dim] (0,-1.55) -- (1.54,-1.55) node[sm,midway,below]{$l_r$};
  \draw[dim] (1.54,-1.55) -- (3.3,-1.55) node[sm,midway,below]{$l_f$};
  \draw[dim] (-1.0,-0.8) -- (-1.0,0.8) node[sm,midway,left]{$b$};
  \node[sm,text=black!50] at (1.65,1.15) {braking + left turn};
\end{tikzpicture}
\end{figside}
\begin{eqside}[0.46]
\[\xD{F_{z,ij}}=\xP{F^{\mathrm{st}}_{z,ij}}+\tfrac12\xP{c_{l,i}}\xS v^2
\mp\tfrac{\xP h}{2\xP l}\big(\xU{F_d}{+}\xU{F_b}{-}\xP{c_d}\xS v^2{-}\xP{F_{\mathrm{roll}}}\big)
\mp\xP{k_i}\,\xU{\gamma_y}\]
{\footnotesize long.: $-$front/$+$rear;\ lat.: $-$left/$+$right;\ $k_f{=}k_{\mathrm{roll}}$, $k_r{=}1{-}k_{\mathrm{roll}}$}
\[\xU{\gamma_y}=\frac{\xP h}{\xP{\bar b}}\,\xP m\,\xD{a_y}\quad(\text{equality constraint})\]
{\footnotesize $\xD{F_{x,ij}}$: drive/brake split $\xP{k_{\mathrm{drive}}},\xP{k_{\mathrm{brake}}}$ minus rolling resistance.}
\deps{$F_z\leftarrow\xS v,\ \xU{F_d,F_b,\gamma_y}$}
\restr{each wheel's load (hence grip) follows the accelerations; $\gamma_y$ is a decision variable closed by an equality.}
\end{eqside}
\end{panel}

% ---------------------------------------------------------------- B8
\begin{panel}{B8}{mintime: tire model and friction circles}
\begin{figside}
\begin{tikzpicture}
  \begin{scope}[xscale=7,yscale=1.6]
    \draw[ax] (0,0) -- (0.38,0) node[sm,right]{$\alpha$};
    \draw[ax] (0,0) -- (0,1.3) node[sm,above]{$F_y/D$};
    \draw[guide] (0,1) -- (0.36,1);
    \draw[cD,line width=1pt,domain=0:0.35,samples=150,smooth] plot (\x,{sin(2.5*atan(10*\x-1.0*(10*\x-rad(atan(10*\x)))))});
    \node[sm,left] at (0,1) {$1$};
  \end{scope}
  \begin{scope}[shift={(3.55,0)},xscale=0.25,yscale=0.25]
    \draw[ax] (0,0) -- (10,0) node[sm,right]{$F_z$};
    \draw[ax] (0,0) -- (0,9.6) node[sm,above]{$D$};
    \draw[guide] (0,0) -- (9,9) node[sm,above]{$\mu F_z$};
    \draw[cD,line width=1pt,domain=0:9,samples=30] plot (\x,{\x*(1-0.1*\x/3)});
    \node[sm,text=cD] at (8.2,4.9) {$\varepsilon<0$};
  \end{scope}
\end{tikzpicture}
\end{figside}
\begin{eqside}
\[\xD{D_{ij}}=\xP\mu\,\xD{F_{z,ij}}\Big(1+\xP\varepsilon\,\frac{\xD{F_{z,ij}}}{\xP{F_{z0}}}\Big)\]
\[\xD{F_{y,ij}}=\xD{D_{ij}}\sin\!\Big(\xP C\arctan\big(\xP B\xD\alpha-\xP E(\xP B\xD\alpha-\arctan\xP B\xD\alpha)\big)\Big)\]
\[\Big(\frac{\xD{F_{x,ij}}}{\xP\mu\,\xD{F_{z,ij}}}\Big)^2+\Big(\frac{\xD{F_{y,ij}}}{\xP\mu\,\xD{F_{z,ij}}}\Big)^2\le\xL{1}\qquad\forall\,ij\]
\deps{$\alpha_{ij}\leftarrow\xS{v,\beta,\omega_z},\xU\delta$;\ \ $F_{z,ij}$ from B7}
\restr{each tire inside its own circle sized by its load (B7 circles); load transfer costs total grip ($\varepsilon<0$).}
\end{eqside}
\end{panel}

% ---------------------------------------------------------------- B9
\begin{panel}{B9}{mintime: the NLP}
\begin{figside}[0.40]
\begin{tikzpicture}[x=1cm,y=1cm]
  \draw[ax] (-0.2,0) -- (5.3,0) node[sm,right]{$s$};
  \foreach \t in {0.1127,0.5,0.8873} \draw[black!25] ({0.4+4.2*\t},0) -- ({0.4+4.2*\t},2.2);
  \draw[cS,line width=1pt,domain=0:1,samples=40] plot ({0.4+4.2*\x},{1+2.4*\x-1.4*\x*\x});
  \foreach \t/\m in {0.1127/0.4963,0.5/0.2381,0.8873/-0.0201} {
    \pgfmathsetmacro\xx{0.4+4.2*\t}
    \pgfmathsetmacro\yy{1+2.4*\t-1.4*\t*\t}
    \draw[cD,line width=1.1pt] ({\xx-0.6},{\yy-0.6*\m+0.12}) -- ({\xx+0.6},{\yy+0.6*\m+0.12});
    \node[dot,cS] at (\xx,\yy) {};
  }
  \node[dot,cS] at (0.4,1) {};
  \node[odot,draw=cS] at (4.6,2) {};
  \node[sm,text=cS,left] at (0.4,1) {$X_k$};
  \node[sm,text=cS,right] at (4.6,2) {$X_{k+1}$};
  \node[sm,text=cD,above] at (2.5,2.3) {slope $=h_k\mathbf f(X_{k,j},U_k,\kappa)$};
  \draw[cU,line width=1pt] (-0.1,-0.75) -- (0.4,-0.75) (0.4,-0.5) -- (4.6,-0.5) (4.6,-0.85) -- (5.1,-0.85);
  \node[sm,text=cU,above] at (2.5,-0.5) {$U_k$};
  \draw[dim] (0.4,-1.15) -- (4.6,-1.15) node[sm,midway,below]{$h_k$};
\end{tikzpicture}
\end{figside}
\begin{eqside}[0.58]
\[\min_{\xS{\mathbf x},\xU{\mathbf u}}\ \int_0^{L}\xD{S_f}\,ds+\xP{\lambda_F}\|\Delta\xU F\|^2+\xP{\lambda_\delta}\|\Delta\xU\delta\|^2\]
\footnotesize
\[\begin{aligned}
&\text{dynamics:}  && h_k\,\mathbf f(X_{k,j},\xU{U_k},\xR{\kappa})=\textstyle\sum_r C_{rj}X_{k,r},\ \ 3\text{ Legendre pts}\\
&\text{track:}     && -\xR{w_r}+\tfrac{\xP{w_{\mathrm{opt}}}}{2}\le\xS n\le\xR{w_\ell}-\tfrac{\xP{w_{\mathrm{opt}}}}{2}\\
&\text{grip, loads:}&& \text{B8 circles},\ \ \gamma_y\ \text{equality (B7)}\\
&\text{power:}     && \xS v\,\xU{F_d}\le\xL{P_{\max}},\qquad -\epsilon\le\xU{F_d}\,\xU{F_b}\le 0\\
&\text{actuators:} && \tfrac{\Delta\xU\delta}{\Delta t}\in\pm\tfrac{\xL{\delta_{\max}}}{\xP{t_\delta}},\ \
                      \tfrac{\Delta\xU{F_d}}{\Delta t}\le\tfrac{\xL{F_{d,\max}}}{\xP{t_d}},\ \
                      \tfrac{\Delta\xU{F_b}}{\Delta t}\ge-\tfrac{\xL{F_{b,\max}}}{\xP{t_b}}\\
&\text{closed lap:}&& \xS{\mathbf x}(0)=\xS{\mathbf x}(L)\\
&\text{bounds:}    && \xS v\ge 1,\ \ |\xS\beta|,|\xS{\omega_z}|,|\xS\xi|\le\tfrac{\pi}{2},\ \ \xU{F_d}\ge0\ge\xU{F_b}
\end{aligned}\]
\small
\restr{dynamics enforced only at collocation points; the lap must close; everything else is a path constraint.}
\end{eqside}
\end{panel}
}
. Requires amsmath, amssymb, array.
% =====================================================================
\usepackage{xcolor}
\usepackage{tikz}
\usetikzlibrary{arrows.meta,calc,decorations.pathmorphing,decorations.pathreplacing}

% ---------------- role colours ----------------------------------------
\definecolor{cS}{HTML}{0072B2}  % decision state
\definecolor{cU}{HTML}{C76E00}  % control input
\definecolor{cR}{HTML}{00876C}  % reference / track data
\definecolor{cP}{HTML}{7A7A7A}  % fixed parameter
\definecolor{cD}{HTML}{8E44AD}  % derived quantity
\definecolor{cL}{HTML}{C0392B}  % limit / bound
\newcommand{\xS}[1]{{\color{cS}#1}}
\newcommand{\xU}[1]{{\color{cU}#1}}
\newcommand{\xR}[1]{{\color{cR}#1}}
\newcommand{\xP}[1]{{\color{cP}#1}}
\newcommand{\xD}[1]{{\color{cD}#1}}
\newcommand{\xL}[1]{{\color{cL}#1}}

% ---------------- layout ----------------------------------------------
\newenvironment{panel}[2]{%
  \par\addvspace{5pt}\noindent\begin{minipage}{\linewidth}%
  {\sffamily\small\bfseries\color{black!70}#1}\enspace{\sffamily\small #2}\par\vspace{4pt}}%
  {\end{minipage}\par\vspace{6pt}\noindent{\color{black!12}\rule{\linewidth}{0.5pt}}\par}
\newenvironment{figside}[1][0.53]{\begin{minipage}[c]{#1\linewidth}\centering}{\end{minipage}\hfill}
\newenvironment{eqside}[1][0.45]{\begin{minipage}[c]{#1\linewidth}\small
  \setlength{\abovedisplayskip}{3pt}\setlength{\belowdisplayskip}{3pt}%
  \setlength{\abovedisplayshortskip}{2pt}\setlength{\belowdisplayshortskip}{2pt}}{\end{minipage}}
\newcommand{\deps}[1]{\par\smallskip{\sffamily\scriptsize\bfseries\color{black!45}DEPENDS}\enspace{\small#1}}
\newcommand{\restr}[1]{\par\smallskip{\sffamily\scriptsize\bfseries\color{cL}RESTRAINS}\enspace{\small#1}}
\newcommand{\chip}[2]{\tikz[baseline=-0.6ex]\fill[#1] (0,0) circle (3pt);\,{\small\sffamily\color{#1}#2}\quad}
\providecommand{\sect}[1]{\par\vspace{8pt}{\large\sffamily\bfseries #1}\par\vspace{2pt}\noindent{\color{black!40}\rule{\linewidth}{0.8pt}}\par}

% ---------------- tikz styles -----------------------------------------
\tikzset{
  every picture/.style={line cap=round,line join=round,>={Stealth[length=4.5pt,width=3.2pt]}},
  ax/.style={->,black!45,line width=0.5pt},
  body/.style={draw=black!65,fill=black!5,line width=0.6pt,rounded corners=1.2pt},
  whl/.style={fill=black!75,draw=none,rounded corners=0.6pt,minimum width=5.5mm,minimum height=2mm,inner sep=0pt},
  vS/.style={->,cS,line width=0.9pt},
  vD/.style={->,cD,line width=0.9pt},
  vec/.style={->,line width=0.8pt},
  refl/.style={cR,line width=1pt},
  bnd/.style={black!55,line width=0.7pt},
  liml/.style={cL,line width=0.8pt},
  guide/.style={black!35,dash pattern=on 2pt off 1.5pt,line width=0.4pt},
  dim/.style={<->,black!50,line width=0.4pt,>={Stealth[length=3pt,width=2.4pt]}},
  lab/.style={font=\footnotesize,inner sep=1pt},
  sm/.style={font=\scriptsize,inner sep=1pt},
  dot/.style={circle,fill,inner sep=0pt,minimum size=3.4pt},
  odot/.style={circle,draw,fill=white,inner sep=0pt,minimum size=3.8pt,line width=0.6pt},
}
% small car icon: \caricon{(x,y)}{heading deg}{scale}
\newcommand{\caricon}[3]{%
  \begin{scope}[shift={#1},rotate=#2,scale=#3]
    \draw[body] (-0.45,-0.2) rectangle (0.45,0.2);
    \fill[black!55] (0.45,0) -- (0.28,0.11) -- (0.28,-0.11) -- cycle;
  \end{scope}}
% dependency-matrix row: {row}{label}{direct cols}{via cols}  (0 = none)
\newcommand{\mrow}[4]{%
  \draw[black!7,line width=3.6pt] (0.6,-#1) -- (\ncols+0.4,-#1);
  \node[anchor=east,font=\footnotesize,inner sep=1pt] at (0.45,-#1) {$#2$};
  \foreach \c in {#3}{\ifnum\c>0 \fill[black!85] (\c,-#1) circle (2.5pt);\fi}
  \foreach \c in {#4}{\ifnum\c>0 \draw[black!60,line width=0.6pt,fill=white] (\c,-#1) circle (2.2pt);\fi}}
\newcommand{\mcol}[3]{\node[font=\footnotesize,text=#3] at (#1,0.85) {$#2$};}

% =====================================================================
.
%  \sect{} renders the two part headings; redefine it before % =====================================================================
%  preamble_maps.tex -- packages, colours, macros and TikZ styles for the
%  equation maps. \input from the preamble of any document that
%  % =====================================================================
%  background_maps.tex -- equation maps (body only).
%  Online: Liniger MPCC (C++ master). Offline: TUM global_racetrajectory_optimization.
%  \input inside a document whose preamble \input{preamble_maps}.
%  \sect{} renders the two part headings; redefine it before \input{preamble_maps}
%  to change how they look (e.g. as \subsection).
% =====================================================================
{\setlength{\parindent}{0pt}%
\chip{cS}{decision state}\chip{cU}{control input}\chip{cR}{reference / track data}%
\chip{cP}{fixed parameter}\chip{cD}{derived}\chip{cL}{limit / bound}

% =====================================================================
\sect{Online \quad MPCC \quad (receding horizon, 60 steps $\times$ 20 ms)}

% ---------------------------------------------------------------- A0
\begin{panel}{A0}{Variables and who depends on whom}
\begin{figside}[0.60]
\begin{tikzpicture}[x=0.5cm,y=0.37cm]
  \def\ncols{14}
  \mcol{1}{X}{cS}\mcol{2}{Y}{cS}\mcol{3}{\varphi}{cS}\mcol{4}{v_x}{cS}\mcol{5}{v_y}{cS}
  \mcol{6}{r}{cS}\mcol{7}{s}{cS}\mcol{8}{D}{cS}\mcol{9}{\delta}{cS}\mcol{10}{v_s}{cS}
  \mcol{11}{\dot D}{cU}\mcol{12}{\dot\delta}{cU}\mcol{13}{\dot v_s}{cU}\mcol{14}{\mathbf p_{\!\mathrm{ref}}}{cR}
  \mrow{1}{\xD{\alpha_f}}{4,5,6,9}{0}
  \mrow{2}{\xD{\alpha_r}}{4,5,6}{0}
  \mrow{3}{\xD{F_{f,y}}}{0}{4,5,6,9}
  \mrow{4}{\xD{F_{r,y}}}{0}{4,5,6}
  \mrow{5}{\xD{F_{r,x}}}{4,8}{0}
  \mrow{6}{\xD{F_{\mathrm{fric}}}}{4}{0}
  \mrow{7.4}{\xS{\dot X,\dot Y}}{3,4,5}{0}
  \mrow{8.4}{\xS{\dot\varphi}}{6}{0}
  \mrow{9.4}{\xS{\dot v_x}}{5,6,9}{4,8}
  \mrow{10.4}{\xS{\dot v_y}}{4,6,9}{5}
  \mrow{11.4}{\xS{\dot r}}{9}{4,5,6}
  \mrow{12.4}{\xS{\dot s}}{10}{0}
  \mrow{13.4}{\xS{\dot D,\dot\delta,\dot v_s}}{11,12,13}{0}
  \mrow{14.8}{\xD{e_c,\,e_l}}{1,2,7,14}{0}
  \mrow{15.8}{\xL{\text{track}}}{1,2,14}{0}
  \mrow{16.8}{\xL{\text{ellipse}}}{0}{4,5,6,8}
  \mrow{17.8}{\xL{|\alpha_f|}}{0}{4,5,6,9}
  \mrow{18.8}{\beta\ \text{cost}}{4,5}{0}
  \mrow{19.8}{\varphi\ \text{cost}}{3,7,14}{0}
  \mrow{20.8}{\text{progress}}{10}{0}
  \draw[black!30] (0.5,-6.7) -- (14.5,-6.7);
  \draw[black!30] (0.5,-14.1) -- (14.5,-14.1);
  \draw[black!30] (10.5,1.4) -- (10.5,-21.3);
\end{tikzpicture}
\end{figside}
\begin{eqside}[0.38]
\[\xS{\mathbf x} = [\,\xS{X\ Y\ \varphi\ v_x\ v_y\ r\ s\ D\ \delta\ v_s}\,]^{\top}\]
\[\xU{\mathbf u} = [\,\xU{\dot D\ \ \dot\delta\ \ \dot v_s}\,]^{\top}\]
\[\xS{\dot s}=\xS{v_s},\quad \xS{\dot D}=\xU{u_1},\quad \xS{\dot\delta}=\xU{u_2},\quad \xS{\dot v_s}=\xU{u_3}\]
\medskip
{\footnotesize
\tikz\fill[black!85] (0,0) circle (2.5pt);\ appears in the formula\par
\tikz\draw[black!60,line width=0.6pt] (0,0) circle (2.2pt);\ enters through $\alpha$ or $F$\par}
\restr{inputs are \emph{rates}, so rate costs and slew bounds act directly on $\xU{\mathbf u}$.}
\end{eqside}

\par\medskip
{\footnotesize\renewcommand{\arraystretch}{1.12}%
\begin{minipage}[t]{0.335\linewidth}
{\sffamily\scriptsize\bfseries\color{cS}STATE $\mathbf x$}\par\smallskip
\begin{tabular}{@{}l@{\ \ }>{\raggedright\arraybackslash}p{4.35cm}@{}}
$\xS{X,Y}$       & car position, world frame [m]\\
$\xS{\varphi}$   & heading (yaw angle), world frame [rad]\\
$\xS{v_x}$       & forward velocity, body frame [m/s]\\
$\xS{v_y}$       & sideways velocity, body frame [m/s]\\
$\xS{r}$         & yaw rate [rad/s]\\
$\xS{s}$         & progress: virtual arc length along the reference [m]\\
$\xS{D}$         & motor duty cycle, $[-0.1,1]$; $<0$ brakes\\
$\xS{\delta}$    & front steering angle, $|\delta|\le 0.35$ rad\\
$\xS{v_s}$       & progress speed $\dot s\ge 0$ [m/s]
\end{tabular}
\end{minipage}\hfill
\begin{minipage}[t]{0.32\linewidth}
{\sffamily\scriptsize\bfseries\color{cU}INPUT $\mathbf u$}\par\smallskip
\begin{tabular}{@{}l@{\ \ }>{\raggedright\arraybackslash}p{4.05cm}@{}}
$\xU{\dot D}$         & duty-cycle rate\\
$\xU{\dot\delta}$     & steering rate\\
$\xU{\dot v_s}$       & progress acceleration
\end{tabular}\par\medskip
{\sffamily\scriptsize\bfseries\color{cR}REFERENCE}\par\smallskip
\begin{tabular}{@{}l@{\ \ }>{\raggedright\arraybackslash}p{4.05cm}@{}}
$\xR{\mathbf p_{\mathrm{ref}}}$ & centreline spline $(X_{\mathrm{ref}}(s),Y_{\mathrm{ref}}(s))$ with heading $\Phi(s)$; fixed data
\end{tabular}\par\medskip
{\sffamily\scriptsize\bfseries\color{cD}DERIVED}\par\smallskip
\begin{tabular}{@{}l@{\ \ }>{\raggedright\arraybackslash}p{4.05cm}@{}}
$\xD{\alpha_f,\alpha_r}$       & front / rear tire slip angle\\
$\xD{F_{f,y},F_{r,y}}$         & front / rear lateral tire force\\
$\xD{F_{r,x}}$                 & rear drive force from the motor\\
$\xD{F_{\mathrm{fric}}}$       & rolling resistance + air drag
\end{tabular}
\end{minipage}\hfill
\begin{minipage}[t]{0.31\linewidth}
{\sffamily\scriptsize\bfseries\color{black!60}COST / CONSTRAINT}\par\smallskip
\begin{tabular}{@{}l@{\ \ }>{\raggedright\arraybackslash}p{3.55cm}@{}}
$\dot{(\,\cdot\,)}$        & time derivative: the state equations\\
$\xD{e_c}$                 & contouring error: sideways offset from $\mathbf p_{\mathrm{ref}}(s)$\\
$\xD{e_l}$                 & lag error: along-track gap between $s$ and the car\\
\xL{track}                 & car inside the slab around the reference\\
\xL{ellipse}               & rear friction ellipse\\
$\xL{|\alpha_f|}$          & front slip-angle limit\\
$\beta$ cost               & sideslip $\arctan(v_y/v_x)$ penalty\\
$\varphi$ cost             & heading vs.\ track heading $\Phi(s)$\\
progress                   & reward $-q_{v_s}v_s$
\end{tabular}
\end{minipage}}
\end{panel}

% ---------------------------------------------------------------- A1
\begin{panel}{A1}{Body velocity $\to$ world pose}
\begin{figside}
\begin{tikzpicture}
  \draw[ax] (0,0) -- (1.0,0) node[sm,right]{$X$};
  \draw[ax] (0,0) -- (0,1.0) node[sm,above]{$Y$};
  \coordinate (C) at (2.9,0.45);
  % world components of v (dotted)
  \draw[black!55,dotted,line width=0.9pt] (C) -- ++(1.282,0) node[sm,below,midway,yshift=-2pt]{$\dot X$}
        -- ++(0,1.779) node[sm,right,pos=0.8]{$\dot Y$};
  % heading angle, measured behind the car
  \draw[guide] (C) -- ++(180:1.5);
  \draw[guide] (C) -- ++(210:1.5);
  \draw[cS,line width=0.6pt] ($(C)+(180:1.15)$) arc (180:210:1.15);
  \node[lab,text=cS] at ($(C)+(195:1.42)$) {$\varphi$};
  \begin{scope}[shift={(C)},rotate=30]
    \draw[body] (-0.55,-0.22) rectangle (0.55,0.22);
    \fill[black!55] (0.55,0) -- (0.4,0.1) -- (0.4,-0.1) -- cycle;
    \draw[vS] (0,0) -- (2.0,0) node[lab,below right=0pt]{$v_x$};
    \draw[vS] (0,0) -- (0,0.9) node[lab,above left=0pt]{$v_y$};
    \draw[->,black!75,line width=0.6pt] (0,0) -- (2.0,0.9) node[lab,above=1pt]{$\mathbf v$};
  \end{scope}
  \draw[cS,->,line width=0.7pt] ($(C)+(245:0.7)$) arc (245:305:0.7);
  \node[lab,text=cS] at ($(C)+(275:0.92)$) {$r$};
\end{tikzpicture}
\end{figside}
\begin{eqside}
\[\begin{aligned}
\xS{\dot X} &= \xS{v_x}\cos\xS{\varphi} - \xS{v_y}\sin\xS{\varphi}\\
\xS{\dot Y} &= \xS{v_x}\sin\xS{\varphi} + \xS{v_y}\cos\xS{\varphi}\\
\xS{\dot\varphi} &= \xS{r}
\end{aligned}\]
\deps{$(\dot X,\dot Y)\leftarrow \xS{v_x,v_y,\varphi}$;\quad $\dot\varphi\leftarrow\xS{r}$}
\restr{pose = integral of the body velocity rotated by the heading.}
\end{eqside}
\end{panel}

% ---------------------------------------------------------------- A2
\begin{panel}{A2}{Slip angles}
\begin{figside}
\begin{tikzpicture}
  \coordinate (R) at (0,0); \coordinate (G) at (1.5,0); \coordinate (F) at (3.2,0);
  \draw[black!40,line width=1.4pt] (R) -- (F);
  \draw[guide] (F) -- ++(0:1.7);
  \draw[guide] ($(F)+(208:0.9)$) -- ++(28:2.7);
  \node[whl] at (R) {};
  \node[whl,rotate=28] at (F) {};
  \fill (G) circle (1.6pt);
  % rear
  \draw[vS] (R) -- ++(-16:1.55) node[sm,below=2pt]{$(v_x,\ v_y{-}l_r r)$};
  \draw[cD,line width=0.8pt] ($(R)+(-16:1.1)$) arc (-16:0:1.1);
  \node[lab,text=cD] at ($(R)+(-6:1.35)$) {$\alpha_r$};
  % front
  \draw[vS] (F) -- ++(10:1.6) node[sm,right]{$(v_x,\ v_y{+}l_f r)$};
  \draw[cD,line width=0.8pt] ($(F)+(10:1.25)$) arc (10:28:1.25);
  \node[lab,text=cD] at ($(F)+(19:1.5)$) {$\alpha_f$};
  \draw[cS,line width=0.6pt] ($(F)+(180:0.75)$) arc (180:208:0.75);
  \node[lab,text=cS] at ($(F)+(194:1.0)$) {$\delta$};
  % dims
  \draw[dim] (0,-1.05) -- (1.5,-1.05) node[sm,midway,below]{$l_r$};
  \draw[dim] (1.5,-1.05) -- (3.2,-1.05) node[sm,midway,below]{$l_f$};
\end{tikzpicture}
\end{figside}
\begin{eqside}
\[\begin{aligned}
\xD{\alpha_f} &= \xS{\delta} - \arctan\frac{\xS{v_y}+\xP{l_f}\,\xS{r}}{\xS{v_x}}\\[2pt]
\xD{\alpha_r} &= -\arctan\frac{\xS{v_y}-\xP{l_r}\,\xS{r}}{\xS{v_x}}
\end{aligned}\]
\deps{$\alpha_f\leftarrow\xS{v_x,v_y,r,\delta}$;\quad $\alpha_r\leftarrow\xS{v_x,v_y,r}$}
\restr{nothing alone: angle between where a tire points and where it moves (bounded in A8).}
\end{eqside}
\end{panel}

% ---------------------------------------------------------------- A3
\begin{panel}{A3}{Tire and drivetrain forces}
\begin{figside}[0.56]
\begin{tikzpicture}
  \begin{scope}[xscale=1.35,yscale=7]
    \fill[cL!10] (0.6,-0.215) rectangle (1.65,0.215);
    \fill[cL!10] (-1.65,-0.215) rectangle (-0.6,0.215);
    \draw[ax] (-1.7,0) -- (1.8,0) node[sm,right]{$\alpha$};
    \draw[ax] (0,-0.22) -- (0,0.25) node[sm,above]{$F_y$};
    \draw[guide] (-1.65,0.192) -- (1.65,0.192);
    \node[sm,text=cP,anchor=south east] at (-0.05,0.193) {$D_f$};
    \draw[black!50,line width=0.5pt] (-0.15,-0.0891) -- (0.26,0.1545);
    \node[sm,text=cP,anchor=east] at (-0.05,0.12) {$B_fC_fD_f$};
    \draw[cD,line width=1pt,domain=-1.6:1.6,samples=100,smooth] plot (\x,{0.192*sin(1.2*atan(2.579*\x))});
    \draw[liml] (0.6,-0.215) -- (0.6,0.215);
    \draw[liml] (-0.6,-0.215) -- (-0.6,0.215);
    \node[sm,text=cL,above] at (0.6,0.215) {$\alpha_{\max}$};
  \end{scope}
  \begin{scope}[shift={(3.35,-0.65)},xscale=0.6,yscale=7]
    \draw[ax] (0,0) -- (3.9,0) node[sm,right]{$v_x$};
    \draw[ax] (0,-0.09) -- (0,0.34) node[sm,above]{$F_x$};
    \draw[cD,line width=1pt,domain=0:3.5] plot (\x,{0.287-0.0545*\x}) node[sm,right]{$D{=}1$};
    \draw[cD!50,line width=1pt,domain=0:3.5] plot (\x,{0.5*(0.287-0.0545*\x)}) node[sm,right]{$D{=}\tfrac12$};
    \draw[cD,dashed,line width=0.8pt,domain=0:3.5] plot (\x,{-0.0518-0.00035*\x*\x}) node[sm,right]{$F_{\mathrm{fric}}$};
  \end{scope}
\end{tikzpicture}
\end{figside}
\begin{eqside}[0.42]
\[\begin{aligned}
\xD{F_{f,y}} &= \xP{D_f}\sin\!\big(\xP{C_f}\arctan(\xP{B_f}\,\xD{\alpha_f})\big)\\
\xD{F_{r,y}} &= \xP{D_r}\sin\!\big(\xP{C_r}\arctan(\xP{B_r}\,\xD{\alpha_r})\big)\\
\xD{F_{r,x}} &= (\xP{C_{m1}}-\xP{C_{m2}}\,\xS{v_x})\,\xS{D}\\
\xD{F_{\mathrm{fric}}} &= -\xP{C_{r0}}-\xP{C_{r2}}\,\xS{v_x}^2
\end{aligned}\]
{\footnotesize $\xP{D_f},\xP{D_r}$: peak force $\approx\mu F_z$ \quad $\xS D$: duty cycle $\in[-0.1,1]$}
\deps{$F_y\leftarrow\xD{\alpha}$;\quad $F_{r,x}\leftarrow\xS{D,v_x}$;\quad $F_{\mathrm{fric}}\leftarrow\xS{v_x}$}
\restr{lateral force saturates at $D$; drive force fades with speed; front tire has no $F_x$.}
\end{eqside}
\end{panel}

% ---------------------------------------------------------------- A4
\begin{panel}{A4}{Newton--Euler in the body frame}
\begin{figside}
\begin{tikzpicture}
  \coordinate (R) at (0,0); \coordinate (G) at (1.6,0); \coordinate (F) at (3.4,0);
  \draw[black!40,line width=1.4pt] (R) -- (F);
  \node[whl] at (R) {};
  \node[whl,rotate=22] at (F) {};
  \fill (G) circle (1.6pt);
  \draw[vD] (R) -- ++(0,1.0) node[lab,above]{$F_{r,y}$};
  \draw[vD] (-1.1,0) -- (-0.33,0);
  \node[lab,text=cD,below] at (-0.75,-0.05) {$F_{r,x}$};
  \draw[vD] (F) -- ++(112:1.05) node[lab,above]{$F_{f,y}$};
  \draw[vD,dashed] (G) -- ++(-0.75,0);
  \node[lab,text=cD,above] at (1.15,0.05) {$F_{\mathrm{fric}}$};
  \draw[cS,->,line width=0.7pt] ($(G)+(-150:0.5)$) arc (-150:-30:0.5);
  \node[lab,text=cS] at ($(G)+(-90:0.72)$) {$r$};
  \draw[dim] (0,-1.15) -- (1.6,-1.15) node[sm,midway,below]{$l_r$};
  \draw[dim] (1.6,-1.15) -- (3.4,-1.15) node[sm,midway,below]{$l_f$};
\end{tikzpicture}
\end{figside}
\begin{eqside}
\[\begin{aligned}
\xP m\,\xS{\dot v_x} &= \xD{F_{r,x}}+\xD{F_{\mathrm{fric}}}-\xD{F_{f,y}}\sin\xS\delta+\xP m\,\xS{v_y}\,\xS r\\
\xP m\,\xS{\dot v_y} &= \xD{F_{r,y}}+\xD{F_{f,y}}\cos\xS\delta-\xP m\,\xS{v_x}\,\xS r\\
\xP{I_z}\,\xS{\dot r} &= \xD{F_{f,y}}\,\xP{l_f}\cos\xS\delta-\xD{F_{r,y}}\,\xP{l_r}
\end{aligned}\]
{\footnotesize $\pm m\,v\,r$ terms: the body frame itself rotates.}
\deps{$(\dot v_x,\dot v_y,\dot r)\leftarrow\xD{F}$'s, $\xS{\delta}$, $\xS{v_x,v_y,r}$}
\restr{accelerations = forces / inertia. Linearized + ZOH-discretized each QP (A9).}
\end{eqside}
\end{panel}

% ---------------------------------------------------------------- A5
\begin{panel}{A5}{Progress $s$, contouring error $e_c$, lag error $e_l$}
\begin{figside}
\begin{tikzpicture}
  \begin{scope}[shift={(0,-4.4)}]
    \draw[refl] (127:5) arc (127:58:5) node[sm,right]{$\mathbf p_{\mathrm{ref}}(\cdot)$};
    \coordinate (P) at (102:5);
    \coordinate (S) at (88:5);
    \coordinate (p) at (88:4.15);
    \coordinate (N) at ($(P)+(-78:1)$);
    \coordinate (Q) at ($(P)!(p)!(N)$);
    \draw[guide] ($(P)+(192:0.6)$) -- ($(P)+(12:2.9)$);
    \draw[vec,black!70] (P) -- ++(12:0.65) node[sm,above]{$\mathbf t$};
    \draw[black!45,dotted,line width=0.8pt] (p) -- (S);
    \draw[cD,line width=1.2pt] (P) -- (Q) node[lab,midway,left=1pt]{$e_c$};
    \draw[cD,line width=1.2pt] (Q) -- (p) node[lab,midway,below=1pt]{$e_l$};
    \caricon{(p)}{-2}{0.65}
    \node[dot,cR] at (P) {};
    \node[sm,text=cR,above left] at (P) {$s$};
    \node[odot,draw=cR] at (S) {};
    \node[sm,text=cR,above] at (S) {$s^{\ast}$};
    \node[sm,below right=4pt] at (p) {$(X,Y)$};
  \end{scope}
\end{tikzpicture}
\end{figside}
\begin{eqside}
\[\begin{aligned}
\xD{e_c} &= \phantom{-}\xR{\mathbf n(\xS s)}^{\!\top}\big(\xS{\mathbf p}-\xR{\mathbf p_{\mathrm{ref}}(\xS s)}\big)\\
\xD{e_l} &= -\xR{\mathbf t(\xS s)}^{\!\top}\big(\xS{\mathbf p}-\xR{\mathbf p_{\mathrm{ref}}(\xS s)}\big)
\end{aligned}\]
{\footnotesize $\xR{\mathbf t}=(\cos\Phi,\sin\Phi)$,\ \ $\xR{\mathbf n}=(\sin\Phi,-\cos\Phi)$ (right),\ \ $\xS{\mathbf p}=(\xS X,\xS Y)$}
\deps{$e_c,e_l\leftarrow\xS{X,Y,s}$ and the spline $\xR{\mathbf p_{\mathrm{ref}}}$}
\restr{nothing alone; $e_l\!\approx\!0 \Leftrightarrow s\!\approx\!s^{\ast}$ (projection), with no inner $\arg\min$.}
\end{eqside}
\end{panel}

% ---------------------------------------------------------------- A6
\begin{panel}{A6}{Stage cost as springs}
\begin{figside}
\begin{tikzpicture}
  \draw[refl] (-0.2,0) -- (5.6,0);
  \coordinate (P) at (3.5,0); \coordinate (Fp) at (1.9,0); \coordinate (p) at (1.9,1.15);
  \draw[cD,line width=1.1pt,decorate,decoration={zigzag,segment length=3.2pt,amplitude=2.6pt}] (Fp) -- (P);
  \draw[cD!75,line width=0.5pt,decorate,decoration={coil,aspect=0.45,segment length=6pt,amplitude=3pt}] (Fp) -- (p);
  \caricon{(p)}{0}{0.8}
  \node[dot,cR] at (P) {};
  \node[sm,text=cR,below=2pt] at (P) {$s$};
  \draw[vS,line width=1.2pt] (P) -- ++(1.15,0);
  \node[lab,text=cS,above] at ($(P)+(0.62,0.05)$) {$q_{v_s}v_s$};
  \node[lab,text=cD,below] at ($(Fp)!0.5!(P)+(0,-0.12)$) {$q_l\,e_l^2$\ \ \textcolor{cP}{\scriptsize$q_l{=}500$}};
  \node[lab,text=cD,left=5pt] at ($(Fp)!0.5!(p)+(0,0.12)$) {$q_c\,e_c^2$};
  \node[sm,text=cP,left=5pt] at ($(Fp)!0.5!(p)+(0,-0.22)$) {$q_c{=}0.1$};
\end{tikzpicture}
\end{figside}
\begin{eqside}
\[J=\sum_{k=0}^{N}\Big(\xP{q_c}\xD{e_c}^2+\xP{q_l}\xD{e_l}^2-\xP{q_{v_s}}\xS{v_s}\Big)+\text{reg.}\]
{\footnotesize reg.\ $= \xP{q_\beta}\beta^2+\xP{q_r}\xS r^2+\xP{q_\mu}(\xS\varphi-\xR\Phi)^2+\|(\xS{D,\delta,v_s})\|^2_{\xP{R}}+\|\xU{\mathbf u}\|^2_{\xP{R_\Delta}}$\\
\phantom{reg.\ }$\beta=\arctan(\xS{v_y}/\xS{v_x})$;\ \ $q_c,q_r\times10$ at $k=N$;\ \ $q_{v_s}{=}0.02$}
\deps{$e_c,e_l$ (A5), $\xS{v_s}$, $\xS{v_x,v_y,r}$, $\xU{\mathbf u}$}
\restr{soft only: $s$ pinned to the car (stiff), lateral position nearly free (weak), progress rewarded.}
\end{eqside}
\end{panel}

% ---------------------------------------------------------------- A7
\begin{panel}{A7}{Track constraint (linearized slab)}
\begin{figside}
\begin{tikzpicture}
  \begin{scope}[shift={(0,-4.3)}]
    \coordinate (P) at (95:4.5);
    \coordinate (Ao) at (95:5.25);
    \coordinate (Ai) at (95:3.75);
    \path ($(Ao)+(185:2.6)$) coordinate (o1) ($(Ao)+(5:2.6)$) coordinate (o2)
          ($(Ai)+(185:2.6)$) coordinate (i1) ($(Ai)+(5:2.6)$) coordinate (i2);
    \fill[cS!7] (o1) -- (o2) -- (i2) -- (i1) -- cycle;
    \begin{scope}
      \clip (o1) -- (o2) -- (i2) -- (i1) -- cycle;
      \fill[cL!20,even odd rule] (0,0) circle (7) (0,0) circle (5.25);
    \end{scope}
    \draw[bnd] (128:5.25) arc (128:60:5.25);
    \draw[bnd] (128:3.75) arc (128:60:3.75);
    \draw[cR,dash pattern=on 4pt off 2pt,line width=0.8pt] (128:4.5) arc (128:60:4.5);
    \draw[liml] (o1) -- (o2);
    \draw[liml] (i1) -- (i2);
    \draw[vec,black!70] (P) -- ++(95:0.5) node[sm,right]{$\bar{\mathbf n}$};
    \node[dot,cR] at (P) {};
    \node[sm,text=cR,below right] at (P) {$\bar s_k$};
    \draw[dim] ($(P)+(185:1.7)$) -- ($(Ao)+(185:1.7)$) node[sm,midway,left]{$R_{\mathrm{out}}$};
    \draw[dim] ($(P)+(185:1.7)$) -- ($(Ai)+(185:1.7)$) node[sm,midway,left]{$R_{\mathrm{in}}$};
    \node[dot,cS] at ($(P)+(5:1.3)+(95:-0.42)$) {};
    \node[sm,text=cS,below] at ($(P)+(5:1.3)+(95:-0.42)$) {$\mathbf p_k$};
    \node[sm,text=cL] at ($(o2)+(-0.15,-0.35)$) {\tiny lin.\ error};
  \end{scope}
\end{tikzpicture}
\end{figside}
\begin{eqside}
\[-\xL{R_{\mathrm{in}}}\ \le\ \xR{\bar{\mathbf n}}^{\!\top}\big(\xS{\mathbf p_k}-\xR{\mathbf p_{\mathrm{ref}}(\bar s_k)}\big)\ \le\ \xL{R_{\mathrm{out}}}\quad(+\,\xi_k\ge 0)\]
{\footnotesize $\bar s_k$: progress of the \emph{previous guess} (fixed in this QP); $\xR{\bar{\mathbf n}}$: left normal there;\\
$R_{\mathrm{in}}{=}R_{\mathrm{out}}{=}0.14$\,m constant; car treated as a point.}
\deps{$\xS{X_k,Y_k}$ only}
\restr{car point inside a straight slab; soft (slack $\xi$, penalized); can overshoot the real edge far from $\bar s_k$.}
\end{eqside}
\end{panel}

% ---------------------------------------------------------------- A8
\begin{panel}{A8}{Tire constraints}
\begin{figside}
\begin{tikzpicture}[scale=0.72]
  \fill[cS!7] (0,0) ellipse (1.83 and 1.65);
  \draw[ax] (-2.3,0) -- (2.9,0) node[sm,right]{$F_{r,x}$};
  \draw[ax] (0,-2.0) -- (0,2.3) node[sm,above]{$F_{r,y}$};
  \draw[guide] (0,0) circle (1.737);
  \draw[liml] (0,0) ellipse (1.83 and 1.65);
  \draw[vD] (0,0) -- (0.9,1.2);
  \node[dot,cD] at (0.9,1.2) {};
  \node[sm,text=cD,right=2pt] at (0.9,1.2) {$(F_{r,x},F_{r,y})$};
  \node[sm,text=cL,anchor=north west] at (1.86,-0.3) {$e_\varepsilon D_r/e_{\mathrm{long}}$};
  \node[sm,text=cL,anchor=south east] at (-0.05,1.68) {$e_\varepsilon D_r$};
  \node[sm,text=black!55,anchor=north east] at (-1.25,-1.25) {$D_r$};
\end{tikzpicture}
\end{figside}
\begin{eqside}
\[\Big(\tfrac{\xP{e_{\mathrm{long}}}\,\xD{F_{r,x}}}{\xP{F_{N,r}}}\Big)^2+\Big(\tfrac{\xD{F_{r,y}}}{\xP{F_{N,r}}}\Big)^2\ \le\ \Big(\tfrac{\xL{e_\varepsilon}\,\xP{D_r}}{\xP{F_{N,r}}}\Big)^2\]
\[|\xD{\alpha_f}|\ \le\ \xL{\alpha_{\max}}\]
{\footnotesize $e_{\mathrm{long}}{=}0.9$,\ $e_\varepsilon{=}0.95$,\ $\alpha_{\max}{=}0.6$; $F_{N,r}$ static load. Both first-order in $\xS{\mathbf x}$ around the guess.}
\deps{ellipse $\leftarrow\xS{v_x,v_y,r,D}$;\quad $|\alpha_f|\leftarrow\xS{v_x,v_y,r,\delta}$}
\restr{rear: drive + lateral force share one ellipse at 95\,\% of peak. Front: slip below $\alpha_{\max}$ (A3, red band). Soft.}
\end{eqside}
\end{panel}

% ---------------------------------------------------------------- A9
\begin{panel}{A9}{One QP per SQP step, trust region on $s$}
\begin{figside}
\begin{tikzpicture}[x=0.52cm,y=1cm]
  \fill[cL!12] (0,0.0) -- (11,2.18) -- (11,2.68) -- (0,0.5) -- cycle;
  \draw[ax] (-0.3,0) -- (12.4,0) node[sm,right]{$k$};
  \draw[ax] (0,-0.1) -- (0,2.8) node[sm,left]{$s_k$};
  \draw[black!40,line width=0.5pt] (0,0.25) -- (11,2.43);
  \foreach \k in {1,...,7,11} {
    \pgfmathsetmacro\sb{0.25+0.198*\k}
    \pgfmathsetmacro\so{\sb+0.17*sin(40*\k)}
    \node[odot,draw=black!45] at (\k,\sb) {};
    \node[dot,cS] at (\k,\so) {};
  }
  \node[dot,cS] at (0,0.25) {};
  \node[sm,left] at (0,0.25) {$\hat{\mathbf x}$};
  \node[sm] at (9,1.4) {$\cdots$};
  \foreach \k/\l in {0/0,1/1,2/2,11/N} \node[sm,below] at (\k,-0.05) {$\l$};
  \draw[decorate,decoration={brace,amplitude=4pt},black!50] (0.2,3.05) -- (11,3.05) node[sm,midway,above=4pt]{$N\,T_s=60\times20\,\text{ms}$};
  \node[sm,text=cL] at (6.3,2.55) {$|s_k-\bar s_k|\le 0.2$\,m};
  \node[sm,text=black!50] at (8.6,0.95) {$\bar s_k$ (guess)};
\end{tikzpicture}
\end{figside}
\begin{eqside}
\[\begin{aligned}
\min\ &\textstyle\sum_k \tfrac12\xS{\mathbf x_k}^{\!\top}\!Q_k\xS{\mathbf x_k}+\mathbf q_k^{\!\top}\xS{\mathbf x_k}+\tfrac12\xU{\mathbf u_k}^{\!\top}\!R\,\xU{\mathbf u_k}+\text{slack}\\
\text{s.t. }&\xS{\mathbf x_0}=\hat{\mathbf x},\quad \xS{\mathbf x_{k+1}}=A_k\xS{\mathbf x_k}+B_k\xU{\mathbf u_k}+\mathbf g_k\\
&\text{A7, A8 (soft)},\ \ \text{box bounds on }\xS{\mathbf x},\xU{\mathbf u}\\
&|\xS{s_k}-\bar s_k|\le 0.2\,\text{m}
\end{aligned}\]
{\footnotesize $(A_k,B_k,\mathbf g_k)$: exact ZOH of the Jacobian at $(\bar{\mathbf x}_k,\bar{\mathbf u}_k)$.\\
$\bar{\mathbf x}\leftarrow 0.8\,\mathbf x^{\ast}+0.2\,\bar{\mathbf x}$, \ 2 iterations, apply $\xU{\mathbf u_0}$, shift.}
\restr{each QP trusts its linearization only near the previous guess.}
\end{eqside}
\end{panel}

% =====================================================================
\clearpage
\sect{Offline \quad TUM \texttt{global\_racetrajectory\_optimization} \quad (whole lap)}

% ---------------------------------------------------------------- B0
\begin{panel}{B0}{Two pipelines}
\centering
\begin{tikzpicture}[
  bx/.style={draw=black!35,rounded corners=2pt,inner sep=3.5pt,font=\footnotesize,align=center,fill=#1!9},
  fl/.style={->,black!55,line width=0.6pt}]
  \node[font=\sffamily\small\bfseries,anchor=west] at (-1.8,0.95) {mincurv};
  \node[font=\footnotesize\itshape,text=black!50,anchor=west] at (-0.2,0.95) {path first, speed second};
  \node[bx=cR] (t1) at (0,0) {$\xR{\mathbf p_i,\ \mathbf n_i}$\\ $\xR{w_{\ell,i},\ w_{r,i}}$};
  \node[bx=cS] (q1) at (3.4,0) {QP over $\xS{\boldsymbol\alpha}$\\ $\min\sum\kappa_i^2$};
  \node[bx=cD] (r1) at (6.6,0) {$\xD{\mathbf r(s)},\ \xD{\kappa(s)}$};
  \node[bx=cS] (v1) at (10.2,0) {fwd/bwd pass\\ $\xS{v(s)},\ \xD{a_x(s)}$};
  \node[bx=cP] (g1) at (10.2,-1.3) {$\text{ggv}(v),\ a_{x,\mathrm{mach}}(v),\ c_d,\ m,\ \mu$};
  \draw[fl] (t1) -- (q1); \draw[fl] (q1) -- (r1); \draw[fl] (r1) -- (v1); \draw[fl] (g1) -- (v1);
  \draw[fl,dashed] (r1.north) to[bend right=35] node[sm,above]{IQP: re-linearize} (q1.north);
  \node[font=\sffamily\small\bfseries,anchor=west] at (-1.8,-2.35) {mintime};
  \node[font=\footnotesize\itshape,text=black!50,anchor=west] at (-0.2,-2.35) {path and speed in one problem};
  \node[bx=cR] (t2) at (0,-3.4) {$\xR{\kappa_{\mathrm{ref}}(s)}$\\ $\xR{w_\ell(s),\ w_r(s)}$};
  \node[bx=cS] (n2) at (5.0,-3.4) {NLP (collocation + IPOPT)\\ $\min\int_0^L \tfrac{dt}{ds}\,ds$\\ over $\xS{\mathbf x(s)},\ \xU{\mathbf u(s)}$};
  \node[bx=cD] (o2) at (10.2,-3.4) {$\xD{\mathbf r(s)},\ \xS{v(s)}$};
  \node[bx=cP] (p2) at (5.0,-4.95) {double track, Pacejka, load transfer, aero, $\mu$};
  \draw[fl] (t2) -- (n2); \draw[fl] (n2) -- (o2); \draw[fl] (p2) -- (n2);
\end{tikzpicture}
\restr{\texttt{mincurv}: speed never feeds back into geometry. \texttt{mintime}: geometry and speed traded off jointly.}
\end{panel}

% ---------------------------------------------------------------- B1
\begin{panel}{B1}{mincurv: the only variable is a lateral shift $\alpha_i$}
\begin{figside}
\begin{tikzpicture}
  \begin{scope}[shift={(0,-5.65)}]
    \draw[bnd] (121:7) arc (121:59:7);
    \draw[bnd] (121:5) arc (121:59:5);
    \draw[refl,dash pattern=on 4pt off 2pt] (121:6) arc (121:59:6);
    \draw[cS,line width=1pt] plot[smooth] coordinates {(121:6.62) (115:6.6) (105:6.15) (95:5.45) (85:5.45) (75:6.15) (65:6.6) (59:6.62)};
    \foreach \t/\a in {115/-0.6,105/-0.15,95/0.55,85/0.55,75/-0.15,65/-0.6} {
      \draw[->,black!40,line width=0.5pt] (\t:6) -- (\t:5.5);
      \node[dot,cR] at (\t:6) {};
      \pgfmathsetmacro\rr{6-\a}
      \node[dot,cS] at (\t:\rr) {};
    }
    \draw[cL!60,line width=3pt,opacity=0.5] (77.5:6.9) -- (77.5:5.1);
    \draw[cL,line width=0.6pt] ($(77.5:6.9)+(167.5:0.12)$) -- ($(77.5:6.9)+(-12.5:0.12)$);
    \draw[cL,line width=0.6pt] ($(77.5:5.1)+(167.5:0.12)$) -- ($(77.5:5.1)+(-12.5:0.12)$);
    \node[sm,text=cL,right=3pt] at (77.5:5.25) {$[-\bar w_\ell,\ \bar w_r]$};
    \node[sm,text=cR,above] at (95:6.02) {$\mathbf p_i$};
    \node[sm,text=cS,below] at (95:5.42) {$\mathbf r_i$};
    \node[sm,text=black!55,left] at (115:5.6) {$\mathbf n_i$};
    \node[lab,text=cS,right] at (91.5:5.75) {$\alpha_i$};
    \draw[dim] (119:6) -- (119:7) node[sm,midway,left]{$w_\ell$};
    \draw[dim] (119:6) -- (119:5) node[sm,midway,left]{$w_r$};
    \draw[->,cR,line width=0.8pt] (63:6) arc (63:58:6);
  \end{scope}
\end{tikzpicture}
\end{figside}
\begin{eqside}
\[\xD{\mathbf r_i}=\xR{\mathbf p_i}+\xS{\alpha_i}\,\xR{\mathbf n_i}\]
\[-\big(\xR{w_{\ell,i}}-\tfrac{\xP{w_{\mathrm{veh}}}}{2}\big)\ \le\ \xS{\alpha_i}\ \le\ \xR{w_{r,i}}-\tfrac{\xP{w_{\mathrm{veh}}}}{2}\]
{\footnotesize $\xR{\mathbf n_i}$ points \emph{right};\quad $\bar w=w-w_{\mathrm{veh}}/2$;\quad decision vector $\xS{\boldsymbol\alpha}\in\mathbb R^{N}$.}
\deps{$\mathbf r_i\leftarrow\xS{\alpha_i}$ and fixed track data}
\restr{each point slides only along its own normal, inside the track minus half the car width.}
\end{eqside}
\end{panel}

% ---------------------------------------------------------------- B2
\begin{panel}{B2}{mincurv: curvature, linearized}
\begin{figside}
\begin{tikzpicture}[scale=1.05]
  \draw[cR,line width=0.8pt,dash pattern=on 4pt off 2pt,domain=-2.2:2.5,samples=50,smooth] plot (\x,{0.22*\x*\x-0.75});
  \draw[cS,line width=1pt,domain=-2.2:2.5,samples=50,smooth] plot (\x,{0.22*\x*\x});
  \coordinate (Ar) at (1.2,-0.4332);
  \coordinate (A) at (1.2,0.3168);
  \draw[vec,black!55] (Ar) -- ++(0.95,0.5016) node[sm,right]{$(\xP{x'},\xP{y'})$ frozen};
  \node[dot,cR] at (Ar) {};
  \draw[vec,cS] (A) -- ++(0,0.85) node[sm,above]{$(\xS{x''},\xS{y''})$};
  \draw[guide] (A) -- ++(-0.7004,1.3265) node[sm,left]{$1/\kappa_i$};
  \draw[cD,dashed,line width=0.6pt] (-0.335,3.224) ++(-98:3.287) arc (-98:-26:3.287);
  \node[dot,cS] at (A) {};
  \node[sm,text=cS,right=2pt] at (A) {$\mathbf r_i$};
  \node[sm,text=cR] at (-1.6,-0.05) {ref.};
  \node[sm,text=cS] at (-1.75,1.25) {raceline};
\end{tikzpicture}
\end{figside}
\begin{eqside}
\[\xD{\kappa_i}=\frac{\xP{x_i'}\,\xS{y_i''}-\xP{y_i'}\,\xS{x_i''}}{\big(\xP{x_i'}^2+\xP{y_i'}^2\big)^{3/2}}
\ \Rightarrow\ \xD{\boldsymbol\kappa}\approx\xR{\boldsymbol\kappa_{\mathrm{ref}}}+\xP{E_\kappa}\,\xS{\boldsymbol\alpha}\]
\[\min_{\xS{\boldsymbol\alpha}}\ \sum_i\xD{\kappa_i}^2=\tfrac12\xS{\boldsymbol\alpha}^{\!\top}\!H\xS{\boldsymbol\alpha}+\mathbf f^{\top}\xS{\boldsymbol\alpha}
\quad\text{s.t.}\ \ |\xD{\kappa_i}|\le\xL{\kappa_{\max}}\ \ \&\ \text{B1}\]
{\footnotesize $x',y'$ taken from the reference spline (frozen); $x'',y''$ are linear in $\boldsymbol\alpha$ through the spline system. IQP: refreeze at the new line, re-solve.}
\deps{$\kappa_i\leftarrow\xS{\boldsymbol\alpha}$ (neighbours too, via splines)}
\restr{a QP: smoothest line through the corridor; no speed, mass or tire in it.}
\end{eqside}
\end{panel}

% ---------------------------------------------------------------- B3
\begin{panel}{B3}{Velocity profile: acceleration left after cornering}
\begin{figside}
\begin{tikzpicture}
  \begin{scope}
    \clip (-2.0,-1.6) rectangle (2.0,0.8);
    \fill[cS!7] (0,0) ellipse (1.6 and 1.3);
  \end{scope}
  \draw[ax] (-2.1,0) -- (2.4,0) node[sm,right]{$a_y$};
  \draw[ax] (0,-1.7) -- (0,1.85) node[sm,above]{$a_x$};
  \draw[cL,line width=0.9pt] (0,0) ellipse (1.6 and 1.3);
  \draw[cL!70,dashed,line width=0.7pt] (1.6,0) -- (0,1.3) -- (-1.6,0) -- (0,-1.3) -- cycle;
  \draw[liml] (-1.9,0.8) -- (1.9,0.8) node[sm,right]{$a_{x,\mathrm{mach}}(v)$};
  \draw[cD,line width=0.8pt] (0.75,-1.6) -- (0.75,1.6) node[sm,above]{$a_y{=}v^2\kappa$};
  \node[odot,draw=cD] at (0.75,1.148) {};
  \node[dot,cD] at (0.75,-1.148) {};
  \node[dot,cD] at (0.75,0.8) {};
  \node[sm,text=cL,anchor=north west] at (1.62,-0.18) {$\mu a_{y,\max}$};
  \node[sm,text=cL,above left] at (0,1.3) {$\mu a_{x,\max}$};
  \node[sm,text=cL] at (-1.45,1.05) {$p{=}2$};
  \node[sm,text=cL!70] at (-0.55,-0.6) {$p{=}1$};
\end{tikzpicture}
\end{figside}
\begin{eqside}
\[\xD{a_y}=\xS v^2\,\xD{\kappa}\]
\[\xD{a_x^{\mathrm{avail}}}=\xL{\mu\,a_{x,\max}(v)}\Big[1-\Big(\frac{\xD{a_y}}{\xL{\mu\,a_{y,\max}(v)}}\Big)^{\xP p}\Big]^{1/\xP p}\]
\[\begin{aligned}
\text{accel: }&\min\!\big(\xD{a_x^{\mathrm{avail}}},\,\xL{a_{x,\mathrm{mach}}(v)}\big)-\tfrac{\xP{c_d}}{\xP m}\xS v^2\\
\text{brake: }&\xD{a_x^{\mathrm{avail}}}+\tfrac{\xP{c_d}}{\xP m}\xS v^2\quad(\text{magnitude})
\end{aligned}\]
{\footnotesize ggv table: $(v,a_{x,\max},a_{y,\max})$; $p$ = \texttt{dyn\_model\_exp} $\in[1,2]$ (default 1).}
\deps{$a_x\leftarrow\xS v,\ \xD\kappa$ (from B2)}
\restr{longitudinal accel.\ allowed by grip after cornering, and by the motor.}
\end{eqside}
\end{panel}

% ---------------------------------------------------------------- B4
\begin{panel}{B4}{Velocity profile: forward--backward pass}
\begin{figside}
\begin{tikzpicture}[x=0.72cm,y=0.36cm]
  \draw[ax] (0,0) -- (10.7,0) node[sm,right]{$s$};
  \draw[ax] (0,0) -- (0,7.2) node[sm,above]{$v$};
  \draw[cL,dash pattern=on 3pt off 2pt,line width=0.6pt,domain=0:10,samples=150,smooth]
     plot (\x,{6-4*exp(-(((\x-3)/0.6)^2))-3*exp(-(((\x-7.5)/0.6)^2))});
  \draw[cS!45,dashed,line width=0.7pt,domain=3:7.5,samples=40] plot (\x,{sqrt(4+5*(\x-3))});
  \draw[cS!45,dashed,line width=0.7pt,domain=7.5:10,samples=30] plot (\x,{sqrt(9+5*(\x-7.5))});
  \draw[cL!55,dashed,line width=0.7pt,domain=0.3:3,samples=30] plot (\x,{sqrt(4+8*(3-\x))});
  \draw[cL!55,dashed,line width=0.7pt,domain=4.2:7.5,samples=30] plot (\x,{sqrt(9+8*(7.5-\x))});
  \draw[cS,line width=1.3pt,domain=0:10,samples=200]
     plot (\x,{min(6-4*exp(-(((\x-3)/0.6)^2))-3*exp(-(((\x-7.5)/0.6)^2)),
                   (\x<3)*sqrt(4+8*abs(3-\x)) + (\x>=3)*(\x<7.5)*min(sqrt(4+5*abs(\x-3)),sqrt(9+8*abs(7.5-\x))) + (\x>=7.5)*sqrt(9+5*abs(\x-7.5)))});
  \node[sm,text=cL,anchor=south] at (5.2,6.05) {$v_\kappa=\sqrt{\mu a_{y,\max}/|\kappa|}$};
  \node[sm,text=cS!70] at (5.35,3.3) {fwd};
  \node[sm,text=cL!80] at (6.45,4.6) {bwd};
  \foreach \s in {3,7.5} \draw[black!40] (\s,-0.25) -- (\s,0.25);
  \node[sm,below] at (3,-0.2) {apex};
  \node[sm,below] at (7.5,-0.2) {apex};
\end{tikzpicture}
\end{figside}
\begin{eqside}
\[\text{fwd: }\ \xS{v_{i+1}}=\sqrt{\xS{v_i}^2+2\,\xD{a_x^{+}}(\xS{v_i},\xD{\kappa_i})\,\xR{\ell_i}}\]
\[\text{bwd: }\ \xS{v_{i}}=\sqrt{\xS{v_{i+1}}^2+2\,|\xD{a_x^{-}}(\xS{v_{i+1}},\xD{\kappa_{i+1}})|\,\xR{\ell_i}}\]
\[\xS{v}=\min\big(v_\kappa,\ \text{fwd},\ \text{bwd}\big)\]
{\footnotesize $\ell_i$: segment length; $a_x^{\pm}$ from B3. No optimizer: a sweep.}
\deps{$v\leftarrow\xD\kappa(s)$, ggv, motor, drag}
\restr{speed $\le$ cornering limit, and reachable by accelerating out of / braking into each apex.}
\end{eqside}
\end{panel}

% ---------------------------------------------------------------- B5
\begin{panel}{B5}{mintime: who depends on whom}
\begin{figside}[0.56]
\begin{tikzpicture}[x=0.55cm,y=0.37cm]
  \def\ncols{11}
  \mcol{1}{v}{cS}\mcol{2}{\beta}{cS}\mcol{3}{\omega_z}{cS}\mcol{4}{n}{cS}\mcol{5}{\xi}{cS}
  \mcol{6}{\delta}{cU}\mcol{7}{F_d}{cU}\mcol{8}{F_b}{cU}\mcol{9}{\gamma_y}{cU}
  \mcol{10}{\kappa}{cR}\mcol{11}{w}{cR}
  \mrow{1}{\xD{S_f}}{1,2,4,5,10}{0}
  \mrow{2}{\xD{\alpha_{ij}}}{1,2,3,6}{0}
  \mrow{3}{\xD{F_{z,ij}}}{1,7,8,9}{0}
  \mrow{4}{\xD{F_{x,ij}}}{7,8}{0}
  \mrow{5}{\xD{F_{y,ij}}}{0}{1,2,3,6,7,8,9}
  \mrow{6.4}{\xS{v'}}{2,6}{1,3,4,5,7,8,9,10}
  \mrow{7.4}{\xS{\beta'}}{1,2,3,6}{4,5,7,8,9,10}
  \mrow{8.4}{\xS{\omega_z'}}{6}{1,2,3,4,5,7,8,9,10}
  \mrow{9.4}{\xS{n'}}{1,2,5}{4,10}
  \mrow{10.4}{\xS{\xi'}}{3,10}{1,2,4,5}
  \mrow{11.8}{T_{\mathrm{lap}}}{0}{1,2,4,5,10}
  \mrow{12.8}{\xL{\text{track}}}{4,11}{0}
  \mrow{13.8}{\xL{\text{circles}}}{0}{1,2,3,6,7,8,9}
  \mrow{14.8}{\xL{\gamma_y\ \text{eq.}}}{9}{1,2,3,6,7,8}
  \mrow{15.8}{\xL{\text{power}}}{1,7}{0}
  \mrow{16.8}{\xL{F_dF_b}}{7,8}{0}
  \mrow{17.8}{\xL{\text{rates}}}{6,7,8}{1,4,10}
  \mrow{18.8}{\xL{\text{periodic}}}{1,2,3,4,5}{0}
  \draw[black!30] (0.5,-5.7) -- (11.5,-5.7);
  \draw[black!30] (0.5,-11.1) -- (11.5,-11.1);
  \draw[black!30] (5.5,1.4) -- (5.5,-19.3);
  \draw[black!30] (9.5,1.4) -- (9.5,-19.3);
\end{tikzpicture}
\end{figside}
\begin{eqside}[0.40]
{\footnotesize
$(\cdot)'=d(\cdot)/ds$\par\medskip
\tikz\fill[black!85] (0,0) circle (2.5pt);\ appears in the formula\par
\tikz\draw[black!60,line width=0.6pt] (0,0) circle (2.2pt);\ enters through $S_f$, $\alpha$ or $F$\par\medskip
Every $d/ds$ row inherits $n,\xi,\kappa$ through $S_f=dt/ds$.}
\restr{compared with A0: no $X,Y$ states and no progress variable; position is $(s,n)$ on the fixed reference.}
\end{eqside}
\end{panel}

% ---------------------------------------------------------------- B6
\begin{panel}{B6}{mintime: curvilinear coordinates, time as a function of $s$}
\begin{figside}
\begin{tikzpicture}
  \coordinate (O) at (0,0);
  \coordinate (P) at (0:4);
  \coordinate (C) at (0:2.85);
  \draw[refl] (-32:4) arc (-32:42:4) node[sm,above]{ref.};
  \node[dot,black!55] at (O) {};
  \draw[guide] (O) -- (-22:4) node[sm,pos=0.33,below,text=cR]{$1/\kappa(s)$};
  \node[dot,cR] at (P) {};
  \node[sm,text=cR,right] at (P) {$s$};
  \draw[guide] (2.85,-1.0) -- (2.85,2.4);
  \draw[guide] (4,-0.25) -- (4,-1.15);
  \draw[guide] (C) -- ++(108:2.3);
  \draw[dim] (2.85,-1.1) -- (4,-1.1) node[sm,midway,below,text=cS]{$n$};
  \begin{scope}[shift={(C)},rotate=108]
    \draw[body] (-0.72,-0.28) rectangle (0.72,0.28);
    \fill[black!55] (0.72,0) -- (0.52,0.12) -- (0.52,-0.12) -- cycle;
  \end{scope}
  \draw[vS] (C) -- ++(122:2.15) node[lab,above]{$v$};
  \draw[cS,line width=0.6pt] ($(C)+(90:1.3)$) arc (90:108:1.3);
  \node[lab,text=cS] at ($(C)+(99:1.08)$) {$\xi$};
  \draw[cS,line width=0.6pt] ($(C)+(108:1.75)$) arc (108:122:1.75);
  \node[lab,text=cS] at ($(C)+(115:2.0)$) {$\beta$};
  \draw[cS,->,line width=0.7pt] ($(C)+(195:0.95)$) arc (195:255:0.95);
  \node[lab,text=cS] at ($(C)+(222:1.22)$) {$\omega_z$};
\end{tikzpicture}
\end{figside}
\begin{eqside}
\[\xD{S_f}\equiv\frac{dt}{ds}=\frac{1-\xS n\,\xR\kappa}{\xS v\cos(\xS\xi+\xS\beta)}\]
\[\frac{d\xS n}{ds}=\xD{S_f}\,\xS v\sin(\xS\xi+\xS\beta),\qquad \frac{d\xS\xi}{ds}=\xD{S_f}\,\xS{\omega_z}-\xR\kappa\]
\[T_{\mathrm{lap}}=\int_0^{L}\xD{S_f}\,ds\]
{\footnotesize $\xS{\mathbf x}=(v,\beta,\omega_z,n,\xi)$,\ $\xU{\mathbf u}=(\delta,F_d,F_b,\gamma_y)$;\ \ $n>0$ = \emph{left} (opposite sign to $\alpha$ in B1).}
\deps{$S_f\leftarrow\xS{v,\beta,n,\xi},\ \xR\kappa$}
\restr{time is not a variable: lap time integrates $dt/ds$ over the fixed reference length $L$.}
\end{eqside}
\end{panel}

% ---------------------------------------------------------------- B7
\begin{panel}{B7}{mintime: double-track loads}
\begin{figside}
\begin{tikzpicture}
  \draw[body] (-0.55,-0.55) rectangle (3.85,0.55);
  \draw[black!35] (0,-0.8) -- (0,0.8);
  \draw[black!35] (3.3,-0.8) -- (3.3,0.8);
  \draw[cL!80,line width=0.6pt,fill=cL!6] (3.3,-0.8) circle (0.55);
  \draw[cL!80,line width=0.6pt,fill=cL!6] (3.3,0.8) circle (0.38);
  \draw[cL!80,line width=0.6pt,fill=cL!6] (0,-0.8) circle (0.40);
  \draw[cL!80,line width=0.6pt,fill=cL!6] (0,0.8) circle (0.24);
  \node[whl] at (0,0.8) {}; \node[whl] at (0,-0.8) {};
  \node[whl,rotate=15] at (3.3,0.8) {}; \node[whl,rotate=15] at (3.3,-0.8) {};
  \fill (1.54,0) circle (1.6pt);
  \draw[vD] (1.54,0) -- ++(-0.8,0) node[sm,left]{$a_x$};
  \draw[vD] (1.54,0) -- ++(0,0.5) node[sm,above]{$a_y$};
  \node[sm,text=cL] at (4.45,-0.8) {$\propto F_{z,fr}$};
  \draw[dim] (0,-1.55) -- (1.54,-1.55) node[sm,midway,below]{$l_r$};
  \draw[dim] (1.54,-1.55) -- (3.3,-1.55) node[sm,midway,below]{$l_f$};
  \draw[dim] (-1.0,-0.8) -- (-1.0,0.8) node[sm,midway,left]{$b$};
  \node[sm,text=black!50] at (1.65,1.15) {braking + left turn};
\end{tikzpicture}
\end{figside}
\begin{eqside}[0.46]
\[\xD{F_{z,ij}}=\xP{F^{\mathrm{st}}_{z,ij}}+\tfrac12\xP{c_{l,i}}\xS v^2
\mp\tfrac{\xP h}{2\xP l}\big(\xU{F_d}{+}\xU{F_b}{-}\xP{c_d}\xS v^2{-}\xP{F_{\mathrm{roll}}}\big)
\mp\xP{k_i}\,\xU{\gamma_y}\]
{\footnotesize long.: $-$front/$+$rear;\ lat.: $-$left/$+$right;\ $k_f{=}k_{\mathrm{roll}}$, $k_r{=}1{-}k_{\mathrm{roll}}$}
\[\xU{\gamma_y}=\frac{\xP h}{\xP{\bar b}}\,\xP m\,\xD{a_y}\quad(\text{equality constraint})\]
{\footnotesize $\xD{F_{x,ij}}$: drive/brake split $\xP{k_{\mathrm{drive}}},\xP{k_{\mathrm{brake}}}$ minus rolling resistance.}
\deps{$F_z\leftarrow\xS v,\ \xU{F_d,F_b,\gamma_y}$}
\restr{each wheel's load (hence grip) follows the accelerations; $\gamma_y$ is a decision variable closed by an equality.}
\end{eqside}
\end{panel}

% ---------------------------------------------------------------- B8
\begin{panel}{B8}{mintime: tire model and friction circles}
\begin{figside}
\begin{tikzpicture}
  \begin{scope}[xscale=7,yscale=1.6]
    \draw[ax] (0,0) -- (0.38,0) node[sm,right]{$\alpha$};
    \draw[ax] (0,0) -- (0,1.3) node[sm,above]{$F_y/D$};
    \draw[guide] (0,1) -- (0.36,1);
    \draw[cD,line width=1pt,domain=0:0.35,samples=150,smooth] plot (\x,{sin(2.5*atan(10*\x-1.0*(10*\x-rad(atan(10*\x)))))});
    \node[sm,left] at (0,1) {$1$};
  \end{scope}
  \begin{scope}[shift={(3.55,0)},xscale=0.25,yscale=0.25]
    \draw[ax] (0,0) -- (10,0) node[sm,right]{$F_z$};
    \draw[ax] (0,0) -- (0,9.6) node[sm,above]{$D$};
    \draw[guide] (0,0) -- (9,9) node[sm,above]{$\mu F_z$};
    \draw[cD,line width=1pt,domain=0:9,samples=30] plot (\x,{\x*(1-0.1*\x/3)});
    \node[sm,text=cD] at (8.2,4.9) {$\varepsilon<0$};
  \end{scope}
\end{tikzpicture}
\end{figside}
\begin{eqside}
\[\xD{D_{ij}}=\xP\mu\,\xD{F_{z,ij}}\Big(1+\xP\varepsilon\,\frac{\xD{F_{z,ij}}}{\xP{F_{z0}}}\Big)\]
\[\xD{F_{y,ij}}=\xD{D_{ij}}\sin\!\Big(\xP C\arctan\big(\xP B\xD\alpha-\xP E(\xP B\xD\alpha-\arctan\xP B\xD\alpha)\big)\Big)\]
\[\Big(\frac{\xD{F_{x,ij}}}{\xP\mu\,\xD{F_{z,ij}}}\Big)^2+\Big(\frac{\xD{F_{y,ij}}}{\xP\mu\,\xD{F_{z,ij}}}\Big)^2\le\xL{1}\qquad\forall\,ij\]
\deps{$\alpha_{ij}\leftarrow\xS{v,\beta,\omega_z},\xU\delta$;\ \ $F_{z,ij}$ from B7}
\restr{each tire inside its own circle sized by its load (B7 circles); load transfer costs total grip ($\varepsilon<0$).}
\end{eqside}
\end{panel}

% ---------------------------------------------------------------- B9
\begin{panel}{B9}{mintime: the NLP}
\begin{figside}[0.40]
\begin{tikzpicture}[x=1cm,y=1cm]
  \draw[ax] (-0.2,0) -- (5.3,0) node[sm,right]{$s$};
  \foreach \t in {0.1127,0.5,0.8873} \draw[black!25] ({0.4+4.2*\t},0) -- ({0.4+4.2*\t},2.2);
  \draw[cS,line width=1pt,domain=0:1,samples=40] plot ({0.4+4.2*\x},{1+2.4*\x-1.4*\x*\x});
  \foreach \t/\m in {0.1127/0.4963,0.5/0.2381,0.8873/-0.0201} {
    \pgfmathsetmacro\xx{0.4+4.2*\t}
    \pgfmathsetmacro\yy{1+2.4*\t-1.4*\t*\t}
    \draw[cD,line width=1.1pt] ({\xx-0.6},{\yy-0.6*\m+0.12}) -- ({\xx+0.6},{\yy+0.6*\m+0.12});
    \node[dot,cS] at (\xx,\yy) {};
  }
  \node[dot,cS] at (0.4,1) {};
  \node[odot,draw=cS] at (4.6,2) {};
  \node[sm,text=cS,left] at (0.4,1) {$X_k$};
  \node[sm,text=cS,right] at (4.6,2) {$X_{k+1}$};
  \node[sm,text=cD,above] at (2.5,2.3) {slope $=h_k\mathbf f(X_{k,j},U_k,\kappa)$};
  \draw[cU,line width=1pt] (-0.1,-0.75) -- (0.4,-0.75) (0.4,-0.5) -- (4.6,-0.5) (4.6,-0.85) -- (5.1,-0.85);
  \node[sm,text=cU,above] at (2.5,-0.5) {$U_k$};
  \draw[dim] (0.4,-1.15) -- (4.6,-1.15) node[sm,midway,below]{$h_k$};
\end{tikzpicture}
\end{figside}
\begin{eqside}[0.58]
\[\min_{\xS{\mathbf x},\xU{\mathbf u}}\ \int_0^{L}\xD{S_f}\,ds+\xP{\lambda_F}\|\Delta\xU F\|^2+\xP{\lambda_\delta}\|\Delta\xU\delta\|^2\]
\footnotesize
\[\begin{aligned}
&\text{dynamics:}  && h_k\,\mathbf f(X_{k,j},\xU{U_k},\xR{\kappa})=\textstyle\sum_r C_{rj}X_{k,r},\ \ 3\text{ Legendre pts}\\
&\text{track:}     && -\xR{w_r}+\tfrac{\xP{w_{\mathrm{opt}}}}{2}\le\xS n\le\xR{w_\ell}-\tfrac{\xP{w_{\mathrm{opt}}}}{2}\\
&\text{grip, loads:}&& \text{B8 circles},\ \ \gamma_y\ \text{equality (B7)}\\
&\text{power:}     && \xS v\,\xU{F_d}\le\xL{P_{\max}},\qquad -\epsilon\le\xU{F_d}\,\xU{F_b}\le 0\\
&\text{actuators:} && \tfrac{\Delta\xU\delta}{\Delta t}\in\pm\tfrac{\xL{\delta_{\max}}}{\xP{t_\delta}},\ \
                      \tfrac{\Delta\xU{F_d}}{\Delta t}\le\tfrac{\xL{F_{d,\max}}}{\xP{t_d}},\ \
                      \tfrac{\Delta\xU{F_b}}{\Delta t}\ge-\tfrac{\xL{F_{b,\max}}}{\xP{t_b}}\\
&\text{closed lap:}&& \xS{\mathbf x}(0)=\xS{\mathbf x}(L)\\
&\text{bounds:}    && \xS v\ge 1,\ \ |\xS\beta|,|\xS{\omega_z}|,|\xS\xi|\le\tfrac{\pi}{2},\ \ \xU{F_d}\ge0\ge\xU{F_b}
\end{aligned}\]
\small
\restr{dynamics enforced only at collocation points; the lap must close; everything else is a path constraint.}
\end{eqside}
\end{panel}
}
. Requires amsmath, amssymb, array.
% =====================================================================
\usepackage{xcolor}
\usepackage{tikz}
\usetikzlibrary{arrows.meta,calc,decorations.pathmorphing,decorations.pathreplacing}

% ---------------- role colours ----------------------------------------
\definecolor{cS}{HTML}{0072B2}  % decision state
\definecolor{cU}{HTML}{C76E00}  % control input
\definecolor{cR}{HTML}{00876C}  % reference / track data
\definecolor{cP}{HTML}{7A7A7A}  % fixed parameter
\definecolor{cD}{HTML}{8E44AD}  % derived quantity
\definecolor{cL}{HTML}{C0392B}  % limit / bound
\newcommand{\xS}[1]{{\color{cS}#1}}
\newcommand{\xU}[1]{{\color{cU}#1}}
\newcommand{\xR}[1]{{\color{cR}#1}}
\newcommand{\xP}[1]{{\color{cP}#1}}
\newcommand{\xD}[1]{{\color{cD}#1}}
\newcommand{\xL}[1]{{\color{cL}#1}}

% ---------------- layout ----------------------------------------------
\newenvironment{panel}[2]{%
  \par\addvspace{5pt}\noindent\begin{minipage}{\linewidth}%
  {\sffamily\small\bfseries\color{black!70}#1}\enspace{\sffamily\small #2}\par\vspace{4pt}}%
  {\end{minipage}\par\vspace{6pt}\noindent{\color{black!12}\rule{\linewidth}{0.5pt}}\par}
\newenvironment{figside}[1][0.53]{\begin{minipage}[c]{#1\linewidth}\centering}{\end{minipage}\hfill}
\newenvironment{eqside}[1][0.45]{\begin{minipage}[c]{#1\linewidth}\small
  \setlength{\abovedisplayskip}{3pt}\setlength{\belowdisplayskip}{3pt}%
  \setlength{\abovedisplayshortskip}{2pt}\setlength{\belowdisplayshortskip}{2pt}}{\end{minipage}}
\newcommand{\deps}[1]{\par\smallskip{\sffamily\scriptsize\bfseries\color{black!45}DEPENDS}\enspace{\small#1}}
\newcommand{\restr}[1]{\par\smallskip{\sffamily\scriptsize\bfseries\color{cL}RESTRAINS}\enspace{\small#1}}
\newcommand{\chip}[2]{\tikz[baseline=-0.6ex]\fill[#1] (0,0) circle (3pt);\,{\small\sffamily\color{#1}#2}\quad}
\providecommand{\sect}[1]{\par\vspace{8pt}{\large\sffamily\bfseries #1}\par\vspace{2pt}\noindent{\color{black!40}\rule{\linewidth}{0.8pt}}\par}

% ---------------- tikz styles -----------------------------------------
\tikzset{
  every picture/.style={line cap=round,line join=round,>={Stealth[length=4.5pt,width=3.2pt]}},
  ax/.style={->,black!45,line width=0.5pt},
  body/.style={draw=black!65,fill=black!5,line width=0.6pt,rounded corners=1.2pt},
  whl/.style={fill=black!75,draw=none,rounded corners=0.6pt,minimum width=5.5mm,minimum height=2mm,inner sep=0pt},
  vS/.style={->,cS,line width=0.9pt},
  vD/.style={->,cD,line width=0.9pt},
  vec/.style={->,line width=0.8pt},
  refl/.style={cR,line width=1pt},
  bnd/.style={black!55,line width=0.7pt},
  liml/.style={cL,line width=0.8pt},
  guide/.style={black!35,dash pattern=on 2pt off 1.5pt,line width=0.4pt},
  dim/.style={<->,black!50,line width=0.4pt,>={Stealth[length=3pt,width=2.4pt]}},
  lab/.style={font=\footnotesize,inner sep=1pt},
  sm/.style={font=\scriptsize,inner sep=1pt},
  dot/.style={circle,fill,inner sep=0pt,minimum size=3.4pt},
  odot/.style={circle,draw,fill=white,inner sep=0pt,minimum size=3.8pt,line width=0.6pt},
}
% small car icon: \caricon{(x,y)}{heading deg}{scale}
\newcommand{\caricon}[3]{%
  \begin{scope}[shift={#1},rotate=#2,scale=#3]
    \draw[body] (-0.45,-0.2) rectangle (0.45,0.2);
    \fill[black!55] (0.45,0) -- (0.28,0.11) -- (0.28,-0.11) -- cycle;
  \end{scope}}
% dependency-matrix row: {row}{label}{direct cols}{via cols}  (0 = none)
\newcommand{\mrow}[4]{%
  \draw[black!7,line width=3.6pt] (0.6,-#1) -- (\ncols+0.4,-#1);
  \node[anchor=east,font=\footnotesize,inner sep=1pt] at (0.45,-#1) {$#2$};
  \foreach \c in {#3}{\ifnum\c>0 \fill[black!85] (\c,-#1) circle (2.5pt);\fi}
  \foreach \c in {#4}{\ifnum\c>0 \draw[black!60,line width=0.6pt,fill=white] (\c,-#1) circle (2.2pt);\fi}}
\newcommand{\mcol}[3]{\node[font=\footnotesize,text=#3] at (#1,0.85) {$#2$};}

% =====================================================================

%  to change how they look (e.g. as \subsection).
% =====================================================================
{\setlength{\parindent}{0pt}%
\chip{cS}{decision state}\chip{cU}{control input}\chip{cR}{reference / track data}%
\chip{cP}{fixed parameter}\chip{cD}{derived}\chip{cL}{limit / bound}

% =====================================================================
\sect{Online \quad MPCC \quad (receding horizon, 60 steps $\times$ 20 ms)}

% ---------------------------------------------------------------- A0
\begin{panel}{A0}{Variables and who depends on whom}
\begin{figside}[0.60]
\begin{tikzpicture}[x=0.5cm,y=0.37cm]
  \def\ncols{14}
  \mcol{1}{X}{cS}\mcol{2}{Y}{cS}\mcol{3}{\varphi}{cS}\mcol{4}{v_x}{cS}\mcol{5}{v_y}{cS}
  \mcol{6}{r}{cS}\mcol{7}{s}{cS}\mcol{8}{D}{cS}\mcol{9}{\delta}{cS}\mcol{10}{v_s}{cS}
  \mcol{11}{\dot D}{cU}\mcol{12}{\dot\delta}{cU}\mcol{13}{\dot v_s}{cU}\mcol{14}{\mathbf p_{\!\mathrm{ref}}}{cR}
  \mrow{1}{\xD{\alpha_f}}{4,5,6,9}{0}
  \mrow{2}{\xD{\alpha_r}}{4,5,6}{0}
  \mrow{3}{\xD{F_{f,y}}}{0}{4,5,6,9}
  \mrow{4}{\xD{F_{r,y}}}{0}{4,5,6}
  \mrow{5}{\xD{F_{r,x}}}{4,8}{0}
  \mrow{6}{\xD{F_{\mathrm{fric}}}}{4}{0}
  \mrow{7.4}{\xS{\dot X,\dot Y}}{3,4,5}{0}
  \mrow{8.4}{\xS{\dot\varphi}}{6}{0}
  \mrow{9.4}{\xS{\dot v_x}}{5,6,9}{4,8}
  \mrow{10.4}{\xS{\dot v_y}}{4,6,9}{5}
  \mrow{11.4}{\xS{\dot r}}{9}{4,5,6}
  \mrow{12.4}{\xS{\dot s}}{10}{0}
  \mrow{13.4}{\xS{\dot D,\dot\delta,\dot v_s}}{11,12,13}{0}
  \mrow{14.8}{\xD{e_c,\,e_l}}{1,2,7,14}{0}
  \mrow{15.8}{\xL{\text{track}}}{1,2,14}{0}
  \mrow{16.8}{\xL{\text{ellipse}}}{0}{4,5,6,8}
  \mrow{17.8}{\xL{|\alpha_f|}}{0}{4,5,6,9}
  \mrow{18.8}{\beta\ \text{cost}}{4,5}{0}
  \mrow{19.8}{\varphi\ \text{cost}}{3,7,14}{0}
  \mrow{20.8}{\text{progress}}{10}{0}
  \draw[black!30] (0.5,-6.7) -- (14.5,-6.7);
  \draw[black!30] (0.5,-14.1) -- (14.5,-14.1);
  \draw[black!30] (10.5,1.4) -- (10.5,-21.3);
\end{tikzpicture}
\end{figside}
\begin{eqside}[0.38]
\[\xS{\mathbf x} = [\,\xS{X\ Y\ \varphi\ v_x\ v_y\ r\ s\ D\ \delta\ v_s}\,]^{\top}\]
\[\xU{\mathbf u} = [\,\xU{\dot D\ \ \dot\delta\ \ \dot v_s}\,]^{\top}\]
\[\xS{\dot s}=\xS{v_s},\quad \xS{\dot D}=\xU{u_1},\quad \xS{\dot\delta}=\xU{u_2},\quad \xS{\dot v_s}=\xU{u_3}\]
\medskip
{\footnotesize
\tikz\fill[black!85] (0,0) circle (2.5pt);\ appears in the formula\par
\tikz\draw[black!60,line width=0.6pt] (0,0) circle (2.2pt);\ enters through $\alpha$ or $F$\par}
\restr{inputs are \emph{rates}, so rate costs and slew bounds act directly on $\xU{\mathbf u}$.}
\end{eqside}

\par\medskip
{\footnotesize\renewcommand{\arraystretch}{1.12}%
\begin{minipage}[t]{0.335\linewidth}
{\sffamily\scriptsize\bfseries\color{cS}STATE $\mathbf x$}\par\smallskip
\begin{tabular}{@{}l@{\ \ }>{\raggedright\arraybackslash}p{4.35cm}@{}}
$\xS{X,Y}$       & car position, world frame [m]\\
$\xS{\varphi}$   & heading (yaw angle), world frame [rad]\\
$\xS{v_x}$       & forward velocity, body frame [m/s]\\
$\xS{v_y}$       & sideways velocity, body frame [m/s]\\
$\xS{r}$         & yaw rate [rad/s]\\
$\xS{s}$         & progress: virtual arc length along the reference [m]\\
$\xS{D}$         & motor duty cycle, $[-0.1,1]$; $<0$ brakes\\
$\xS{\delta}$    & front steering angle, $|\delta|\le 0.35$ rad\\
$\xS{v_s}$       & progress speed $\dot s\ge 0$ [m/s]
\end{tabular}
\end{minipage}\hfill
\begin{minipage}[t]{0.32\linewidth}
{\sffamily\scriptsize\bfseries\color{cU}INPUT $\mathbf u$}\par\smallskip
\begin{tabular}{@{}l@{\ \ }>{\raggedright\arraybackslash}p{4.05cm}@{}}
$\xU{\dot D}$         & duty-cycle rate\\
$\xU{\dot\delta}$     & steering rate\\
$\xU{\dot v_s}$       & progress acceleration
\end{tabular}\par\medskip
{\sffamily\scriptsize\bfseries\color{cR}REFERENCE}\par\smallskip
\begin{tabular}{@{}l@{\ \ }>{\raggedright\arraybackslash}p{4.05cm}@{}}
$\xR{\mathbf p_{\mathrm{ref}}}$ & centreline spline $(X_{\mathrm{ref}}(s),Y_{\mathrm{ref}}(s))$ with heading $\Phi(s)$; fixed data
\end{tabular}\par\medskip
{\sffamily\scriptsize\bfseries\color{cD}DERIVED}\par\smallskip
\begin{tabular}{@{}l@{\ \ }>{\raggedright\arraybackslash}p{4.05cm}@{}}
$\xD{\alpha_f,\alpha_r}$       & front / rear tire slip angle\\
$\xD{F_{f,y},F_{r,y}}$         & front / rear lateral tire force\\
$\xD{F_{r,x}}$                 & rear drive force from the motor\\
$\xD{F_{\mathrm{fric}}}$       & rolling resistance + air drag
\end{tabular}
\end{minipage}\hfill
\begin{minipage}[t]{0.31\linewidth}
{\sffamily\scriptsize\bfseries\color{black!60}COST / CONSTRAINT}\par\smallskip
\begin{tabular}{@{}l@{\ \ }>{\raggedright\arraybackslash}p{3.55cm}@{}}
$\dot{(\,\cdot\,)}$        & time derivative: the state equations\\
$\xD{e_c}$                 & contouring error: sideways offset from $\mathbf p_{\mathrm{ref}}(s)$\\
$\xD{e_l}$                 & lag error: along-track gap between $s$ and the car\\
\xL{track}                 & car inside the slab around the reference\\
\xL{ellipse}               & rear friction ellipse\\
$\xL{|\alpha_f|}$          & front slip-angle limit\\
$\beta$ cost               & sideslip $\arctan(v_y/v_x)$ penalty\\
$\varphi$ cost             & heading vs.\ track heading $\Phi(s)$\\
progress                   & reward $-q_{v_s}v_s$
\end{tabular}
\end{minipage}}
\end{panel}

% ---------------------------------------------------------------- A1
\begin{panel}{A1}{Body velocity $\to$ world pose}
\begin{figside}
\begin{tikzpicture}
  \draw[ax] (0,0) -- (1.0,0) node[sm,right]{$X$};
  \draw[ax] (0,0) -- (0,1.0) node[sm,above]{$Y$};
  \coordinate (C) at (2.9,0.45);
  % world components of v (dotted)
  \draw[black!55,dotted,line width=0.9pt] (C) -- ++(1.282,0) node[sm,below,midway,yshift=-2pt]{$\dot X$}
        -- ++(0,1.779) node[sm,right,pos=0.8]{$\dot Y$};
  % heading angle, measured behind the car
  \draw[guide] (C) -- ++(180:1.5);
  \draw[guide] (C) -- ++(210:1.5);
  \draw[cS,line width=0.6pt] ($(C)+(180:1.15)$) arc (180:210:1.15);
  \node[lab,text=cS] at ($(C)+(195:1.42)$) {$\varphi$};
  \begin{scope}[shift={(C)},rotate=30]
    \draw[body] (-0.55,-0.22) rectangle (0.55,0.22);
    \fill[black!55] (0.55,0) -- (0.4,0.1) -- (0.4,-0.1) -- cycle;
    \draw[vS] (0,0) -- (2.0,0) node[lab,below right=0pt]{$v_x$};
    \draw[vS] (0,0) -- (0,0.9) node[lab,above left=0pt]{$v_y$};
    \draw[->,black!75,line width=0.6pt] (0,0) -- (2.0,0.9) node[lab,above=1pt]{$\mathbf v$};
  \end{scope}
  \draw[cS,->,line width=0.7pt] ($(C)+(245:0.7)$) arc (245:305:0.7);
  \node[lab,text=cS] at ($(C)+(275:0.92)$) {$r$};
\end{tikzpicture}
\end{figside}
\begin{eqside}
\[\begin{aligned}
\xS{\dot X} &= \xS{v_x}\cos\xS{\varphi} - \xS{v_y}\sin\xS{\varphi}\\
\xS{\dot Y} &= \xS{v_x}\sin\xS{\varphi} + \xS{v_y}\cos\xS{\varphi}\\
\xS{\dot\varphi} &= \xS{r}
\end{aligned}\]
\deps{$(\dot X,\dot Y)\leftarrow \xS{v_x,v_y,\varphi}$;\quad $\dot\varphi\leftarrow\xS{r}$}
\restr{pose = integral of the body velocity rotated by the heading.}
\end{eqside}
\end{panel}

% ---------------------------------------------------------------- A2
\begin{panel}{A2}{Slip angles}
\begin{figside}
\begin{tikzpicture}
  \coordinate (R) at (0,0); \coordinate (G) at (1.5,0); \coordinate (F) at (3.2,0);
  \draw[black!40,line width=1.4pt] (R) -- (F);
  \draw[guide] (F) -- ++(0:1.7);
  \draw[guide] ($(F)+(208:0.9)$) -- ++(28:2.7);
  \node[whl] at (R) {};
  \node[whl,rotate=28] at (F) {};
  \fill (G) circle (1.6pt);
  % rear
  \draw[vS] (R) -- ++(-16:1.55) node[sm,below=2pt]{$(v_x,\ v_y{-}l_r r)$};
  \draw[cD,line width=0.8pt] ($(R)+(-16:1.1)$) arc (-16:0:1.1);
  \node[lab,text=cD] at ($(R)+(-6:1.35)$) {$\alpha_r$};
  % front
  \draw[vS] (F) -- ++(10:1.6) node[sm,right]{$(v_x,\ v_y{+}l_f r)$};
  \draw[cD,line width=0.8pt] ($(F)+(10:1.25)$) arc (10:28:1.25);
  \node[lab,text=cD] at ($(F)+(19:1.5)$) {$\alpha_f$};
  \draw[cS,line width=0.6pt] ($(F)+(180:0.75)$) arc (180:208:0.75);
  \node[lab,text=cS] at ($(F)+(194:1.0)$) {$\delta$};
  % dims
  \draw[dim] (0,-1.05) -- (1.5,-1.05) node[sm,midway,below]{$l_r$};
  \draw[dim] (1.5,-1.05) -- (3.2,-1.05) node[sm,midway,below]{$l_f$};
\end{tikzpicture}
\end{figside}
\begin{eqside}
\[\begin{aligned}
\xD{\alpha_f} &= \xS{\delta} - \arctan\frac{\xS{v_y}+\xP{l_f}\,\xS{r}}{\xS{v_x}}\\[2pt]
\xD{\alpha_r} &= -\arctan\frac{\xS{v_y}-\xP{l_r}\,\xS{r}}{\xS{v_x}}
\end{aligned}\]
\deps{$\alpha_f\leftarrow\xS{v_x,v_y,r,\delta}$;\quad $\alpha_r\leftarrow\xS{v_x,v_y,r}$}
\restr{nothing alone: angle between where a tire points and where it moves (bounded in A8).}
\end{eqside}
\end{panel}

% ---------------------------------------------------------------- A3
\begin{panel}{A3}{Tire and drivetrain forces}
\begin{figside}[0.56]
\begin{tikzpicture}
  \begin{scope}[xscale=1.35,yscale=7]
    \fill[cL!10] (0.6,-0.215) rectangle (1.65,0.215);
    \fill[cL!10] (-1.65,-0.215) rectangle (-0.6,0.215);
    \draw[ax] (-1.7,0) -- (1.8,0) node[sm,right]{$\alpha$};
    \draw[ax] (0,-0.22) -- (0,0.25) node[sm,above]{$F_y$};
    \draw[guide] (-1.65,0.192) -- (1.65,0.192);
    \node[sm,text=cP,anchor=south east] at (-0.05,0.193) {$D_f$};
    \draw[black!50,line width=0.5pt] (-0.15,-0.0891) -- (0.26,0.1545);
    \node[sm,text=cP,anchor=east] at (-0.05,0.12) {$B_fC_fD_f$};
    \draw[cD,line width=1pt,domain=-1.6:1.6,samples=100,smooth] plot (\x,{0.192*sin(1.2*atan(2.579*\x))});
    \draw[liml] (0.6,-0.215) -- (0.6,0.215);
    \draw[liml] (-0.6,-0.215) -- (-0.6,0.215);
    \node[sm,text=cL,above] at (0.6,0.215) {$\alpha_{\max}$};
  \end{scope}
  \begin{scope}[shift={(3.35,-0.65)},xscale=0.6,yscale=7]
    \draw[ax] (0,0) -- (3.9,0) node[sm,right]{$v_x$};
    \draw[ax] (0,-0.09) -- (0,0.34) node[sm,above]{$F_x$};
    \draw[cD,line width=1pt,domain=0:3.5] plot (\x,{0.287-0.0545*\x}) node[sm,right]{$D{=}1$};
    \draw[cD!50,line width=1pt,domain=0:3.5] plot (\x,{0.5*(0.287-0.0545*\x)}) node[sm,right]{$D{=}\tfrac12$};
    \draw[cD,dashed,line width=0.8pt,domain=0:3.5] plot (\x,{-0.0518-0.00035*\x*\x}) node[sm,right]{$F_{\mathrm{fric}}$};
  \end{scope}
\end{tikzpicture}
\end{figside}
\begin{eqside}[0.42]
\[\begin{aligned}
\xD{F_{f,y}} &= \xP{D_f}\sin\!\big(\xP{C_f}\arctan(\xP{B_f}\,\xD{\alpha_f})\big)\\
\xD{F_{r,y}} &= \xP{D_r}\sin\!\big(\xP{C_r}\arctan(\xP{B_r}\,\xD{\alpha_r})\big)\\
\xD{F_{r,x}} &= (\xP{C_{m1}}-\xP{C_{m2}}\,\xS{v_x})\,\xS{D}\\
\xD{F_{\mathrm{fric}}} &= -\xP{C_{r0}}-\xP{C_{r2}}\,\xS{v_x}^2
\end{aligned}\]
{\footnotesize $\xP{D_f},\xP{D_r}$: peak force $\approx\mu F_z$ \quad $\xS D$: duty cycle $\in[-0.1,1]$}
\deps{$F_y\leftarrow\xD{\alpha}$;\quad $F_{r,x}\leftarrow\xS{D,v_x}$;\quad $F_{\mathrm{fric}}\leftarrow\xS{v_x}$}
\restr{lateral force saturates at $D$; drive force fades with speed; front tire has no $F_x$.}
\end{eqside}
\end{panel}

% ---------------------------------------------------------------- A4
\begin{panel}{A4}{Newton--Euler in the body frame}
\begin{figside}
\begin{tikzpicture}
  \coordinate (R) at (0,0); \coordinate (G) at (1.6,0); \coordinate (F) at (3.4,0);
  \draw[black!40,line width=1.4pt] (R) -- (F);
  \node[whl] at (R) {};
  \node[whl,rotate=22] at (F) {};
  \fill (G) circle (1.6pt);
  \draw[vD] (R) -- ++(0,1.0) node[lab,above]{$F_{r,y}$};
  \draw[vD] (-1.1,0) -- (-0.33,0);
  \node[lab,text=cD,below] at (-0.75,-0.05) {$F_{r,x}$};
  \draw[vD] (F) -- ++(112:1.05) node[lab,above]{$F_{f,y}$};
  \draw[vD,dashed] (G) -- ++(-0.75,0);
  \node[lab,text=cD,above] at (1.15,0.05) {$F_{\mathrm{fric}}$};
  \draw[cS,->,line width=0.7pt] ($(G)+(-150:0.5)$) arc (-150:-30:0.5);
  \node[lab,text=cS] at ($(G)+(-90:0.72)$) {$r$};
  \draw[dim] (0,-1.15) -- (1.6,-1.15) node[sm,midway,below]{$l_r$};
  \draw[dim] (1.6,-1.15) -- (3.4,-1.15) node[sm,midway,below]{$l_f$};
\end{tikzpicture}
\end{figside}
\begin{eqside}
\[\begin{aligned}
\xP m\,\xS{\dot v_x} &= \xD{F_{r,x}}+\xD{F_{\mathrm{fric}}}-\xD{F_{f,y}}\sin\xS\delta+\xP m\,\xS{v_y}\,\xS r\\
\xP m\,\xS{\dot v_y} &= \xD{F_{r,y}}+\xD{F_{f,y}}\cos\xS\delta-\xP m\,\xS{v_x}\,\xS r\\
\xP{I_z}\,\xS{\dot r} &= \xD{F_{f,y}}\,\xP{l_f}\cos\xS\delta-\xD{F_{r,y}}\,\xP{l_r}
\end{aligned}\]
{\footnotesize $\pm m\,v\,r$ terms: the body frame itself rotates.}
\deps{$(\dot v_x,\dot v_y,\dot r)\leftarrow\xD{F}$'s, $\xS{\delta}$, $\xS{v_x,v_y,r}$}
\restr{accelerations = forces / inertia. Linearized + ZOH-discretized each QP (A9).}
\end{eqside}
\end{panel}

% ---------------------------------------------------------------- A5
\begin{panel}{A5}{Progress $s$, contouring error $e_c$, lag error $e_l$}
\begin{figside}
\begin{tikzpicture}
  \begin{scope}[shift={(0,-4.4)}]
    \draw[refl] (127:5) arc (127:58:5) node[sm,right]{$\mathbf p_{\mathrm{ref}}(\cdot)$};
    \coordinate (P) at (102:5);
    \coordinate (S) at (88:5);
    \coordinate (p) at (88:4.15);
    \coordinate (N) at ($(P)+(-78:1)$);
    \coordinate (Q) at ($(P)!(p)!(N)$);
    \draw[guide] ($(P)+(192:0.6)$) -- ($(P)+(12:2.9)$);
    \draw[vec,black!70] (P) -- ++(12:0.65) node[sm,above]{$\mathbf t$};
    \draw[black!45,dotted,line width=0.8pt] (p) -- (S);
    \draw[cD,line width=1.2pt] (P) -- (Q) node[lab,midway,left=1pt]{$e_c$};
    \draw[cD,line width=1.2pt] (Q) -- (p) node[lab,midway,below=1pt]{$e_l$};
    \caricon{(p)}{-2}{0.65}
    \node[dot,cR] at (P) {};
    \node[sm,text=cR,above left] at (P) {$s$};
    \node[odot,draw=cR] at (S) {};
    \node[sm,text=cR,above] at (S) {$s^{\ast}$};
    \node[sm,below right=4pt] at (p) {$(X,Y)$};
  \end{scope}
\end{tikzpicture}
\end{figside}
\begin{eqside}
\[\begin{aligned}
\xD{e_c} &= \phantom{-}\xR{\mathbf n(\xS s)}^{\!\top}\big(\xS{\mathbf p}-\xR{\mathbf p_{\mathrm{ref}}(\xS s)}\big)\\
\xD{e_l} &= -\xR{\mathbf t(\xS s)}^{\!\top}\big(\xS{\mathbf p}-\xR{\mathbf p_{\mathrm{ref}}(\xS s)}\big)
\end{aligned}\]
{\footnotesize $\xR{\mathbf t}=(\cos\Phi,\sin\Phi)$,\ \ $\xR{\mathbf n}=(\sin\Phi,-\cos\Phi)$ (right),\ \ $\xS{\mathbf p}=(\xS X,\xS Y)$}
\deps{$e_c,e_l\leftarrow\xS{X,Y,s}$ and the spline $\xR{\mathbf p_{\mathrm{ref}}}$}
\restr{nothing alone; $e_l\!\approx\!0 \Leftrightarrow s\!\approx\!s^{\ast}$ (projection), with no inner $\arg\min$.}
\end{eqside}
\end{panel}

% ---------------------------------------------------------------- A6
\begin{panel}{A6}{Stage cost as springs}
\begin{figside}
\begin{tikzpicture}
  \draw[refl] (-0.2,0) -- (5.6,0);
  \coordinate (P) at (3.5,0); \coordinate (Fp) at (1.9,0); \coordinate (p) at (1.9,1.15);
  \draw[cD,line width=1.1pt,decorate,decoration={zigzag,segment length=3.2pt,amplitude=2.6pt}] (Fp) -- (P);
  \draw[cD!75,line width=0.5pt,decorate,decoration={coil,aspect=0.45,segment length=6pt,amplitude=3pt}] (Fp) -- (p);
  \caricon{(p)}{0}{0.8}
  \node[dot,cR] at (P) {};
  \node[sm,text=cR,below=2pt] at (P) {$s$};
  \draw[vS,line width=1.2pt] (P) -- ++(1.15,0);
  \node[lab,text=cS,above] at ($(P)+(0.62,0.05)$) {$q_{v_s}v_s$};
  \node[lab,text=cD,below] at ($(Fp)!0.5!(P)+(0,-0.12)$) {$q_l\,e_l^2$\ \ \textcolor{cP}{\scriptsize$q_l{=}500$}};
  \node[lab,text=cD,left=5pt] at ($(Fp)!0.5!(p)+(0,0.12)$) {$q_c\,e_c^2$};
  \node[sm,text=cP,left=5pt] at ($(Fp)!0.5!(p)+(0,-0.22)$) {$q_c{=}0.1$};
\end{tikzpicture}
\end{figside}
\begin{eqside}
\[J=\sum_{k=0}^{N}\Big(\xP{q_c}\xD{e_c}^2+\xP{q_l}\xD{e_l}^2-\xP{q_{v_s}}\xS{v_s}\Big)+\text{reg.}\]
{\footnotesize reg.\ $= \xP{q_\beta}\beta^2+\xP{q_r}\xS r^2+\xP{q_\mu}(\xS\varphi-\xR\Phi)^2+\|(\xS{D,\delta,v_s})\|^2_{\xP{R}}+\|\xU{\mathbf u}\|^2_{\xP{R_\Delta}}$\\
\phantom{reg.\ }$\beta=\arctan(\xS{v_y}/\xS{v_x})$;\ \ $q_c,q_r\times10$ at $k=N$;\ \ $q_{v_s}{=}0.02$}
\deps{$e_c,e_l$ (A5), $\xS{v_s}$, $\xS{v_x,v_y,r}$, $\xU{\mathbf u}$}
\restr{soft only: $s$ pinned to the car (stiff), lateral position nearly free (weak), progress rewarded.}
\end{eqside}
\end{panel}

% ---------------------------------------------------------------- A7
\begin{panel}{A7}{Track constraint (linearized slab)}
\begin{figside}
\begin{tikzpicture}
  \begin{scope}[shift={(0,-4.3)}]
    \coordinate (P) at (95:4.5);
    \coordinate (Ao) at (95:5.25);
    \coordinate (Ai) at (95:3.75);
    \path ($(Ao)+(185:2.6)$) coordinate (o1) ($(Ao)+(5:2.6)$) coordinate (o2)
          ($(Ai)+(185:2.6)$) coordinate (i1) ($(Ai)+(5:2.6)$) coordinate (i2);
    \fill[cS!7] (o1) -- (o2) -- (i2) -- (i1) -- cycle;
    \begin{scope}
      \clip (o1) -- (o2) -- (i2) -- (i1) -- cycle;
      \fill[cL!20,even odd rule] (0,0) circle (7) (0,0) circle (5.25);
    \end{scope}
    \draw[bnd] (128:5.25) arc (128:60:5.25);
    \draw[bnd] (128:3.75) arc (128:60:3.75);
    \draw[cR,dash pattern=on 4pt off 2pt,line width=0.8pt] (128:4.5) arc (128:60:4.5);
    \draw[liml] (o1) -- (o2);
    \draw[liml] (i1) -- (i2);
    \draw[vec,black!70] (P) -- ++(95:0.5) node[sm,right]{$\bar{\mathbf n}$};
    \node[dot,cR] at (P) {};
    \node[sm,text=cR,below right] at (P) {$\bar s_k$};
    \draw[dim] ($(P)+(185:1.7)$) -- ($(Ao)+(185:1.7)$) node[sm,midway,left]{$R_{\mathrm{out}}$};
    \draw[dim] ($(P)+(185:1.7)$) -- ($(Ai)+(185:1.7)$) node[sm,midway,left]{$R_{\mathrm{in}}$};
    \node[dot,cS] at ($(P)+(5:1.3)+(95:-0.42)$) {};
    \node[sm,text=cS,below] at ($(P)+(5:1.3)+(95:-0.42)$) {$\mathbf p_k$};
    \node[sm,text=cL] at ($(o2)+(-0.15,-0.35)$) {\tiny lin.\ error};
  \end{scope}
\end{tikzpicture}
\end{figside}
\begin{eqside}
\[-\xL{R_{\mathrm{in}}}\ \le\ \xR{\bar{\mathbf n}}^{\!\top}\big(\xS{\mathbf p_k}-\xR{\mathbf p_{\mathrm{ref}}(\bar s_k)}\big)\ \le\ \xL{R_{\mathrm{out}}}\quad(+\,\xi_k\ge 0)\]
{\footnotesize $\bar s_k$: progress of the \emph{previous guess} (fixed in this QP); $\xR{\bar{\mathbf n}}$: left normal there;\\
$R_{\mathrm{in}}{=}R_{\mathrm{out}}{=}0.14$\,m constant; car treated as a point.}
\deps{$\xS{X_k,Y_k}$ only}
\restr{car point inside a straight slab; soft (slack $\xi$, penalized); can overshoot the real edge far from $\bar s_k$.}
\end{eqside}
\end{panel}

% ---------------------------------------------------------------- A8
\begin{panel}{A8}{Tire constraints}
\begin{figside}
\begin{tikzpicture}[scale=0.72]
  \fill[cS!7] (0,0) ellipse (1.83 and 1.65);
  \draw[ax] (-2.3,0) -- (2.9,0) node[sm,right]{$F_{r,x}$};
  \draw[ax] (0,-2.0) -- (0,2.3) node[sm,above]{$F_{r,y}$};
  \draw[guide] (0,0) circle (1.737);
  \draw[liml] (0,0) ellipse (1.83 and 1.65);
  \draw[vD] (0,0) -- (0.9,1.2);
  \node[dot,cD] at (0.9,1.2) {};
  \node[sm,text=cD,right=2pt] at (0.9,1.2) {$(F_{r,x},F_{r,y})$};
  \node[sm,text=cL,anchor=north west] at (1.86,-0.3) {$e_\varepsilon D_r/e_{\mathrm{long}}$};
  \node[sm,text=cL,anchor=south east] at (-0.05,1.68) {$e_\varepsilon D_r$};
  \node[sm,text=black!55,anchor=north east] at (-1.25,-1.25) {$D_r$};
\end{tikzpicture}
\end{figside}
\begin{eqside}
\[\Big(\tfrac{\xP{e_{\mathrm{long}}}\,\xD{F_{r,x}}}{\xP{F_{N,r}}}\Big)^2+\Big(\tfrac{\xD{F_{r,y}}}{\xP{F_{N,r}}}\Big)^2\ \le\ \Big(\tfrac{\xL{e_\varepsilon}\,\xP{D_r}}{\xP{F_{N,r}}}\Big)^2\]
\[|\xD{\alpha_f}|\ \le\ \xL{\alpha_{\max}}\]
{\footnotesize $e_{\mathrm{long}}{=}0.9$,\ $e_\varepsilon{=}0.95$,\ $\alpha_{\max}{=}0.6$; $F_{N,r}$ static load. Both first-order in $\xS{\mathbf x}$ around the guess.}
\deps{ellipse $\leftarrow\xS{v_x,v_y,r,D}$;\quad $|\alpha_f|\leftarrow\xS{v_x,v_y,r,\delta}$}
\restr{rear: drive + lateral force share one ellipse at 95\,\% of peak. Front: slip below $\alpha_{\max}$ (A3, red band). Soft.}
\end{eqside}
\end{panel}

% ---------------------------------------------------------------- A9
\begin{panel}{A9}{One QP per SQP step, trust region on $s$}
\begin{figside}
\begin{tikzpicture}[x=0.52cm,y=1cm]
  \fill[cL!12] (0,0.0) -- (11,2.18) -- (11,2.68) -- (0,0.5) -- cycle;
  \draw[ax] (-0.3,0) -- (12.4,0) node[sm,right]{$k$};
  \draw[ax] (0,-0.1) -- (0,2.8) node[sm,left]{$s_k$};
  \draw[black!40,line width=0.5pt] (0,0.25) -- (11,2.43);
  \foreach \k in {1,...,7,11} {
    \pgfmathsetmacro\sb{0.25+0.198*\k}
    \pgfmathsetmacro\so{\sb+0.17*sin(40*\k)}
    \node[odot,draw=black!45] at (\k,\sb) {};
    \node[dot,cS] at (\k,\so) {};
  }
  \node[dot,cS] at (0,0.25) {};
  \node[sm,left] at (0,0.25) {$\hat{\mathbf x}$};
  \node[sm] at (9,1.4) {$\cdots$};
  \foreach \k/\l in {0/0,1/1,2/2,11/N} \node[sm,below] at (\k,-0.05) {$\l$};
  \draw[decorate,decoration={brace,amplitude=4pt},black!50] (0.2,3.05) -- (11,3.05) node[sm,midway,above=4pt]{$N\,T_s=60\times20\,\text{ms}$};
  \node[sm,text=cL] at (6.3,2.55) {$|s_k-\bar s_k|\le 0.2$\,m};
  \node[sm,text=black!50] at (8.6,0.95) {$\bar s_k$ (guess)};
\end{tikzpicture}
\end{figside}
\begin{eqside}
\[\begin{aligned}
\min\ &\textstyle\sum_k \tfrac12\xS{\mathbf x_k}^{\!\top}\!Q_k\xS{\mathbf x_k}+\mathbf q_k^{\!\top}\xS{\mathbf x_k}+\tfrac12\xU{\mathbf u_k}^{\!\top}\!R\,\xU{\mathbf u_k}+\text{slack}\\
\text{s.t. }&\xS{\mathbf x_0}=\hat{\mathbf x},\quad \xS{\mathbf x_{k+1}}=A_k\xS{\mathbf x_k}+B_k\xU{\mathbf u_k}+\mathbf g_k\\
&\text{A7, A8 (soft)},\ \ \text{box bounds on }\xS{\mathbf x},\xU{\mathbf u}\\
&|\xS{s_k}-\bar s_k|\le 0.2\,\text{m}
\end{aligned}\]
{\footnotesize $(A_k,B_k,\mathbf g_k)$: exact ZOH of the Jacobian at $(\bar{\mathbf x}_k,\bar{\mathbf u}_k)$.\\
$\bar{\mathbf x}\leftarrow 0.8\,\mathbf x^{\ast}+0.2\,\bar{\mathbf x}$, \ 2 iterations, apply $\xU{\mathbf u_0}$, shift.}
\restr{each QP trusts its linearization only near the previous guess.}
\end{eqside}
\end{panel}

% =====================================================================
\clearpage
\sect{Offline \quad TUM \texttt{global\_racetrajectory\_optimization} \quad (whole lap)}

% ---------------------------------------------------------------- B0
\begin{panel}{B0}{Two pipelines}
\centering
\begin{tikzpicture}[
  bx/.style={draw=black!35,rounded corners=2pt,inner sep=3.5pt,font=\footnotesize,align=center,fill=#1!9},
  fl/.style={->,black!55,line width=0.6pt}]
  \node[font=\sffamily\small\bfseries,anchor=west] at (-1.8,0.95) {mincurv};
  \node[font=\footnotesize\itshape,text=black!50,anchor=west] at (-0.2,0.95) {path first, speed second};
  \node[bx=cR] (t1) at (0,0) {$\xR{\mathbf p_i,\ \mathbf n_i}$\\ $\xR{w_{\ell,i},\ w_{r,i}}$};
  \node[bx=cS] (q1) at (3.4,0) {QP over $\xS{\boldsymbol\alpha}$\\ $\min\sum\kappa_i^2$};
  \node[bx=cD] (r1) at (6.6,0) {$\xD{\mathbf r(s)},\ \xD{\kappa(s)}$};
  \node[bx=cS] (v1) at (10.2,0) {fwd/bwd pass\\ $\xS{v(s)},\ \xD{a_x(s)}$};
  \node[bx=cP] (g1) at (10.2,-1.3) {$\text{ggv}(v),\ a_{x,\mathrm{mach}}(v),\ c_d,\ m,\ \mu$};
  \draw[fl] (t1) -- (q1); \draw[fl] (q1) -- (r1); \draw[fl] (r1) -- (v1); \draw[fl] (g1) -- (v1);
  \draw[fl,dashed] (r1.north) to[bend right=35] node[sm,above]{IQP: re-linearize} (q1.north);
  \node[font=\sffamily\small\bfseries,anchor=west] at (-1.8,-2.35) {mintime};
  \node[font=\footnotesize\itshape,text=black!50,anchor=west] at (-0.2,-2.35) {path and speed in one problem};
  \node[bx=cR] (t2) at (0,-3.4) {$\xR{\kappa_{\mathrm{ref}}(s)}$\\ $\xR{w_\ell(s),\ w_r(s)}$};
  \node[bx=cS] (n2) at (5.0,-3.4) {NLP (collocation + IPOPT)\\ $\min\int_0^L \tfrac{dt}{ds}\,ds$\\ over $\xS{\mathbf x(s)},\ \xU{\mathbf u(s)}$};
  \node[bx=cD] (o2) at (10.2,-3.4) {$\xD{\mathbf r(s)},\ \xS{v(s)}$};
  \node[bx=cP] (p2) at (5.0,-4.95) {double track, Pacejka, load transfer, aero, $\mu$};
  \draw[fl] (t2) -- (n2); \draw[fl] (n2) -- (o2); \draw[fl] (p2) -- (n2);
\end{tikzpicture}
\restr{\texttt{mincurv}: speed never feeds back into geometry. \texttt{mintime}: geometry and speed traded off jointly.}
\end{panel}

% ---------------------------------------------------------------- B1
\begin{panel}{B1}{mincurv: the only variable is a lateral shift $\alpha_i$}
\begin{figside}
\begin{tikzpicture}
  \begin{scope}[shift={(0,-5.65)}]
    \draw[bnd] (121:7) arc (121:59:7);
    \draw[bnd] (121:5) arc (121:59:5);
    \draw[refl,dash pattern=on 4pt off 2pt] (121:6) arc (121:59:6);
    \draw[cS,line width=1pt] plot[smooth] coordinates {(121:6.62) (115:6.6) (105:6.15) (95:5.45) (85:5.45) (75:6.15) (65:6.6) (59:6.62)};
    \foreach \t/\a in {115/-0.6,105/-0.15,95/0.55,85/0.55,75/-0.15,65/-0.6} {
      \draw[->,black!40,line width=0.5pt] (\t:6) -- (\t:5.5);
      \node[dot,cR] at (\t:6) {};
      \pgfmathsetmacro\rr{6-\a}
      \node[dot,cS] at (\t:\rr) {};
    }
    \draw[cL!60,line width=3pt,opacity=0.5] (77.5:6.9) -- (77.5:5.1);
    \draw[cL,line width=0.6pt] ($(77.5:6.9)+(167.5:0.12)$) -- ($(77.5:6.9)+(-12.5:0.12)$);
    \draw[cL,line width=0.6pt] ($(77.5:5.1)+(167.5:0.12)$) -- ($(77.5:5.1)+(-12.5:0.12)$);
    \node[sm,text=cL,right=3pt] at (77.5:5.25) {$[-\bar w_\ell,\ \bar w_r]$};
    \node[sm,text=cR,above] at (95:6.02) {$\mathbf p_i$};
    \node[sm,text=cS,below] at (95:5.42) {$\mathbf r_i$};
    \node[sm,text=black!55,left] at (115:5.6) {$\mathbf n_i$};
    \node[lab,text=cS,right] at (91.5:5.75) {$\alpha_i$};
    \draw[dim] (119:6) -- (119:7) node[sm,midway,left]{$w_\ell$};
    \draw[dim] (119:6) -- (119:5) node[sm,midway,left]{$w_r$};
    \draw[->,cR,line width=0.8pt] (63:6) arc (63:58:6);
  \end{scope}
\end{tikzpicture}
\end{figside}
\begin{eqside}
\[\xD{\mathbf r_i}=\xR{\mathbf p_i}+\xS{\alpha_i}\,\xR{\mathbf n_i}\]
\[-\big(\xR{w_{\ell,i}}-\tfrac{\xP{w_{\mathrm{veh}}}}{2}\big)\ \le\ \xS{\alpha_i}\ \le\ \xR{w_{r,i}}-\tfrac{\xP{w_{\mathrm{veh}}}}{2}\]
{\footnotesize $\xR{\mathbf n_i}$ points \emph{right};\quad $\bar w=w-w_{\mathrm{veh}}/2$;\quad decision vector $\xS{\boldsymbol\alpha}\in\mathbb R^{N}$.}
\deps{$\mathbf r_i\leftarrow\xS{\alpha_i}$ and fixed track data}
\restr{each point slides only along its own normal, inside the track minus half the car width.}
\end{eqside}
\end{panel}

% ---------------------------------------------------------------- B2
\begin{panel}{B2}{mincurv: curvature, linearized}
\begin{figside}
\begin{tikzpicture}[scale=1.05]
  \draw[cR,line width=0.8pt,dash pattern=on 4pt off 2pt,domain=-2.2:2.5,samples=50,smooth] plot (\x,{0.22*\x*\x-0.75});
  \draw[cS,line width=1pt,domain=-2.2:2.5,samples=50,smooth] plot (\x,{0.22*\x*\x});
  \coordinate (Ar) at (1.2,-0.4332);
  \coordinate (A) at (1.2,0.3168);
  \draw[vec,black!55] (Ar) -- ++(0.95,0.5016) node[sm,right]{$(\xP{x'},\xP{y'})$ frozen};
  \node[dot,cR] at (Ar) {};
  \draw[vec,cS] (A) -- ++(0,0.85) node[sm,above]{$(\xS{x''},\xS{y''})$};
  \draw[guide] (A) -- ++(-0.7004,1.3265) node[sm,left]{$1/\kappa_i$};
  \draw[cD,dashed,line width=0.6pt] (-0.335,3.224) ++(-98:3.287) arc (-98:-26:3.287);
  \node[dot,cS] at (A) {};
  \node[sm,text=cS,right=2pt] at (A) {$\mathbf r_i$};
  \node[sm,text=cR] at (-1.6,-0.05) {ref.};
  \node[sm,text=cS] at (-1.75,1.25) {raceline};
\end{tikzpicture}
\end{figside}
\begin{eqside}
\[\xD{\kappa_i}=\frac{\xP{x_i'}\,\xS{y_i''}-\xP{y_i'}\,\xS{x_i''}}{\big(\xP{x_i'}^2+\xP{y_i'}^2\big)^{3/2}}
\ \Rightarrow\ \xD{\boldsymbol\kappa}\approx\xR{\boldsymbol\kappa_{\mathrm{ref}}}+\xP{E_\kappa}\,\xS{\boldsymbol\alpha}\]
\[\min_{\xS{\boldsymbol\alpha}}\ \sum_i\xD{\kappa_i}^2=\tfrac12\xS{\boldsymbol\alpha}^{\!\top}\!H\xS{\boldsymbol\alpha}+\mathbf f^{\top}\xS{\boldsymbol\alpha}
\quad\text{s.t.}\ \ |\xD{\kappa_i}|\le\xL{\kappa_{\max}}\ \ \&\ \text{B1}\]
{\footnotesize $x',y'$ taken from the reference spline (frozen); $x'',y''$ are linear in $\boldsymbol\alpha$ through the spline system. IQP: refreeze at the new line, re-solve.}
\deps{$\kappa_i\leftarrow\xS{\boldsymbol\alpha}$ (neighbours too, via splines)}
\restr{a QP: smoothest line through the corridor; no speed, mass or tire in it.}
\end{eqside}
\end{panel}

% ---------------------------------------------------------------- B3
\begin{panel}{B3}{Velocity profile: acceleration left after cornering}
\begin{figside}
\begin{tikzpicture}
  \begin{scope}
    \clip (-2.0,-1.6) rectangle (2.0,0.8);
    \fill[cS!7] (0,0) ellipse (1.6 and 1.3);
  \end{scope}
  \draw[ax] (-2.1,0) -- (2.4,0) node[sm,right]{$a_y$};
  \draw[ax] (0,-1.7) -- (0,1.85) node[sm,above]{$a_x$};
  \draw[cL,line width=0.9pt] (0,0) ellipse (1.6 and 1.3);
  \draw[cL!70,dashed,line width=0.7pt] (1.6,0) -- (0,1.3) -- (-1.6,0) -- (0,-1.3) -- cycle;
  \draw[liml] (-1.9,0.8) -- (1.9,0.8) node[sm,right]{$a_{x,\mathrm{mach}}(v)$};
  \draw[cD,line width=0.8pt] (0.75,-1.6) -- (0.75,1.6) node[sm,above]{$a_y{=}v^2\kappa$};
  \node[odot,draw=cD] at (0.75,1.148) {};
  \node[dot,cD] at (0.75,-1.148) {};
  \node[dot,cD] at (0.75,0.8) {};
  \node[sm,text=cL,anchor=north west] at (1.62,-0.18) {$\mu a_{y,\max}$};
  \node[sm,text=cL,above left] at (0,1.3) {$\mu a_{x,\max}$};
  \node[sm,text=cL] at (-1.45,1.05) {$p{=}2$};
  \node[sm,text=cL!70] at (-0.55,-0.6) {$p{=}1$};
\end{tikzpicture}
\end{figside}
\begin{eqside}
\[\xD{a_y}=\xS v^2\,\xD{\kappa}\]
\[\xD{a_x^{\mathrm{avail}}}=\xL{\mu\,a_{x,\max}(v)}\Big[1-\Big(\frac{\xD{a_y}}{\xL{\mu\,a_{y,\max}(v)}}\Big)^{\xP p}\Big]^{1/\xP p}\]
\[\begin{aligned}
\text{accel: }&\min\!\big(\xD{a_x^{\mathrm{avail}}},\,\xL{a_{x,\mathrm{mach}}(v)}\big)-\tfrac{\xP{c_d}}{\xP m}\xS v^2\\
\text{brake: }&\xD{a_x^{\mathrm{avail}}}+\tfrac{\xP{c_d}}{\xP m}\xS v^2\quad(\text{magnitude})
\end{aligned}\]
{\footnotesize ggv table: $(v,a_{x,\max},a_{y,\max})$; $p$ = \texttt{dyn\_model\_exp} $\in[1,2]$ (default 1).}
\deps{$a_x\leftarrow\xS v,\ \xD\kappa$ (from B2)}
\restr{longitudinal accel.\ allowed by grip after cornering, and by the motor.}
\end{eqside}
\end{panel}

% ---------------------------------------------------------------- B4
\begin{panel}{B4}{Velocity profile: forward--backward pass}
\begin{figside}
\begin{tikzpicture}[x=0.72cm,y=0.36cm]
  \draw[ax] (0,0) -- (10.7,0) node[sm,right]{$s$};
  \draw[ax] (0,0) -- (0,7.2) node[sm,above]{$v$};
  \draw[cL,dash pattern=on 3pt off 2pt,line width=0.6pt,domain=0:10,samples=150,smooth]
     plot (\x,{6-4*exp(-(((\x-3)/0.6)^2))-3*exp(-(((\x-7.5)/0.6)^2))});
  \draw[cS!45,dashed,line width=0.7pt,domain=3:7.5,samples=40] plot (\x,{sqrt(4+5*(\x-3))});
  \draw[cS!45,dashed,line width=0.7pt,domain=7.5:10,samples=30] plot (\x,{sqrt(9+5*(\x-7.5))});
  \draw[cL!55,dashed,line width=0.7pt,domain=0.3:3,samples=30] plot (\x,{sqrt(4+8*(3-\x))});
  \draw[cL!55,dashed,line width=0.7pt,domain=4.2:7.5,samples=30] plot (\x,{sqrt(9+8*(7.5-\x))});
  \draw[cS,line width=1.3pt,domain=0:10,samples=200]
     plot (\x,{min(6-4*exp(-(((\x-3)/0.6)^2))-3*exp(-(((\x-7.5)/0.6)^2)),
                   (\x<3)*sqrt(4+8*abs(3-\x)) + (\x>=3)*(\x<7.5)*min(sqrt(4+5*abs(\x-3)),sqrt(9+8*abs(7.5-\x))) + (\x>=7.5)*sqrt(9+5*abs(\x-7.5)))});
  \node[sm,text=cL,anchor=south] at (5.2,6.05) {$v_\kappa=\sqrt{\mu a_{y,\max}/|\kappa|}$};
  \node[sm,text=cS!70] at (5.35,3.3) {fwd};
  \node[sm,text=cL!80] at (6.45,4.6) {bwd};
  \foreach \s in {3,7.5} \draw[black!40] (\s,-0.25) -- (\s,0.25);
  \node[sm,below] at (3,-0.2) {apex};
  \node[sm,below] at (7.5,-0.2) {apex};
\end{tikzpicture}
\end{figside}
\begin{eqside}
\[\text{fwd: }\ \xS{v_{i+1}}=\sqrt{\xS{v_i}^2+2\,\xD{a_x^{+}}(\xS{v_i},\xD{\kappa_i})\,\xR{\ell_i}}\]
\[\text{bwd: }\ \xS{v_{i}}=\sqrt{\xS{v_{i+1}}^2+2\,|\xD{a_x^{-}}(\xS{v_{i+1}},\xD{\kappa_{i+1}})|\,\xR{\ell_i}}\]
\[\xS{v}=\min\big(v_\kappa,\ \text{fwd},\ \text{bwd}\big)\]
{\footnotesize $\ell_i$: segment length; $a_x^{\pm}$ from B3. No optimizer: a sweep.}
\deps{$v\leftarrow\xD\kappa(s)$, ggv, motor, drag}
\restr{speed $\le$ cornering limit, and reachable by accelerating out of / braking into each apex.}
\end{eqside}
\end{panel}

% ---------------------------------------------------------------- B5
\begin{panel}{B5}{mintime: who depends on whom}
\begin{figside}[0.56]
\begin{tikzpicture}[x=0.55cm,y=0.37cm]
  \def\ncols{11}
  \mcol{1}{v}{cS}\mcol{2}{\beta}{cS}\mcol{3}{\omega_z}{cS}\mcol{4}{n}{cS}\mcol{5}{\xi}{cS}
  \mcol{6}{\delta}{cU}\mcol{7}{F_d}{cU}\mcol{8}{F_b}{cU}\mcol{9}{\gamma_y}{cU}
  \mcol{10}{\kappa}{cR}\mcol{11}{w}{cR}
  \mrow{1}{\xD{S_f}}{1,2,4,5,10}{0}
  \mrow{2}{\xD{\alpha_{ij}}}{1,2,3,6}{0}
  \mrow{3}{\xD{F_{z,ij}}}{1,7,8,9}{0}
  \mrow{4}{\xD{F_{x,ij}}}{7,8}{0}
  \mrow{5}{\xD{F_{y,ij}}}{0}{1,2,3,6,7,8,9}
  \mrow{6.4}{\xS{v'}}{2,6}{1,3,4,5,7,8,9,10}
  \mrow{7.4}{\xS{\beta'}}{1,2,3,6}{4,5,7,8,9,10}
  \mrow{8.4}{\xS{\omega_z'}}{6}{1,2,3,4,5,7,8,9,10}
  \mrow{9.4}{\xS{n'}}{1,2,5}{4,10}
  \mrow{10.4}{\xS{\xi'}}{3,10}{1,2,4,5}
  \mrow{11.8}{T_{\mathrm{lap}}}{0}{1,2,4,5,10}
  \mrow{12.8}{\xL{\text{track}}}{4,11}{0}
  \mrow{13.8}{\xL{\text{circles}}}{0}{1,2,3,6,7,8,9}
  \mrow{14.8}{\xL{\gamma_y\ \text{eq.}}}{9}{1,2,3,6,7,8}
  \mrow{15.8}{\xL{\text{power}}}{1,7}{0}
  \mrow{16.8}{\xL{F_dF_b}}{7,8}{0}
  \mrow{17.8}{\xL{\text{rates}}}{6,7,8}{1,4,10}
  \mrow{18.8}{\xL{\text{periodic}}}{1,2,3,4,5}{0}
  \draw[black!30] (0.5,-5.7) -- (11.5,-5.7);
  \draw[black!30] (0.5,-11.1) -- (11.5,-11.1);
  \draw[black!30] (5.5,1.4) -- (5.5,-19.3);
  \draw[black!30] (9.5,1.4) -- (9.5,-19.3);
\end{tikzpicture}
\end{figside}
\begin{eqside}[0.40]
{\footnotesize
$(\cdot)'=d(\cdot)/ds$\par\medskip
\tikz\fill[black!85] (0,0) circle (2.5pt);\ appears in the formula\par
\tikz\draw[black!60,line width=0.6pt] (0,0) circle (2.2pt);\ enters through $S_f$, $\alpha$ or $F$\par\medskip
Every $d/ds$ row inherits $n,\xi,\kappa$ through $S_f=dt/ds$.}
\restr{compared with A0: no $X,Y$ states and no progress variable; position is $(s,n)$ on the fixed reference.}
\end{eqside}
\end{panel}

% ---------------------------------------------------------------- B6
\begin{panel}{B6}{mintime: curvilinear coordinates, time as a function of $s$}
\begin{figside}
\begin{tikzpicture}
  \coordinate (O) at (0,0);
  \coordinate (P) at (0:4);
  \coordinate (C) at (0:2.85);
  \draw[refl] (-32:4) arc (-32:42:4) node[sm,above]{ref.};
  \node[dot,black!55] at (O) {};
  \draw[guide] (O) -- (-22:4) node[sm,pos=0.33,below,text=cR]{$1/\kappa(s)$};
  \node[dot,cR] at (P) {};
  \node[sm,text=cR,right] at (P) {$s$};
  \draw[guide] (2.85,-1.0) -- (2.85,2.4);
  \draw[guide] (4,-0.25) -- (4,-1.15);
  \draw[guide] (C) -- ++(108:2.3);
  \draw[dim] (2.85,-1.1) -- (4,-1.1) node[sm,midway,below,text=cS]{$n$};
  \begin{scope}[shift={(C)},rotate=108]
    \draw[body] (-0.72,-0.28) rectangle (0.72,0.28);
    \fill[black!55] (0.72,0) -- (0.52,0.12) -- (0.52,-0.12) -- cycle;
  \end{scope}
  \draw[vS] (C) -- ++(122:2.15) node[lab,above]{$v$};
  \draw[cS,line width=0.6pt] ($(C)+(90:1.3)$) arc (90:108:1.3);
  \node[lab,text=cS] at ($(C)+(99:1.08)$) {$\xi$};
  \draw[cS,line width=0.6pt] ($(C)+(108:1.75)$) arc (108:122:1.75);
  \node[lab,text=cS] at ($(C)+(115:2.0)$) {$\beta$};
  \draw[cS,->,line width=0.7pt] ($(C)+(195:0.95)$) arc (195:255:0.95);
  \node[lab,text=cS] at ($(C)+(222:1.22)$) {$\omega_z$};
\end{tikzpicture}
\end{figside}
\begin{eqside}
\[\xD{S_f}\equiv\frac{dt}{ds}=\frac{1-\xS n\,\xR\kappa}{\xS v\cos(\xS\xi+\xS\beta)}\]
\[\frac{d\xS n}{ds}=\xD{S_f}\,\xS v\sin(\xS\xi+\xS\beta),\qquad \frac{d\xS\xi}{ds}=\xD{S_f}\,\xS{\omega_z}-\xR\kappa\]
\[T_{\mathrm{lap}}=\int_0^{L}\xD{S_f}\,ds\]
{\footnotesize $\xS{\mathbf x}=(v,\beta,\omega_z,n,\xi)$,\ $\xU{\mathbf u}=(\delta,F_d,F_b,\gamma_y)$;\ \ $n>0$ = \emph{left} (opposite sign to $\alpha$ in B1).}
\deps{$S_f\leftarrow\xS{v,\beta,n,\xi},\ \xR\kappa$}
\restr{time is not a variable: lap time integrates $dt/ds$ over the fixed reference length $L$.}
\end{eqside}
\end{panel}

% ---------------------------------------------------------------- B7
\begin{panel}{B7}{mintime: double-track loads}
\begin{figside}
\begin{tikzpicture}
  \draw[body] (-0.55,-0.55) rectangle (3.85,0.55);
  \draw[black!35] (0,-0.8) -- (0,0.8);
  \draw[black!35] (3.3,-0.8) -- (3.3,0.8);
  \draw[cL!80,line width=0.6pt,fill=cL!6] (3.3,-0.8) circle (0.55);
  \draw[cL!80,line width=0.6pt,fill=cL!6] (3.3,0.8) circle (0.38);
  \draw[cL!80,line width=0.6pt,fill=cL!6] (0,-0.8) circle (0.40);
  \draw[cL!80,line width=0.6pt,fill=cL!6] (0,0.8) circle (0.24);
  \node[whl] at (0,0.8) {}; \node[whl] at (0,-0.8) {};
  \node[whl,rotate=15] at (3.3,0.8) {}; \node[whl,rotate=15] at (3.3,-0.8) {};
  \fill (1.54,0) circle (1.6pt);
  \draw[vD] (1.54,0) -- ++(-0.8,0) node[sm,left]{$a_x$};
  \draw[vD] (1.54,0) -- ++(0,0.5) node[sm,above]{$a_y$};
  \node[sm,text=cL] at (4.45,-0.8) {$\propto F_{z,fr}$};
  \draw[dim] (0,-1.55) -- (1.54,-1.55) node[sm,midway,below]{$l_r$};
  \draw[dim] (1.54,-1.55) -- (3.3,-1.55) node[sm,midway,below]{$l_f$};
  \draw[dim] (-1.0,-0.8) -- (-1.0,0.8) node[sm,midway,left]{$b$};
  \node[sm,text=black!50] at (1.65,1.15) {braking + left turn};
\end{tikzpicture}
\end{figside}
\begin{eqside}[0.46]
\[\xD{F_{z,ij}}=\xP{F^{\mathrm{st}}_{z,ij}}+\tfrac12\xP{c_{l,i}}\xS v^2
\mp\tfrac{\xP h}{2\xP l}\big(\xU{F_d}{+}\xU{F_b}{-}\xP{c_d}\xS v^2{-}\xP{F_{\mathrm{roll}}}\big)
\mp\xP{k_i}\,\xU{\gamma_y}\]
{\footnotesize long.: $-$front/$+$rear;\ lat.: $-$left/$+$right;\ $k_f{=}k_{\mathrm{roll}}$, $k_r{=}1{-}k_{\mathrm{roll}}$}
\[\xU{\gamma_y}=\frac{\xP h}{\xP{\bar b}}\,\xP m\,\xD{a_y}\quad(\text{equality constraint})\]
{\footnotesize $\xD{F_{x,ij}}$: drive/brake split $\xP{k_{\mathrm{drive}}},\xP{k_{\mathrm{brake}}}$ minus rolling resistance.}
\deps{$F_z\leftarrow\xS v,\ \xU{F_d,F_b,\gamma_y}$}
\restr{each wheel's load (hence grip) follows the accelerations; $\gamma_y$ is a decision variable closed by an equality.}
\end{eqside}
\end{panel}

% ---------------------------------------------------------------- B8
\begin{panel}{B8}{mintime: tire model and friction circles}
\begin{figside}
\begin{tikzpicture}
  \begin{scope}[xscale=7,yscale=1.6]
    \draw[ax] (0,0) -- (0.38,0) node[sm,right]{$\alpha$};
    \draw[ax] (0,0) -- (0,1.3) node[sm,above]{$F_y/D$};
    \draw[guide] (0,1) -- (0.36,1);
    \draw[cD,line width=1pt,domain=0:0.35,samples=150,smooth] plot (\x,{sin(2.5*atan(10*\x-1.0*(10*\x-rad(atan(10*\x)))))});
    \node[sm,left] at (0,1) {$1$};
  \end{scope}
  \begin{scope}[shift={(3.55,0)},xscale=0.25,yscale=0.25]
    \draw[ax] (0,0) -- (10,0) node[sm,right]{$F_z$};
    \draw[ax] (0,0) -- (0,9.6) node[sm,above]{$D$};
    \draw[guide] (0,0) -- (9,9) node[sm,above]{$\mu F_z$};
    \draw[cD,line width=1pt,domain=0:9,samples=30] plot (\x,{\x*(1-0.1*\x/3)});
    \node[sm,text=cD] at (8.2,4.9) {$\varepsilon<0$};
  \end{scope}
\end{tikzpicture}
\end{figside}
\begin{eqside}
\[\xD{D_{ij}}=\xP\mu\,\xD{F_{z,ij}}\Big(1+\xP\varepsilon\,\frac{\xD{F_{z,ij}}}{\xP{F_{z0}}}\Big)\]
\[\xD{F_{y,ij}}=\xD{D_{ij}}\sin\!\Big(\xP C\arctan\big(\xP B\xD\alpha-\xP E(\xP B\xD\alpha-\arctan\xP B\xD\alpha)\big)\Big)\]
\[\Big(\frac{\xD{F_{x,ij}}}{\xP\mu\,\xD{F_{z,ij}}}\Big)^2+\Big(\frac{\xD{F_{y,ij}}}{\xP\mu\,\xD{F_{z,ij}}}\Big)^2\le\xL{1}\qquad\forall\,ij\]
\deps{$\alpha_{ij}\leftarrow\xS{v,\beta,\omega_z},\xU\delta$;\ \ $F_{z,ij}$ from B7}
\restr{each tire inside its own circle sized by its load (B7 circles); load transfer costs total grip ($\varepsilon<0$).}
\end{eqside}
\end{panel}

% ---------------------------------------------------------------- B9
\begin{panel}{B9}{mintime: the NLP}
\begin{figside}[0.40]
\begin{tikzpicture}[x=1cm,y=1cm]
  \draw[ax] (-0.2,0) -- (5.3,0) node[sm,right]{$s$};
  \foreach \t in {0.1127,0.5,0.8873} \draw[black!25] ({0.4+4.2*\t},0) -- ({0.4+4.2*\t},2.2);
  \draw[cS,line width=1pt,domain=0:1,samples=40] plot ({0.4+4.2*\x},{1+2.4*\x-1.4*\x*\x});
  \foreach \t/\m in {0.1127/0.4963,0.5/0.2381,0.8873/-0.0201} {
    \pgfmathsetmacro\xx{0.4+4.2*\t}
    \pgfmathsetmacro\yy{1+2.4*\t-1.4*\t*\t}
    \draw[cD,line width=1.1pt] ({\xx-0.6},{\yy-0.6*\m+0.12}) -- ({\xx+0.6},{\yy+0.6*\m+0.12});
    \node[dot,cS] at (\xx,\yy) {};
  }
  \node[dot,cS] at (0.4,1) {};
  \node[odot,draw=cS] at (4.6,2) {};
  \node[sm,text=cS,left] at (0.4,1) {$X_k$};
  \node[sm,text=cS,right] at (4.6,2) {$X_{k+1}$};
  \node[sm,text=cD,above] at (2.5,2.3) {slope $=h_k\mathbf f(X_{k,j},U_k,\kappa)$};
  \draw[cU,line width=1pt] (-0.1,-0.75) -- (0.4,-0.75) (0.4,-0.5) -- (4.6,-0.5) (4.6,-0.85) -- (5.1,-0.85);
  \node[sm,text=cU,above] at (2.5,-0.5) {$U_k$};
  \draw[dim] (0.4,-1.15) -- (4.6,-1.15) node[sm,midway,below]{$h_k$};
\end{tikzpicture}
\end{figside}
\begin{eqside}[0.58]
\[\min_{\xS{\mathbf x},\xU{\mathbf u}}\ \int_0^{L}\xD{S_f}\,ds+\xP{\lambda_F}\|\Delta\xU F\|^2+\xP{\lambda_\delta}\|\Delta\xU\delta\|^2\]
\footnotesize
\[\begin{aligned}
&\text{dynamics:}  && h_k\,\mathbf f(X_{k,j},\xU{U_k},\xR{\kappa})=\textstyle\sum_r C_{rj}X_{k,r},\ \ 3\text{ Legendre pts}\\
&\text{track:}     && -\xR{w_r}+\tfrac{\xP{w_{\mathrm{opt}}}}{2}\le\xS n\le\xR{w_\ell}-\tfrac{\xP{w_{\mathrm{opt}}}}{2}\\
&\text{grip, loads:}&& \text{B8 circles},\ \ \gamma_y\ \text{equality (B7)}\\
&\text{power:}     && \xS v\,\xU{F_d}\le\xL{P_{\max}},\qquad -\epsilon\le\xU{F_d}\,\xU{F_b}\le 0\\
&\text{actuators:} && \tfrac{\Delta\xU\delta}{\Delta t}\in\pm\tfrac{\xL{\delta_{\max}}}{\xP{t_\delta}},\ \
                      \tfrac{\Delta\xU{F_d}}{\Delta t}\le\tfrac{\xL{F_{d,\max}}}{\xP{t_d}},\ \
                      \tfrac{\Delta\xU{F_b}}{\Delta t}\ge-\tfrac{\xL{F_{b,\max}}}{\xP{t_b}}\\
&\text{closed lap:}&& \xS{\mathbf x}(0)=\xS{\mathbf x}(L)\\
&\text{bounds:}    && \xS v\ge 1,\ \ |\xS\beta|,|\xS{\omega_z}|,|\xS\xi|\le\tfrac{\pi}{2},\ \ \xU{F_d}\ge0\ge\xU{F_b}
\end{aligned}\]
\small
\restr{dynamics enforced only at collocation points; the lap must close; everything else is a path constraint.}
\end{eqside}
\end{panel}
}

\restoregeometry

\end{document}
